\documentclass[preprint,12pt]{elsarticle}
\usepackage{iftex}
\ifPDFTeX
  \usepackage[T1]{fontenc}
  \usepackage[utf8]{inputenc}
  \usepackage{lmodern}
\else
  \usepackage[no-math]{fontspec}
\fi
\usepackage{amsmath}
\usepackage{amssymb}
\usepackage{amsfonts}
\usepackage{mathrsfs}
\usepackage{amsthm}
\usepackage{mathtools}
\usepackage{geometry}
\usepackage{enumitem}
\usepackage{longtable}
\usepackage{booktabs}
\usepackage{array}
\usepackage{microtype}
\usepackage{graphicx}
\graphicspath{{../figures/}{figures/}{./}}
\usepackage[dvipsnames]{xcolor}
\usepackage{float}
\usepackage[pagebackref=true]{hyperref}
\usepackage{bookmark}
\usepackage{pdfcomment}
\usepackage{listings}
\lstset{language=Python,basicstyle=\ttfamily\scriptsize,breaklines=true,breakatwhitespace=false,columns=fullflexible,keepspaces=true,showstringspaces=false,frame=single,numbers=left,numberstyle=\tiny,numbersep=4pt,xleftmargin=1.6em,postbreak=\mbox{$\hookrightarrow$\space},literate={á}{{\'a}}1 {é}{{\'e}}1 {í}{{\'\i}}1 {ó}{{\'o}}1 {ú}{{\'u}}1 {ñ}{{\~n}}1 {ó}{{\'o}}1 {ç}{{\c c}}1 {–}{{--}}1 {—}{{---}}1 {≤}{{$\leq$}}1 {≥}{{$\geq$}}1 {→}{{$\rightarrow$}}1 {Δ}{{$\Delta$}}1 {δ}{{$\delta$}}1 {τ}{{$\tau$}}1 {γ}{{$\gamma$}}1 {π}{{$\pi$}}1 {θ}{{$\theta$}}1 {Φ}{{$\Phi$}}1 {ζ}{{$\zeta$}}1 {σ}{{$\sigma$}}1 {λ}{{$\lambda$}}1 {ρ}{{$\rho$}}1 {ε}{{$\varepsilon$}}1 {²}{{$^2$}}1 {³}{{$^3$}}1 {·}{{$\cdot$}}1 {∫}{{$\int$}}1 {∞}{{$\infty$}}1 {≈}{{$\approx$}}1 {≍}{{$\asymp$}}1 {≪}{{$\ll$}}1 {ϑ}{{$\vartheta$}}1 {ϖ}{{$\varpi$}}1 {Λ}{{$\Lambda$}}1 {√}{{$\surd$}}1 {μ}{{$\mu$}}1 {ω}{{$\omega$}}1 {α}{{$\alpha$}}1 {β}{{$\beta$}}1 {η}{{$\eta$}}1 {χ}{{$\chi$}}1 {ξ}{{$\xi$}}1 {Σ}{{$\Sigma$}}1 {−}{{-}}1 {×}{{$\times$}}1 {±}{{$\pm$}}1 {°}{{$^\circ$}}1 {…}{{...}}1}

% Hover previews for internal references (no layout change).
\makeatletter
\let\CG@origref\ref
\let\CG@origeqref\eqref
\newcommand{\CG@gettip}[1]{\@ifundefined{CGTip@#1}{Reference: #1}{\csname CGTip@#1\endcsname}}
\DeclareRobustCommand{\CGref}[1]{\pdftooltip{\CG@origref{#1}}{\CG@gettip{#1}}}
\DeclareRobustCommand{\CGeqref}[1]{\pdftooltip{\CG@origeqref{#1}}{\CG@gettip{#1}}}
\makeatother
% Tooltip texts shown when hovering over internal references.
\makeatletter
\expandafter\gdef\csname CGTip@sec:introduction\endcsname{Section 1: Introduction (p. 3)}
\expandafter\gdef\csname CGTip@subsec:clock_intro\endcsname{Section 1.1: The Riemann--von Mangoldt clock and the ordinates (p. 3)}
\expandafter\gdef\csname CGTip@eq:intro_rvm\endcsname{Eq. (1.1), p. 3: N(T)=N\_0(T)+(1)/(pi)arg zeta ((1)/(2)+iT)+O((1)/(T)), N\_0(T)=(T)/(2pi)log(T)/(2pi e)+(7)/(8),}
\expandafter\gdef\csname CGTip@eq:intro_clock\endcsname{Eq. (1.2), p. 3: N\_0(tau\_k)=k-(1)/(2), tau\_k=(2pi (k-(11)/(8)))/(W\_0((k-(11)/(8))/e))(k>=1),}
\expandafter\gdef\csname CGTip@eq:intro_defect\endcsname{Eq. (1.3), p. 3: s\_k=k-(1)/(2)-N\_0(gamma\_k), u\_k(X)=s\_k-S\_X(gamma\_k), S\_X(t)=-(1)/(pi)sum\_p<= Xarg (1-p\textasciicircum{}-1/2-it).}
\expandafter\gdef\csname CGTip@subsec:program\endcsname{Section 1.2: The approximation of the zeros by the clock and the primes (p. 3)}
\expandafter\gdef\csname CGTip@subsec:known_intro\endcsname{Section 1.3: What is known and only used (p. 3)}
\expandafter\gdef\csname CGTip@subsec:comparison\endcsname{Section 1.4: Comparison with the cited antecedents (p. 4)}
\expandafter\gdef\csname CGTip@tab:antecedents\endcsname{Table 1 (p. 4): Cited antecedents, conclusions proved here, and the scope of the difference.}
\expandafter\gdef\csname CGTip@subsec:main_results\endcsname{Section 1.5: What this paper adds (p. 6)}
\expandafter\gdef\csname CGTip@eq:intro_main\endcsname{Eq. (1.4), p. 6: (1)/(no.I\_T)sum\_kin I\_Tu\_k(X)\textasciicircum{}2=(1)/(2pi\textasciicircum{}2)log(1)/(c)+O(1+sqrt(log(1/c))),}
\expandafter\gdef\csname CGTip@eq:intro_D\endcsname{Eq. (1.5), p. 6: (1)/(no.I\_T)sum\_kin I\_T(tau\_k\textasciicircum{}X-gamma\_k)\textasciicircum{}2=(log log log T+O(sqrt(log log log T)))/(log\textasciicircum{}2T).}
\expandafter\gdef\csname CGTip@eq:intro_E\endcsname{Eq. (1.6), p. 7: limsup\_k->inf|tau\_k-gamma\_k|<=(1)/(log 2)int\_0\textasciicircum{}1(e(sigma) dsigma)/(sqrt(sigma(2-sigma))),}
\expandafter\gdef\csname CGTip@eq:intro_ratio\endcsname{Eq. (1.7), p. 7: (<(tau\_k\textasciicircum{}X-gamma\_k)\textasciicircum{}2>\_I\_T)/(<(tau\_k-gamma\_k)\textasciicircum{}2>\_I\_T)\textasciitilde{}(log log log T)/(2log log T)->0,}
\expandafter\gdef\csname CGTip@subsec:status\endcsname{Section 1.6: Degree of proof (p. 8)}
\expandafter\gdef\csname CGTip@tab:status\endcsname{Table 2 (p. 8): Main results: degree of proof and role.}
\expandafter\gdef\csname CGTip@sec:notation\endcsname{Section 2: Notation (p. 8)}
\expandafter\gdef\csname CGTip@subsec:symbols\endcsname{Section 2.1: List of symbols (p. 8)}
\expandafter\gdef\csname CGTip@subsec:counting\endcsname{Section 2.2: Counting functions (p. 11)}
\expandafter\gdef\csname CGTip@eq:RvM\endcsname{Eq. (2.1), p. 11: N(t)=(1)/(pi)vartheta(t)+1+S(t), S(t)=(1)/(pi)arg zeta ((1)/(2)+it)=O(log t),}
\expandafter\gdef\csname CGTip@eq:N0_def\endcsname{Eq. (2.2), p. 11: Theta\_0(t)=(t)/(2)log(t)/(2pi e)-(pi)/(8), N\_0(t)=(1)/(pi)Theta\_0(t)+1, varpi(t)=vartheta(t)-Theta\_0(t).}
\expandafter\gdef\csname CGTip@eq:N_decomposition\endcsname{Eq. (2.3), p. 11: N(t)=N\_0(t)+S(t)+(varpi(t))/(pi), L(t)=N\_0'(t)=(1)/(2pi)log(t)/(2pi).}
\expandafter\gdef\csname CGTip@eq:IT_size\endcsname{Eq. (2.4), p. 11: no.I\_T=TL(T)+O(T).}
\expandafter\gdef\csname CGTip@subsec:clock\endcsname{Section 2.3: The RvM clock and the displacements (p. 11)}
\expandafter\gdef\csname CGTip@eq:clock\endcsname{Eq. (2.5), p. 11: tau\_x=(2pi (x-(11)/(8)))/(W\_0((x-(11)/(8))/e)), (dtau\_x)/(dx)=(1)/(L(tau\_x))=(2pi)/(log(tau\_x/2pi)),}
\expandafter\gdef\csname CGTip@def:displacements\endcsname{Definition 2.1 (p. 11): For k>=1, s\_k=k-(1)/(2)-N\_0(gamma\_k) and delta\_k=tau\_k-gamma\_k.}
\expandafter\gdef\csname CGTip@eq:s_delta\endcsname{Eq. (2.6), p. 11: s\_k=L(xi\_k) delta\_kfor some xi\_k between gamma\_k and tau\_k .}
\expandafter\gdef\csname CGTip@lem:s_bounds\endcsname{Lemma 2.2 (p. 11): For every k>=1, N(gamma\_k\textasciicircum{}-)-N\_0(gamma\_k)+(1)/(2)<= s\_k<= N(gamma\_k\textasciicircum{}+)-N\_0(gamma\_k)-(1)/(2), and |s\_k-(varpi(gamma\_k))/(pi)|<=max (|S(gamma\_k\textasciicircum{}-)|,|S(gamma\_k\textasciicircum{}+)|)-(1)/(2) .}
\expandafter\gdef\csname CGTip@eq:s_bounds\endcsname{Eq. (2.7), p. 11: N(gamma\_k\textasciicircum{}-)-N\_0(gamma\_k)+(1)/(2)<= s\_k<= N(gamma\_k\textasciicircum{}+)-N\_0(gamma\_k)-(1)/(2),}
\expandafter\gdef\csname CGTip@eq:s_S_bound\endcsname{Eq. (2.8), p. 11: |s\_k-(varpi(gamma\_k))/(pi)|<=max (|S(gamma\_k\textasciicircum{}-)|,|S(gamma\_k\textasciicircum{}+)|)-(1)/(2) .}
\expandafter\gdef\csname CGTip@subsec:euler_phase\endcsname{Section 2.4: Euler phases and the corrected clock (p. 12)}
\expandafter\gdef\csname CGTip@eq:SA_def\endcsname{Eq. (2.9), p. 12: S\_A(t)=sum\_nin Aa\_nsin(tlog n), v\_A=(1)/(2)sum\_nin Aa\_n\textasciicircum{}2 ;}
\expandafter\gdef\csname CGTip@eq:SX_def\endcsname{Eq. (2.10), p. 12: S\_X(t)=-(1)/(pi)sum\_p<= Xarg (1-p\textasciicircum{}-1/2-it)=-(1)/(pi)sum\_p<= Xsum\_m>=1(sin(mtlog p))/(m p\textasciicircum{}m/2)=S\_A\_X(t).}
\expandafter\gdef\csname CGTip@eq:sigma_M_def\endcsname{Eq. (2.11), p. 12: |S\_X(t)|<=sigma\_X:=(1)/(pi)sum\_p<= Xsum\_m>=1(1)/(m p\textasciicircum{}m/2), |S\_X'(t)|<= M\_X:=(1)/(pi)sum\_p<= X(log p)/(sqrt(p)-1),}
\expandafter\gdef\csname CGTip@def:defect\endcsname{Definition 2.3 (p. 12): For X>=2 let F\_X=N\_0+S\_X and u\_k(X)=k-(1)/(2)-F\_X(gamma\_k)=s\_k-S\_X(gamma\_k). The one-step clock and its residual are tau\_k\textasciicircum{}X=tau\_k-(S\_X(tau\_k))/(L(tau\_k)), delta\_k\textasciicircum{}X=tau\_k\textasciicircum{}X-gamma\_k, r\_k(X)=L(tau\_k) delta\_k\textasciicircum{}X .}
\expandafter\gdef\csname CGTip@eq:one_step_clock\endcsname{Eq. (2.12), p. 12: tau\_k\textasciicircum{}X=tau\_k-(S\_X(tau\_k))/(L(tau\_k)), delta\_k\textasciicircum{}X=tau\_k\textasciicircum{}X-gamma\_k, r\_k(X)=L(tau\_k) delta\_k\textasciicircum{}X .}
\expandafter\gdef\csname CGTip@sec:known_results\endcsname{Section 3: External results (p. 12)}
\expandafter\gdef\csname CGTip@thm:known_basic\endcsname{Theorem 3.1 ([ref, Thms. 9.4, 9.9 and S 9.12], [ref]) (p. 12): S(t)=O(log t); gamma\_k+1-gamma\_k is bounded; and S\_1(t)=int\_0\textasciicircum{}tS(y) dy=O(log t).}
\expandafter\gdef\csname CGTip@thm:known_BW\endcsname{Theorem 3.2 (Bellotti--Wong [ref, Thm. 1.1, Cor. 1.5 and (1.2)]) (p. 13): For T>= e, |N(T)-(T)/(2pi)log(T)/(2pi e)|<=BW(T) and |S(T)|<=BWS(T), where BW(t) =min 0.10076log t+0.24460log log t+8.08344,; 0.11200log t+0.12567log log t+3.77417,; BWS(t) =0.10076log t+min 0.24460log log t+7.20844, 1.68845log log t+1.50956. Moreover |S(T)|<=2.5167 for 0<= T<=30 610 046 000.}
\expandafter\gdef\csname CGTip@eq:BWS_def\endcsname{Eq. (3.1), p. 13: BW(t) =min 0.10076log t+0.24460log log t+8.08344,; 0.11200log t+0.12567log log t+3.77417,; BWS(t) =0.10076log t+min 0.24460log log t+7.20844, 1.68845log log t+1.50956.}
\expandafter\gdef\csname CGTip@thm:known_selberg_moments\endcsname{Theorem 3.3 (Selberg [ref]; see [ref, p. 225]) (p. 13): For every integer j>=1, int\_T\textasciicircum{}2TS(t)\textasciicircum{}2j dt=kappa\_j T(log log T)\textasciicircum{}j+O\_j(T(log log T)\textasciicircum{}j-1/2).}
\expandafter\gdef\csname CGTip@thm:known_korolev\endcsname{Theorem 3.4 (Selberg; Korolev [ref], [ref, Thm. 4.2, p. 73]) (p. 13): Let Delta\_k be the integer such that g\_k+Delta\_k-2<gamma\_k<= g\_k+Delta\_k-1 (in [ref] the Gram points are t\_m=g\_m-1 and t\_m-1<gamma\_k<= t\_m with m=k+Delta\_k), and let 0<varepsilon\_0<10\textasciicircum{}-3. For every fixed j>=1, as n\_0->inf, uniformly for n\_0\textasciicircum{}27/82+0.9varepsilon\_0<= h\_0<= n\_0\textasciicircum{}27/82+varepsilon\_0, sum\_n\_0<k<= n\_0+h\_0Del...}
\expandafter\gdef\csname CGTip@eq:korolev_moments\endcsname{Eq. (3.2), p. 13: sum\_n\_0<k<= n\_0+h\_0Delta\_k\textasciicircum{}2j=kappa\_jh\_0(log log n\_0)\textasciicircum{}j(1+o(1)), |sum\_n\_0<k<= n\_0+h\_0Delta\_k\textasciicircum{}2j-1|<<\_j h\_0(log log n\_0)\textasciicircum{}j-1.}
\expandafter\gdef\csname CGTip@thm:known_fujii_selberg\endcsname{Theorem 3.5 (Selberg's formula with Fujii's bounds [ref, S 3, pp. 232--233]) (p. 13): Let 0<a\_2<(1)/(4) be fixed. There are T\_0 and C, depending only on a\_2, such that for T>= T\_0, T\textasciicircum{}1/2+a\_2<H<T and every integer k>=1 with x\_rm F=T\textasciicircum{}a\_2/(60k)>=2, int\_T\textasciicircum{}T+H(S(t)-S\textasciicircum{}rm pr\_x\_rm F\textasciicircum{}3(t))\textasciicircum{}2kdt<=(Ck\textasciicircum{}2)\textasciicircum{}2kH .}
\expandafter\gdef\csname CGTip@eq:fujii_selberg\endcsname{Eq. (3.3), p. 13: int\_T\textasciicircum{}T+H(S(t)-S\textasciicircum{}rm pr\_x\_rm F\textasciicircum{}3(t))\textasciicircum{}2kdt<=(Ck\textasciicircum{}2)\textasciicircum{}2kH .}
\expandafter\gdef\csname CGTip@lem:selberg_uniform\endcsname{Lemma 3.6 (Selberg's approximation uniformly in the order) (p. 14): Let kappa\_rm F=1/(300log 2) and Y\_k=T\textasciicircum{}1/(80k). There are absolute constants C and T\_0 such that for T>= T\_0 and every integer k with 1<= k<=kappa\_rm Flog T, int\_T\textasciicircum{}2T(S(t)-S\textasciicircum{}rm pr\_Y\_k(t))\textasciicircum{}2kdt<=(Ck\textasciicircum{}2)\textasciicircum{}2k T . In particular int\_T\textasciicircum{}2T(S-S\textasciicircum{}rm pr\_Y\_0)\textasciicircum{}4<< T with Y\_0=Y\_2=T\textasciicircum{}1/160.}
\expandafter\gdef\csname CGTip@eq:selberg_uniform\endcsname{Eq. (3.4), p. 14: int\_T\textasciicircum{}2T(S(t)-S\textasciicircum{}rm pr\_Y\_k(t))\textasciicircum{}2kdt<=(Ck\textasciicircum{}2)\textasciicircum{}2k T .}
\expandafter\gdef\csname CGTip@rem:korolev_uniform\endcsname{Remark 3.7 (An independent route to the uniformity in the order) (p. 14): Korolev [ref, Lemma 3.13], as quoted in [ref, Lemma 12, p. 16] (we have not seen the original), proves: for fixed 0<varepsilon\_rm K<10\textasciicircum{}-3, T\_1>= T\_0(varepsilon\_rm K), x\_rm K=T\_1\textasciicircum{}0.1varepsilon\_rm K, T\_1\textasciicircum{}27/82+0.9varepsilon\_rm K<= H<= T\_1\textasciicircum{}27/82+varepsilon\_rm K, integers 1<= m<= c\_rm Klog x\_rm K with an absolute consta...}
\expandafter\gdef\csname CGTip@eq:korolev_lemma\endcsname{Eq. (3.5), p. 14: int\_T\_1\textasciicircum{}T\_1+H|pi S(t)+sum\_p<= Y\_rm K(sin(tlog p))/(sqrt(p))|\textasciicircum{}2mdt<=(c\_0H)/(56 varepsilon\_rm K)((pi c\_0)/(varepsilon\_rm K))\textasciicircum{}2m, c\_0=2880 .}
\expandafter\gdef\csname CGTip@thm:known_fujii_difference\endcsname{Theorem 3.8 (Fujii [ref, Lemma, p. 140], in the form of [ref, (9.25.3), p. 246]) (p. 15): There are absolute constants A and T\_0 such that for every integer k>=1, T>= T\_0 and 0<= h<=(1)/(2)T, [ int\_0\textasciicircum{}T|S(t+h)-S(t)|\textasciicircum{}2kdt<= TAk\textasciicircum{}4log(3+hlog T)\textasciicircum{}k . ]}
\expandafter\gdef\csname CGTip@eq:fujii_difference\endcsname{Eq. (3.6), p. 15: int\_T-1\textasciicircum{}2T(S(t+h)-S(t))\textasciicircum{}2kdt<= C T(log(2+hlog T))\textasciicircum{}k(k=1,2, 0<h<=1, T>= T\_0),}
\expandafter\gdef\csname CGTip@thm:known_MV\endcsname{Theorem 3.9 (Montgomery--Vaughan [ref, Cor. 3]) (p. 15): For complex z\_n, real t\_0 and H>0, int\_t\_0\textasciicircum{}t\_0+H|sum\_n<= Yz\_nn\textasciicircum{}-it|\textasciicircum{}2dt=sum\_n<= Y|z\_n|\textasciicircum{}2(H+O(n)) with an absolute implied constant.}
\expandafter\gdef\csname CGTip@thm:known_mertens\endcsname{Theorem 3.10 (Mertens [ref]; see [ref, Ch. I.1, Thms. 1.8 and 1.10]) (p. 15): For y>=2, sum\_p<= y(log p)/(p)=log y+O(1) and sum\_p<= y(1)/(p)=log log y+O(1).}
\expandafter\gdef\csname CGTip@thm:known_selberg_clt\endcsname{Theorem 3.11 (Selberg [ref]; Tsang [ref]; see [ref, (1.1)]) (p. 15): For fixed real y\_1<y\_2, (1)/(T)meastin[T,2T]:y\_1<S(t)/ s\_T<= y\_2->(1)/(sqrt(2pi))int\_y\_1\textasciicircum{}y\_2e\textasciicircum{}-y\textasciicircum{}2/2dy, where s\_T=((1)/(2pi\textasciicircum{}2)log log T)\textasciicircum{}1/2.}
\expandafter\gdef\csname CGTip@thm:known_maximal\endcsname{Theorem 3.12 (Hardy--Littlewood [ref, Thm. I.1], as quoted in [ref, (1.2)]; Kinnunen [ref, Thm. 1.4]) (p. 15): Let Mf(x)=sup\_y>0(1)/(2y)int\_x-y\textasciicircum{}x+y|f| be the centred maximal function of fin L\textasciicircum{}1\_rm loc(R). For 1<p<inf and fin L\textasciicircum{}p(R), || Mf||\_p<= A\_p||f||\_p, where A\_p is bounded for p>=2. If f and f' belong to L\textasciicircum{}p(R), 1<p<inf, then Mf has a locally absolutely continuous representative, (Mf)'in L\textasciicircum{}p(R), and |(Mf)'|<= M(f') almos...}
\expandafter\gdef\csname CGTip@thm:known_carleson\endcsname{Theorem 3.13 (Carleson [ref], Hunt [ref]; on R see [ref, (1.1) and p. 1]) (p. 15): For fin L\textasciicircum{}2(R)n L\textasciicircum{}4(R) with Fourier transform Ff(omega)=int\_R f(t)e\textasciicircum{}-iomega tdt, let Cf(x)=sup\_omega\_1in R|(1)/(2pi)int\_-inf\textasciicircum{}omega\_1 Ff(omega)e\textasciicircum{}iomega xdomega |. Then || Cf||\_4<<||f||\_4.}
\expandafter\gdef\csname CGTip@thm:known_exponents\endcsname{Theorem 3.14 ([ref, Lemma 2.8 and (7.57)]; Bourgain [ref, S 5, Thm. 6]) (p. 16): (a) If (k, l) is an exponent pair, so is A(k, l)=((k)/(2 k+2),(1)/(2)+(l)/(2 k+2)). (b) If (k, l) is an exponent pair with l- k>=(1)/(2), then zeta(l- k+it)<<|t|\textasciicircum{} klog|t| for |t|>=2. (c) For every varepsilon>0, ((13)/(84)+varepsilon,(55)/(84)+varepsilon) is an exponent pair.}
\expandafter\gdef\csname CGTip@thm:known_PL\endcsname{Theorem 3.15 (Phragmen--Lindelof principle in a strip [ref, Lemma 2.5], [ref, Lemma 3]) (p. 16): Let sigma\_1<sigma\_2 and e\_1>= e\_2>=0. If |zeta(sigma\_i+it)|<= C\_i|t|\textasciicircum{}e\_ilog|t| for |t|>=2 and i=1,2, then |zeta(sigma+it)|<<|t|\textasciicircum{}e(sigma)log|t| uniformly for sigma\_1<=sigma <=sigma\_2 and |t|>=2, where e is the linear function with e(sigma\_i)=e\_i.}
\expandafter\gdef\csname CGTip@thm:known_backlund\endcsname{Theorem 3.16 (Backlund's method [ref, Lemmas 2.1, 2.3 and 2.4], after [ref]) (p. 16): Let T be large, varkappa in N, f\_varkappa(s)=(1)/(2)[((s+iT-1)zeta(s+iT))\textasciicircum{}varkappa+((s-iT-1)zeta(s-iT))\textasciicircum{}varkappa ], and for a disc D(sigma\_o,varrho) let Z\_varkappa(sigma\_o,varrho)=(1)/(varkappa)sum\_zlog(varrho)/(|z-sigma\_o|) over the zeros z of f\_varkappa in D(sigma\_o,varrho). Let 1<sigma\_o<varrho, sigma\textasciicircum{}sharp=sigma...}
\expandafter\gdef\csname CGTip@rem:korolev\endcsname{Remark 3.17 (Known consequences for the clock) (p. 16): (a) Moments and law of s\_k. Since vartheta(gamma\_k)/pi in(k+Delta\_k-2,k+Delta\_k-1] and vartheta/pi=N\_0-1+varpi/pi, one has s\_k+Delta\_k=k+Delta\_k-(1)/(2)-N\_0(gamma\_k)in[-(1)/(2)+(varpi(gamma\_k))/(pi),(1)/(2)+(varpi(gamma\_k))/(pi)), so s\_k=-Delta\_k+ d\_k with | d\_k|<1. The index set I\_T is the integer interval (N(T),N(...}
\expandafter\gdef\csname CGTip@eq:s_moments\endcsname{Eq. (3.7), p. 16: < s\_k\textasciicircum{}2j>\_I\_T=kappa\_j(log log T)\textasciicircum{}j(1+o(1)), < s\_k\textasciicircum{}2j-1>\_I\_T=O\_j((log log T)\textasciicircum{}j-1).}
\expandafter\gdef\csname CGTip@eq:BW_pointwise\endcsname{Eq. (3.8), p. 17: -(BW(gamma\_k)+3/8)/(L(tau\_k))<=delta\_k<=(BW(gamma\_k)-11/8)/(L(gamma\_k))(k>=1).}
\expandafter\gdef\csname CGTip@sec:identities\endcsname{Section 4: Theorem A: an exact defect identity at the zeros (p. 17)}
\expandafter\gdef\csname CGTip@thm:identity_general\endcsname{Theorem 4.1 (p. 17): Let 1<= k\_1<k\_2 be indices, F a real C\textasciicircum{}1 function on an open interval containing [gamma\_k\_1,gamma\_k\_2], and u\_k=k-(1)/(2)-F(gamma\_k) for k\_1<= k<= k\_2. For every integer j>=1, int\_gamma\_k\_1\textasciicircum{}gamma\_k\_2(N(t)-F(t))\textasciicircum{}2jF'(t) dt=sum\_k=k\_1+1\textasciicircum{}k\_2-1Psi\_j(u\_k)+(((1)/(2)+u\_k\_1)\textasciicircum{}2j+1+((1)/(2)-u\_k\_2)\textasciicircum{}2j+1)/(2j+1), where Psi\_j(u)=...}
\expandafter\gdef\csname CGTip@eq:identity_general\endcsname{Eq. (4.1), p. 17: int\_gamma\_k\_1\textasciicircum{}gamma\_k\_2(N(t)-F(t))\textasciicircum{}2jF'(t) dt=sum\_k=k\_1+1\textasciicircum{}k\_2-1Psi\_j(u\_k)+(((1)/(2)+u\_k\_1)\textasciicircum{}2j+1+((1)/(2)-u\_k\_2)\textasciicircum{}2j+1)/(2j+1),}
\expandafter\gdef\csname CGTip@eq:Psi_def\endcsname{Eq. (4.2), p. 17: Psi\_j(u)=(((1)/(2)+u)\textasciicircum{}2j+1+((1)/(2)-u)\textasciicircum{}2j+1)/(2j+1)=sum\_l=0\textasciicircum{}jbinom(2j+1,2l)(4\textasciicircum{}l-j)/(2j+1) u\textasciicircum{}2l=u\textasciicircum{}2j+sum\_l<jbinom(2j+1,2l)(4\textasciicircum{}l-j)/(2j+1)u\textasciicircum{}2l.}
\expandafter\gdef\csname CGTip@eq:identity_quadratic\endcsname{Eq. (4.3), p. 17: sum\_k=k\_1+1\textasciicircum{}k\_2-1u\_k\textasciicircum{}2=int\_gamma\_k\_1\textasciicircum{}gamma\_k\_2(N(t)-F(t))\textasciicircum{}2F'(t) dt-(k\_2-k\_1-1)/(12)-(((1)/(2)+u\_k\_1)\textasciicircum{}3+((1)/(2)-u\_k\_2)\textasciicircum{}3)/(3) .}
\expandafter\gdef\csname CGTip@eq:identity_one_gap\endcsname{Eq. (4.4), p. 17: int\_gamma\_k\textasciicircum{}gamma\_k+1(k-F(t))\textasciicircum{}2jF'(t) dt =((k-F(gamma\_k))\textasciicircum{}2j+1-(k-F(gamma\_k+1))\textasciicircum{}2j+1)/(2j+1); =(((1)/(2)+u\_k)\textasciicircum{}2j+1+((1)/(2)-u\_k+1)\textasciicircum{}2j+1)/(2j+1),}
\expandafter\gdef\csname CGTip@thm:identity_g\endcsname{Theorem 4.2 (p. 18): Let k\_1<k\_2, F and u\_k be as in Theorem 4.1, and let g:R -> C be continuous with a primitive G. Then int\_gamma\_k\_1\textasciicircum{}gamma\_k\_2 g(N(t)-F(t))F'(t) dt=sum\_k=k\_1+1\textasciicircum{}k\_2-1[ G(u\_k+(1)/(2))- G(u\_k-(1)/(2))]+ G(u\_k\_1+(1)/(2))- G(u\_k\_2-(1)/(2)).}
\expandafter\gdef\csname CGTip@eq:identity_g\endcsname{Eq. (4.5), p. 18: int\_gamma\_k\_1\textasciicircum{}gamma\_k\_2 g(N(t)-F(t))F'(t) dt=sum\_k=k\_1+1\textasciicircum{}k\_2-1[ G(u\_k+(1)/(2))- G(u\_k-(1)/(2))]+ G(u\_k\_1+(1)/(2))- G(u\_k\_2-(1)/(2)).}
\expandafter\gdef\csname CGTip@cor:identity_exp\endcsname{Corollary 4.3 (p. 18): For real omega !=0, (2sin(omega/2))/(omega)sum\_k=k\_1+1\textasciicircum{}k\_2-1e\textasciicircum{}iomega u\_k=int\_gamma\_k\_1\textasciicircum{}gamma\_k\_2e\textasciicircum{}iomega(N(t)-F(t))F'(t) dt-(e\textasciicircum{}iomega(u\_k\_1+1/2)-e\textasciicircum{}iomega(u\_k\_2-1/2))/(iomega).}
\expandafter\gdef\csname CGTip@eq:identity_exp\endcsname{Eq. (4.6), p. 18: (2sin(omega/2))/(omega)sum\_k=k\_1+1\textasciicircum{}k\_2-1e\textasciicircum{}iomega u\_k=int\_gamma\_k\_1\textasciicircum{}gamma\_k\_2e\textasciicircum{}iomega(N(t)-F(t))F'(t) dt-(e\textasciicircum{}iomega(u\_k\_1+1/2)-e\textasciicircum{}iomega(u\_k\_2-1/2))/(iomega).}
\expandafter\gdef\csname CGTip@cor:signed_law\endcsname{Corollary 4.4 (Signed law of s\_k from the t-law of S) (p. 18): For fixed real xi, < e\textasciicircum{}ixi s\_k/ s\_T>\_I\_T-> e\textasciicircum{}-xi\textasciicircum{}2/2. Hence the empirical law of s\_k/ s\_T on I\_T tends to the standard normal law.}
\expandafter\gdef\csname CGTip@rem:twelfth\endcsname{Remark 4.5 (p. 19): Psi\_j(u)=int\_u-1/2\textasciicircum{}u+1/2x\textasciicircum{}2jdx is the mean of x\textasciicircum{}2j over the unit interval centred at u; in particular Psi\_1(u)=u\textasciicircum{}2+(1)/(12) and (1)/(12)=int\_-1/2\textasciicircum{}1/2x\textasciicircum{}2dx. The constant comes from the unit jump of N-F at each zero, through the primitive in Theorem 4.2; when F is not monotone, N-F need not run through an interval of...}
\expandafter\gdef\csname CGTip@cor:identity_clock\endcsname{Corollary 4.6 (p. 19): For X>=2 and 1<= k\_1<k\_2, sum\_k=k\_1+1\textasciicircum{}k\_2-1u\_k(X)\textasciicircum{}2= int\_gamma\_k\_1\textasciicircum{}gamma\_k\_2(S-S\_X+(varpi)/(pi))\textasciicircum{}2(L+S\_X')dt; -(k\_2-k\_1-1)/(12)-(((1)/(2)+u\_k\_1(X))\textasciicircum{}3+((1)/(2)-u\_k\_2(X))\textasciicircum{}3)/(3) .}
\expandafter\gdef\csname CGTip@eq:identity_corrected\endcsname{Eq. (4.7), p. 19: sum\_k=k\_1+1\textasciicircum{}k\_2-1u\_k(X)\textasciicircum{}2= int\_gamma\_k\_1\textasciicircum{}gamma\_k\_2(S-S\_X+(varpi)/(pi))\textasciicircum{}2(L+S\_X')dt; -(k\_2-k\_1-1)/(12)-(((1)/(2)+u\_k\_1(X))\textasciicircum{}3+((1)/(2)-u\_k\_2(X))\textasciicircum{}3)/(3) .}
\expandafter\gdef\csname CGTip@lem:first_moment_identity\endcsname{Lemma 4.7 (First-moment identity) (p. 19): Let 2pi<t\_1<t\_2 be points which are not ordinates and psi in C\textasciicircum{}1[t\_1,t\_2], real or complex. Then sum\_t\_1<gamma\_k<= t\_2s\_kpsi(gamma\_k)=int\_t\_1\textasciicircum{}t\_2G(t)psi(t)L(t) dt+(1)/(2)[psi G\textasciicircum{}2]\_t\_1\textasciicircum{}t\_2-(1)/(2)int\_t\_1\textasciicircum{}t\_2G(t)\textasciicircum{}2psi'(t) dt .}
\expandafter\gdef\csname CGTip@eq:first_moment_identity\endcsname{Eq. (4.8), p. 19: sum\_t\_1<gamma\_k<= t\_2s\_kpsi(gamma\_k)=int\_t\_1\textasciicircum{}t\_2G(t)psi(t)L(t) dt+(1)/(2)[psi G\textasciicircum{}2]\_t\_1\textasciicircum{}t\_2-(1)/(2)int\_t\_1\textasciicircum{}t\_2G(t)\textasciicircum{}2psi'(t) dt .}
\expandafter\gdef\csname CGTip@eq:first_moment_measures\endcsname{Eq. (4.9), p. 19: d(G\textasciicircum{}2)=-2GL dt+sum\_t\_1<gamma<t\_22G\textasciicircum{}*(gamma)nu(gamma) delta\_gamma ,}
\expandafter\gdef\csname CGTip@sec:main_theorem\endcsname{Section 5: Theorem B: how much prime information the clock needs (p. 19)}
\expandafter\gdef\csname CGTip@thm:main\endcsname{Theorem 5.1 (p. 19): There are absolute constants T\_0 and C such that for T>= T\_0 and every X with 2<= X<= T\textasciicircum{}1/2, writing c=log X/log T, |(1)/(no.I\_T)sum\_kin I\_Tu\_k(X)\textasciicircum{}2-(1)/(2pi\textasciicircum{}2)log(1)/(c)|<= C(1+sqrt(log(1/c))).}
\expandafter\gdef\csname CGTip@eq:main_theorem\endcsname{Eq. (5.1), p. 19: |(1)/(no.I\_T)sum\_kin I\_Tu\_k(X)\textasciicircum{}2-(1)/(2pi\textasciicircum{}2)log(1)/(c)|<= C(1+sqrt(log(1/c))).}
\expandafter\gdef\csname CGTip@cor:fixed_cutoff\endcsname{Corollary 5.2 (fixed cutoff) (p. 20): For fixed X>=2, (1)/(no.I\_T)sum\_kin I\_Tu\_k(X)\textasciicircum{}2=(1)/(2pi\textasciicircum{}2)log log T+O\_X(sqrt(log log T)), and no.kin I\_T:|u\_k(X)|<= y\_0=o(no.I\_T) for every fixed y\_0>0.}
\expandafter\gdef\csname CGTip@cor:power_cutoff\endcsname{Corollary 5.3 (power cutoff) (p. 20): For fixed 0<c<=(1)/(2) and X=T\textasciicircum{}c, limsup\_T->inf(1)/(no.I\_T)sum\_kin I\_Tu\_k(T\textasciicircum{}c)\textasciicircum{}2<=(1)/(2pi\textasciicircum{}2)log(1)/(c)+C(1+sqrt(log(1/c))), and no.kin I\_T:|u\_k(T\textasciicircum{}c)|>w\_T=o(no.I\_T) whenever w\_T->inf. The same holds for X=N(T)\textasciicircum{}c with fixed 0<c<(1)/(2), since then X=T\textasciicircum{}c' with c'=c+O(log log T/log T).}
\expandafter\gdef\csname CGTip@eq:power_cutoff\endcsname{Eq. (5.2), p. 20: limsup\_T->inf(1)/(no.I\_T)sum\_kin I\_Tu\_k(T\textasciicircum{}c)\textasciicircum{}2<=(1)/(2pi\textasciicircum{}2)log(1)/(c)+C(1+sqrt(log(1/c))),}
\expandafter\gdef\csname CGTip@cor:dichotomy\endcsname{Corollary 5.4 (dichotomy) (p. 20): Let X=X(T) with 2<= X<= T\textasciicircum{}1/2 for T>= T\_0. Then [ sup\_T>= T\_0(1)/(no.I\_T)sum\_kin I\_Tu\_k(X(T))\textasciicircum{}2<inf if and only ifliminf\_T->inf(log X)/(log T)>0 . ] For fixed X\_0>=2 and cin(0,(1)/(2)], (1)/(no.I\_T)sum\_kin I\_T(u\_k(X\_0)\textasciicircum{}2-u\_k(T\textasciicircum{}c)\textasciicircum{}2)=(1)/(2pi\textasciicircum{}2)log log T+O(sqrt(log log T)) .}
\expandafter\gdef\csname CGTip@eq:dichotomy\endcsname{Eq. (5.3), p. 20: (1)/(no.I\_T)sum\_kin I\_T(u\_k(X\_0)\textasciicircum{}2-u\_k(T\textasciicircum{}c)\textasciicircum{}2)=(1)/(2pi\textasciicircum{}2)log log T+O(sqrt(log log T)) .}
\expandafter\gdef\csname CGTip@prop:rounding\endcsname{Proposition 5.5 (Rounding the corrected count misassigns a positive proportion) (p. 20): There are absolute constants c\_*>0, kappa\_0>0 and T\_0 such that for T>= T\_0 and 2<= X<= T\textasciicircum{}c\_*, [ (1)/(T)meastin[T,2T]:floor(F\_X(t)+(1)/(2))!= N(t)>=kappa\_0 . ]}
\expandafter\gdef\csname CGTip@rem:comparison_cutoff\endcsname{Remark 5.6 (p. 21): Since logno.p<= X=log X+O(log log X), the condition of Corollary 5.4 says that the logarithm of the number of primes must grow proportionally to log T. A fixed number of primes leaves almost all zeros violating the corrected equation by more than any fixed amount; growth of the number of primes alone is necessary but...}
\expandafter\gdef\csname CGTip@rem:localization\endcsname{Remark 5.7 (Localization of the ordinates) (p. 21): If M\_X<=(1)/(2)L(T), then F\_X'>=L-M\_X>=(1)/(2)L(T) on [T,2T], so F\_X is strictly increasing there. For kin I\_T with F\_X(T)<= k-(1)/(2)<= F\_X(2T) the equation F\_X(t)=k-(1)/(2) has exactly one root t\_k\textasciicircum{}Xin[T,2T] (intermediate value theorem), and u\_k(X)=F\_X(t\_k\textasciicircum{}X)-F\_X(gamma\_k) gives |gamma\_k-t\_k\textasciicircum{}X|<=2|u\_k(X)|/L(T) (mea...}
\expandafter\gdef\csname CGTip@subsec:main_lemmas\endcsname{Section 5.1: Lemmas on prime polynomials (p. 21)}
\expandafter\gdef\csname CGTip@eq:D_def\endcsname{Eq. (5.4), p. 21: W\_1(t)=-(1)/(pi)sum\_pin J(sin(tlog p))/(sqrt(p)), V=sum\_pin J(1)/(p),; W=S\textasciicircum{}rm pr\_Y\_0-S\_X=+- W\_1-S\_X\textasciicircum{}(2)-S\_X\textasciicircum{}(>=3),}
\expandafter\gdef\csname CGTip@lem:V_size\endcsname{Lemma 5.8 (p. 21): V=log(1/c)+O(1) uniformly for log 2/log T<= c<=(1)/(2).}
\expandafter\gdef\csname CGTip@lem:trig_MV\endcsname{Lemma 5.9 (p. 21): For real z\_n, 2<= n<= Y, int\_T\textasciicircum{}2T(sum\_n<= Yz\_nsin(tlog n))\textasciicircum{}2dt=(T)/(2)sum\_n<= Yz\_n\textasciicircum{}2+O(sum\_n<= Ynz\_n\textasciicircum{}2+(sum\_n<= Y|z\_n|)\textasciicircum{}2), and the same with cos in place of sin.}
\expandafter\gdef\csname CGTip@eq:trig_MV\endcsname{Eq. (5.5), p. 21: int\_T\textasciicircum{}2T(sum\_n<= Yz\_nsin(tlog n))\textasciicircum{}2dt=(T)/(2)sum\_n<= Yz\_n\textasciicircum{}2+O(sum\_n<= Ynz\_n\textasciicircum{}2+(sum\_n<= Y|z\_n|)\textasciicircum{}2),}
\expandafter\gdef\csname CGTip@lem:D_moments\endcsname{Lemma 5.10 (p. 21): With absolute implied constants, int\_T\textasciicircum{}2T W\textasciicircum{}2 dt =(T)/(2pi\textasciicircum{}2)V+O(T(1+sqrt(V))), ; int\_T\textasciicircum{}2T W\textasciicircum{}4 dt << T(1+V)\textasciicircum{}2, ; int\_T\textasciicircum{}2TS\_X'(t)\textasciicircum{}2 dt << T(1+log X)\textasciicircum{}2 .}
\expandafter\gdef\csname CGTip@eq:D_second\endcsname{Eq. (5.6), p. 21: int\_T\textasciicircum{}2T W\textasciicircum{}2 dt =(T)/(2pi\textasciicircum{}2)V+O(T(1+sqrt(V))),}
\expandafter\gdef\csname CGTip@eq:D_fourth\endcsname{Eq. (5.7), p. 21: int\_T\textasciicircum{}2T W\textasciicircum{}4 dt << T(1+V)\textasciicircum{}2,}
\expandafter\gdef\csname CGTip@eq:SX_derivative\endcsname{Eq. (5.8), p. 21: int\_T\textasciicircum{}2TS\_X'(t)\textasciicircum{}2 dt << T(1+log X)\textasciicircum{}2 .}
\expandafter\gdef\csname CGTip@subsec:main_proof\endcsname{Section 5.2: Proof of Theorem 5.1 (p. 22)}
\expandafter\gdef\csname CGTip@eq:E0_bounds\endcsname{Eq. (5.9), p. 22: int\_T\textasciicircum{}2TE\_0\textasciicircum{}4 dt<< T,int\_T\textasciicircum{}2TE\_0\textasciicircum{}2 dt<= T\textasciicircum{}1/2(int\_T\textasciicircum{}2TE\_0\textasciicircum{}4 dt)\textasciicircum{}1/2<< T .}
\expandafter\gdef\csname CGTip@eq:J_asymptotic\endcsname{Eq. (5.10), p. 22: E\_2=(T)/(2pi\textasciicircum{}2)V+O(T(1+sqrt(V))),}
\expandafter\gdef\csname CGTip@eq:R_fourth\endcsname{Eq. (5.11), p. 22: int\_T\textasciicircum{}2TE\_X\textasciicircum{}4 dt<< T(1+V)\textasciicircum{}2 .}
\expandafter\gdef\csname CGTip@eq:main_step1\endcsname{Eq. (5.12), p. 22: sum\_kin I\_Tu\_k(X)\textasciicircum{}2=int\_T\textasciicircum{}2TE\_X\textasciicircum{}2L dt+int\_T\textasciicircum{}2TE\_X\textasciicircum{}2S\_X' dt-(no.I\_T)/(12)+O(T\textasciicircum{}3/4).}
\expandafter\gdef\csname CGTip@eq:main_step2\endcsname{Eq. (5.13), p. 23: |int\_T\textasciicircum{}2TE\_X\textasciicircum{}2S\_X' dt|<=(int\_T\textasciicircum{}2TE\_X\textasciicircum{}4)\textasciicircum{}1/2(int\_T\textasciicircum{}2TS\_X'\textasciicircum{}2)\textasciicircum{}1/2<< T(1+V)(1+log X).}
\expandafter\gdef\csname CGTip@eq:main_step3\endcsname{Eq. (5.14), p. 23: (1)/(no.I\_T)sum\_kin I\_Tu\_k(X)\textasciicircum{}2=(V)/(2pi\textasciicircum{}2)+O(1+sqrt(V)+(1+V)/(log T)+(1+V) c),}
\expandafter\gdef\csname CGTip@subsec:variable_cutoff\endcsname{Section 5.3: A cutoff varying with the index (p. 23)}
\expandafter\gdef\csname CGTip@lem:maximal_primes\endcsname{Lemma 5.11 (p. 23): Let T>= T\_0 and 2<= X\_-<X\_+<=(T/3)\textasciicircum{}1/2 with X\_+<=4X\_-; let V\_+-=sum\_X\_-<p<= X\_+p\textasciicircum{}-1 and [ Q(t)=max\_X\_-<= y<= X\_+|sum\_X\_-<p<= yp\textasciicircum{}-1/2-it| . ] Then Q is Lipschitz, | Q'(t)|<=max\_X\_-<= y<= X\_+|sum\_X\_-<p<= ylog p p\textasciicircum{}-1/2-it| for almost every t, and int\_T-3\textasciicircum{}2T+3 Q\textasciicircum{}4 dt<< TV\_+-\textasciicircum{}2,int\_T-3\textasciicircum{}2T+3| Q'|\textasciicircum{}4 dt<< T(log X\_+)\textasciicircum{}4V\_+-\textasciicircum{}2...}
\expandafter\gdef\csname CGTip@eq:maximal_L4\endcsname{Eq. (5.15), p. 23: int\_T-3\textasciicircum{}2T+3 Q\textasciicircum{}4 dt<< TV\_+-\textasciicircum{}2,int\_T-3\textasciicircum{}2T+3| Q'|\textasciicircum{}4 dt<< T(log X\_+)\textasciicircum{}4V\_+-\textasciicircum{}2,}
\expandafter\gdef\csname CGTip@eq:maximal_points\endcsname{Eq. (5.16), p. 23: < Q(gamma\_k)\textasciicircum{}2>\_I\_T<< V\_+-(1+(log X\_+)/(log T)sqrt(log log T))+(X\_+)/(T),< Q(tau\_k)\textasciicircum{}2>\_ I\textasciicircum{}tau\_T<< V\_+- .}
\expandafter\gdef\csname CGTip@eq:maximal_stieltjes\endcsname{Eq. (5.17), p. 24: sum\_kin I\_T Q(gamma\_k)\textasciicircum{}2=int\_T\textasciicircum{}2T Q\textasciicircum{}2L dt+[ Q\textasciicircum{}2G]\_T\textasciicircum{}2T-2int\_T\textasciicircum{}2TG Q Q' dt,}
\expandafter\gdef\csname CGTip@thm:variable_cutoff\endcsname{Theorem 5.12 (p. 24): For (log 2)/(log T)<= c<=(1)/(2) let X\_k=max 2,(tau\_k/2pi)\textasciicircum{}c and tilde u\_k=s\_k-S\_X\_k(gamma\_k). Uniformly in c, |<tilde u\_k\textasciicircum{}2>\_I\_T-(1)/(2pi\textasciicircum{}2)log(1)/(c)|<<1+sqrt(log(1/c)) . For fixed c, <(tilde u\_k-u\_k(X\_-))\textasciicircum{}2>\_I\_T->0, with X\_- as in the proof.}
\expandafter\gdef\csname CGTip@eq:variable_cutoff\endcsname{Eq. (5.18), p. 24: |<tilde u\_k\textasciicircum{}2>\_I\_T-(1)/(2pi\textasciicircum{}2)log(1)/(c)|<<1+sqrt(log(1/c)) .}
\expandafter\gdef\csname CGTip@sec:accounting\endcsname{Section 6: Theorem C: exact accounting for the one-step clock (p. 25)}
\expandafter\gdef\csname CGTip@eq:two_responses\endcsname{Eq. (6.1), p. 25: R\textasciicircum{}gamma\_n(I)=(2)/(a\_n)< s\_ksin(gamma\_klog n)>\_I, R\textasciicircum{}tau\_n(I)=(2)/(a\_n)< s\_ksin(tau\_klog n)>\_I,}
\expandafter\gdef\csname CGTip@eq:Cr_def\endcsname{Eq. (6.2), p. 25: Cr\_A(I)=< S\_A(tau\_k)\textasciicircum{}2>\_I-v\_A .}
\expandafter\gdef\csname CGTip@thm:accounting\endcsname{Theorem 6.1 (Finite-window accounting) (p. 25): Let I be a finite nonempty set of indices and A a set of prime powers with sum\_nin A|a\_n|<inf. item For every k and X>=2, r\_k(X)=u\_k(X)+(S\_X(gamma\_k)-S\_X(tau\_k))+chi\_k,; chi\_k=(gamma\_k)/(2pi)((tau\_k)/(gamma\_k)-1-log(tau\_k)/(gamma\_k)),(delta\_k\textasciicircum{}2)/(4pi max(gamma\_k,tau\_k))<=chi\_k<=(delta\_k\textasciicircum{}2)/(4pi min(gamma\_k,tau\_k)),;...}
\expandafter\gdef\csname CGTip@eq:one_step_decomposition\endcsname{Eq. (6.3), p. 25: r\_k(X)=u\_k(X)+(S\_X(gamma\_k)-S\_X(tau\_k))+chi\_k,; chi\_k=(gamma\_k)/(2pi)((tau\_k)/(gamma\_k)-1-log(tau\_k)/(gamma\_k)),(delta\_k\textasciicircum{}2)/(4pi max(gamma\_k,tau\_k))<=chi\_k<=(delta\_k\textasciicircum{}2)/(4pi min(gamma\_k,tau\_k)),; |S\_X(gamma\_k)-S\_X(tau\_k)|<= M\_X|delta\_k|.}
\expandafter\gdef\csname CGTip@eq:exact_gain\endcsname{Eq. (6.4), p. 25: < s\_k\textasciicircum{}2-(s\_k-lambda S\_A(tau\_k))\textasciicircum{}2>\_I =sum\_nin Aa\_n\textasciicircum{}2(lambda R\textasciicircum{}tau\_n(I)-(lambda\textasciicircum{}2)/(2))-lambda\textasciicircum{}2Cr\_A(I).}
\expandafter\gdef\csname CGTip@eq:full_factor_gain\endcsname{Eq. (6.5), p. 25: < s\_k\textasciicircum{}2-(s\_k-S\_X(tau\_k))\textasciicircum{}2>\_I=D\_X(I)-Cr\_X(I),; D\_X(I)=sum\_nin A\_Xa\_n\textasciicircum{}2(R\textasciicircum{}tau\_n(I)-(1)/(2))=2< s\_kS\_X(tau\_k)>\_I-v\_X,}
\expandafter\gdef\csname CGTip@eq:physical_gain\endcsname{Eq. (6.6), p. 26: <(L(tau\_k)delta\_k)\textasciicircum{}2-r\_k(X)\textasciicircum{}2>\_I=D\_X(I)-Cr\_X(I)+2<chi\_kS\_X(tau\_k)>\_I,; |<chi\_kS\_X(tau\_k)>\_I|<=sigma\_Xmax\_kin Ichi\_k .}
\expandafter\gdef\csname CGTip@rem:accounting_meaning\endcsname{Remark 6.2 (p. 26): (6.5) is an exact sign criterion for the gain of the full Euler phase at the clock points: the gain is positive exactly when the diagonal responses outweigh the cross term. No independence of the phases tau\_klog p is assumed. A threshold R\textasciicircum{}tau\_p=(1)/(2) describes a single diagonal prime contribution and does not det...}
\expandafter\gdef\csname CGTip@rem:matrix_projection\endcsname{Remark 6.3 (Optimal trigonometric correction) (p. 26): For a finite set A of frequencies let varsigma\_n(k)=sin(tau\_klog n), Q=(<varsigma\_n\_1varsigma\_n\_2>\_I)\_n\_1,n\_2in A, b=(< s\_kvarsigma\_n(k)>\_I)\_nin A and Upsilon\_lambda=sum\_nlambda\_nvarsigma\_n. Ordinary least squares gives <(s\_k-Upsilon\_lambda(k))\textasciicircum{}2>\_I=< s\_k\textasciicircum{}2>\_I-2lambda\textasciicircum{} T b+lambda\textasciicircum{} T Qlambda, minimized by lambda= Q\textasciicircum{}d...}
\expandafter\gdef\csname CGTip@rem:transport\endcsname{Remark 6.4 (Clock point against zero) (p. 26): For a prime power n, sin(gamma\_kl\_n)-sin(tau\_kl\_n)=int\_0\textasciicircum{}1(d)/(dy)sin((tau\_k-ydelta\_k)l\_n) dy, whence, exchanging the finite mean and the integral, R\textasciicircum{}gamma\_n(I)-R\textasciicircum{}tau\_n(I)=-(2l\_n)/(a\_n)int\_0\textasciicircum{}1< s\_kdelta\_kcos ((tau\_k-ydelta\_k)l\_n)>\_I dy . For fixed n and I=I\_T the right side is O(<|s\_kdelta\_k|>\_I\_T)=O(log log T/log T...}
\expandafter\gdef\csname CGTip@eq:response_transport\endcsname{Eq. (6.7), p. 26: R\textasciicircum{}gamma\_n(I)-R\textasciicircum{}tau\_n(I)=-(2l\_n)/(a\_n)int\_0\textasciicircum{}1< s\_kdelta\_kcos ((tau\_k-ydelta\_k)l\_n)>\_I dy .}
\expandafter\gdef\csname CGTip@rem:fixed_sup\endcsname{Remark 6.5 (Fixed prime sets) (p. 26): For fixed X, |delta\_k\textasciicircum{}X-delta\_k|<=sigma\_X/L(tau\_k)->0 and |u\_k(X)-s\_k|<=sigma\_X. A fixed set of primes changes neither limsup\_k|delta\_k| nor the unboundedness of the defect.}
\expandafter\gdef\csname CGTip@sec:computable\endcsname{Section 7: Theorem D: a computable clock that approaches the zeros (p. 27)}
\expandafter\gdef\csname CGTip@eq:r_decomposition\endcsname{Eq. (7.1), p. 27: r\_k(X)=u\_k(X)+ t\_k(X)+chi\_k, t\_k(X)=S\_X(gamma\_k)-S\_X(tau\_k),}
\expandafter\gdef\csname CGTip@lem:SX_fourth\endcsname{Lemma 7.1 (p. 27): For T>= T\_0 and 2<= X<= T\textasciicircum{}1/2, [ int\_T-3\textasciicircum{}2T+3S\_X'(t)\textasciicircum{}4 dt<< T(1+log X)\textasciicircum{}4,int\_T-3\textasciicircum{}2T+3S\_X''(t)\textasciicircum{}4 dt<< T(1+log X)\textasciicircum{}8 . ]}
\expandafter\gdef\csname CGTip@lem:transport\endcsname{Lemma 7.2 (Transport) (p. 27): For T>= T\_0 and 2<= X<= T\textasciicircum{}1/2, < t\_k(X)\textasciicircum{}2>\_I\_T<<(log log T (1+log X)\textasciicircum{}2)/(log\textasciicircum{}2T).}
\expandafter\gdef\csname CGTip@eq:transport\endcsname{Eq. (7.2), p. 27: < t\_k(X)\textasciicircum{}2>\_I\_T<<(log log T (1+log X)\textasciicircum{}2)/(log\textasciicircum{}2T).}
\expandafter\gdef\csname CGTip@eq:delta_fourth\endcsname{Eq. (7.3), p. 27: <delta\_k\textasciicircum{}4>\_I\_T=< s\_k\textasciicircum{}4L(xi\_k)\textasciicircum{}-4>\_I\_T<<((log log T)\textasciicircum{}2)/(log\textasciicircum{}4T).}
\expandafter\gdef\csname CGTip@eq:transport_maximal\endcsname{Eq. (7.4), p. 27: | t\_k(X)|=|int\_tau\_k\textasciicircum{}gamma\_k h(y) dy|<=int\_tau\_k-|delta\_k|\textasciicircum{}tau\_k+|delta\_k|| h|<=2|delta\_k| M h(tau\_k)}
\expandafter\gdef\csname CGTip@thm:computable_clock\endcsname{Theorem 7.3 (Theorem D) (p. 28): There are absolute constants T\_0 and C such that for T>= T\_0 and 2<= X<= T\textasciicircum{}1/2, with c=log X/log T, |< r\_k(X)\textasciicircum{}2>\_I\_T-(1)/(2pi\textasciicircum{}2)log(1)/(c)|<= C(1+sqrt(log(1/c)))(1+c\textasciicircum{}2log log T).}
\expandafter\gdef\csname CGTip@eq:computable_clock\endcsname{Eq. (7.5), p. 28: |< r\_k(X)\textasciicircum{}2>\_I\_T-(1)/(2pi\textasciicircum{}2)log(1)/(c)|<= C(1+sqrt(log(1/c)))(1+c\textasciicircum{}2log log T).}
\expandafter\gdef\csname CGTip@cor:computable_clock\endcsname{Corollary 7.4 (p. 28): Let X=T\textasciicircum{}1/sqrt(log log T). Then < r\_k(X)\textasciicircum{}2>\_I\_T =(1)/(4pi\textasciicircum{}2)log log log T+O(sqrt(log log log T)), ; (1)/(no.I\_T)sum\_kin I\_T(tau\_k\textasciicircum{}X-gamma\_k)\textasciicircum{}2 =(log log log T+O(sqrt(log log log T)))/(log\textasciicircum{}2T), while (1)/(no.I\_T)sum\_kin I\_T(tau\_k-gamma\_k)\textasciicircum{}2=(2+o(1))log log T/log\textasciicircum{}2T; the ratio is asymptotic to log log log T/(2log log T).}
\expandafter\gdef\csname CGTip@eq:computable_units\endcsname{Eq. (7.6), p. 28: < r\_k(X)\textasciicircum{}2>\_I\_T =(1)/(4pi\textasciicircum{}2)log log log T+O(sqrt(log log log T)),}
\expandafter\gdef\csname CGTip@eq:computable_height\endcsname{Eq. (7.7), p. 28: (1)/(no.I\_T)sum\_kin I\_T(tau\_k\textasciicircum{}X-gamma\_k)\textasciicircum{}2 =(log log log T+O(sqrt(log log log T)))/(log\textasciicircum{}2T),}
\expandafter\gdef\csname CGTip@cor:local_clock\endcsname{Corollary 7.5 (A cutoff read off the clock) (p. 28): For tau\_k>e\textasciicircum{}e let X\_k\textasciicircum{}rm loc=tau\_k\textasciicircum{}1/sqrt(log log tau\_k) and tau\_k\textasciicircum{}rm loc=tau\_k-S\_X\_k\textasciicircum{}rm loc(tau\_k)/L(tau\_k). Then [ (1)/(no. I\_T\textasciicircum{}tau)sum\_kin I\_T\textasciicircum{}tau(tau\_k\textasciicircum{}rm loc-gamma\_k)\textasciicircum{}2=(log log log T+O(sqrt(log log log T)))/(log\textasciicircum{}2T). ]}
\expandafter\gdef\csname CGTip@prop:monotone\endcsname{Proposition 7.6 (The root of the corrected equation in the monotone regime) (p. 29): Let X=X(T)>=2 with M\_X<=(1)/(2)L(T), and let I'\_T be the set of kin I\_T with F\_X(T)<= k-(1)/(2)<= F\_X(2T), so that the root t\_k\textasciicircum{}Xin[T,2T] of Remark 5.7 exists. Then X<<log\textasciicircum{}2T, no.(I\_T I'\_T)<<log T, and [ (1)/(no.I\_T)sum\_kin I'\_T(t\_k\textasciicircum{}X-gamma\_k)\textasciicircum{}2>=(4)/(9L(2T)\textasciicircum{}2)((1)/(2pi\textasciicircum{}2)log log T-O(sqrt(log log T)))>>(log log T)/(log\textasciicircum{}...}
\expandafter\gdef\csname CGTip@rem:newton\endcsname{Remark 7.7 (The Newton term and the restriction on c) (p. 29): Let xi\textasciicircum{}X\_k be a point between gamma\_k and tau\_k with t\_k(X)=-S\_X'(xi\_k\textasciicircum{}X)delta\_k (mean value theorem). Since delta\_k=(r\_k(X)+S\_X(tau\_k))/L(tau\_k) by (2.12), (7.1) becomes the exact relation r\_k(X)(1+(S\_X'(xi\_k\textasciicircum{}X))/(L(tau\_k)))=u\_k(X)+chi\_k-y\_k S\_X'(xi\_k\textasciicircum{}X). The term y\_kS\_X'(xi\_k\textasciicircum{}X) is the second-order term of Newton'...}
\expandafter\gdef\csname CGTip@eq:newton_exact\endcsname{Eq. (7.8), p. 29: r\_k(X)(1+(S\_X'(xi\_k\textasciicircum{}X))/(L(tau\_k)))=u\_k(X)+chi\_k-y\_k S\_X'(xi\_k\textasciicircum{}X).}
\expandafter\gdef\csname CGTip@sec:iteration\endcsname{Section 8: Iterating the Newton-type step (p. 30)}
\expandafter\gdef\csname CGTip@subsec:recursion\endcsname{Section 8.1: The exact recursion (p. 30)}
\expandafter\gdef\csname CGTip@eq:newton_iteration\endcsname{Eq. (8.1), p. 30: x\textasciicircum{}(j)\_k=x\textasciicircum{}(j-1)\_k-(N\_0(x\textasciicircum{}(j-1)\_k)+S\_X\_j(x\textasciicircum{}(j-1)\_k)-(k-(1)/(2)))/(L(x\textasciicircum{}(j-1)\_k)), e\_j,k=x\textasciicircum{}(j)\_k-gamma\_k,}
\expandafter\gdef\csname CGTip@prop:recursion\endcsname{Proposition 8.1 (Exact recursion) (p. 30): Let j>=1 and suppose x\textasciicircum{}(j-1)\_k>2pi. Then L(x\textasciicircum{}(j-1)\_k) e\_j,k=u\_k(X\_j)+(S\_X\_j(gamma\_k)-S\_X\_j(x\textasciicircum{}(j-1)\_k))+chi\_j,k, 0<=chi\_j,k<=(e\_j-1,k\textasciicircum{}2)/(4pi min(gamma\_k,x\textasciicircum{}(j-1)\_k)) . Let varphi\_jin C\textasciicircum{}inf(R), 0<=varphi\_j<=1, be 1 on an interval containing all segments between gamma\_k and x\textasciicircum{}(j-1)\_k, kin I, and g\_j= M(varphi\_jS'\_X\_j)....}
\expandafter\gdef\csname CGTip@eq:recursion_identity\endcsname{Eq. (8.2), p. 30: L(x\textasciicircum{}(j-1)\_k) e\_j,k=u\_k(X\_j)+(S\_X\_j(gamma\_k)-S\_X\_j(x\textasciicircum{}(j-1)\_k))+chi\_j,k, 0<=chi\_j,k<=(e\_j-1,k\textasciicircum{}2)/(4pi min(gamma\_k,x\textasciicircum{}(j-1)\_k)) .}
\expandafter\gdef\csname CGTip@eq:recursion_norms\endcsname{Eq. (8.3), p. 30: ||L(x\textasciicircum{}(j-1))e\_j||\_2r,I<=||u(X\_j)||\_2r,I+2 ||e\_j-1||\_4r,I || g\_j(x\textasciicircum{}(j-1))||\_4r,I+||chi\_j||\_2r,I,; ||chi\_j||\_2r,I<=(||e\_j-1||\_4r,I\textasciicircum{}2)/(4pi min\_kin Imin(gamma\_k,x\textasciicircum{}(j-1)\_k)) .}
\expandafter\gdef\csname CGTip@lem:higher_moments\endcsname{Lemma 8.2 (p. 31): There is an absolute constant C such that for every integer r>=1, every T>= T\_0 with 2r<=kappa\_rm Flog T and every 2<= X<= T\textasciicircum{}1/(160r), with V as in Lemma 5.8 computed with Y\_2r=T\textasciicircum{}1/(160r) in place of Y\_0, [ int\_T\textasciicircum{}2TE\_X\textasciicircum{}4r dt<=(C(r\textasciicircum{}2+sqrt(r(1+V))))\textasciicircum{}4rT, < u\_k(X)\textasciicircum{}2r>\_I\_T<=(C(r\textasciicircum{}2+sqrt(r(1+V))))\textasciicircum{}2r, ] and for every r>=1...}
\expandafter\gdef\csname CGTip@thm:tails\endcsname{Theorem 8.3 (Exponential tails of the ordinal defect) (p. 32): There are absolute constants C\_8>=1, lambda\_0 and T\_0 such that, for T>= T\_0, every lambda >=lambda\_0 and every X=T\textasciicircum{}c>=2 with exp (-(lambda/C\_8)\textasciicircum{}3/2)<= c<=(1)/(C\_8sqrt(lambda)), no.kin I\_T:|u\_k(X)|>lambda<=2\textasciicircum{} 2-sqrt(lambda/C\_8) no.I\_T .}
\expandafter\gdef\csname CGTip@eq:tails_range\endcsname{Eq. (8.4), p. 32: exp (-(lambda/C\_8)\textasciicircum{}3/2)<= c<=(1)/(C\_8sqrt(lambda)),}
\expandafter\gdef\csname CGTip@eq:tails\endcsname{Eq. (8.5), p. 32: no.kin I\_T:|u\_k(X)|>lambda<=2\textasciicircum{} 2-sqrt(lambda/C\_8) no.I\_T .}
\expandafter\gdef\csname CGTip@lem:multiplicity\endcsname{Lemma 8.4 (Multiplicity) (p. 32): int\_R tilde n(u)\textasciicircum{}2 du=sum\_j,kin I\_T(b\_T-|x\textasciicircum{}(1)\_j-x\textasciicircum{}(1)\_k|)\_+<< b\_T no.I\_T .}
\expandafter\gdef\csname CGTip@eq:close_pair\endcsname{Eq. (8.6), p. 33: |y\_k+m-y\_k|>=(m-1)b .}
\expandafter\gdef\csname CGTip@lem:sampling_rho\endcsname{Lemma 8.5 (Sampling at points of bounded multiplicity) (p. 33): Let x\_k, kin I\_T, be real points of [T-2T\textasciicircum{}1/2,2T+2T\textasciicircum{}1/2] with multiplicity function tilde n(u)=no.kin I\_T:u- b\_T<= x\_k<= u satisfying int\_R tilde n\textasciicircum{}2<= k b\_Tno.I\_T for some k>=1. Let r>=1 be an integer, 2<= X<= T\textasciicircum{}1/(4r), varphi a cutoff equal to 1 on [T-2T\textasciicircum{}1/2,2T+2T\textasciicircum{}1/2] as in Lemma 7.2, and g= M(varphi S'\_X). Then, wit...}
\expandafter\gdef\csname CGTip@thm:second_step\endcsname{Theorem 8.6 (Second Newton-type step) (p. 34): Let H\_0=log log T, H\_1=log H\_0, H\_2=log H\_1, X\_1=T\textasciicircum{}H\_0\textasciicircum{}-1/2, X\_2=T\textasciicircum{}H\_1\textasciicircum{}-1/2, and let x\textasciicircum{}(2)\_k be given by (8.1). Then <L(x\textasciicircum{}(1)\_k)\textasciicircum{}2e\_2,k\textasciicircum{}2>\_I\_T=(1)/(4pi\textasciicircum{}2)H\_2+O(1+sqrt(H\_2)),; (1)/(no.I\_T)sum\_kin I\_T(x\textasciicircum{}(2)\_k-gamma\_k)\textasciicircum{}2=(log log log log T+O(sqrt(log log log log T)))/(log\textasciicircum{}2T).}
\expandafter\gdef\csname CGTip@eq:second_step\endcsname{Eq. (8.7), p. 34: <L(x\textasciicircum{}(1)\_k)\textasciicircum{}2e\_2,k\textasciicircum{}2>\_I\_T=(1)/(4pi\textasciicircum{}2)H\_2+O(1+sqrt(H\_2)),; (1)/(no.I\_T)sum\_kin I\_T(x\textasciicircum{}(2)\_k-gamma\_k)\textasciicircum{}2=(log log log log T+O(sqrt(log log log log T)))/(log\textasciicircum{}2T).}
\expandafter\gdef\csname CGTip@eq:e1_fourth\endcsname{Eq. (8.8), p. 34: < e\_1,k\textasciicircum{}4>\_I\_T<<(H\_1\textasciicircum{}2)/(log\textasciicircum{}4T).}
\expandafter\gdef\csname CGTip@subsec:further_steps\endcsname{Section 8.5: All further steps: the exact identity, the local deformation, and the multiplicity (p. 35)}
\expandafter\gdef\csname CGTip@lem:second_identity\endcsname{Lemma 8.7 (The second-level identity) (p. 35): For j>=1, with y\textasciicircum{}(j)\_k=tau\_k-x\textasciicircum{}(j)\_k and chi''\_j,k=int\_x\textasciicircum{}(j-1)\_k\textasciicircum{}tau\_k(L(tau\_k)-L(y)) dy, y\textasciicircum{}(j)\_k=(1-(L(tau\_k))/(L(x\textasciicircum{}(j-1)\_k)))y\textasciicircum{}(j-1)\_k+(S\_X\_j(x\textasciicircum{}(j-1)\_k)+chi''\_j,k)/(L(x\textasciicircum{}(j-1)\_k)),; henceL(tau\_k)y\textasciicircum{}(j)\_k=S\_X\_j(x\textasciicircum{}(j-1)\_k)+ r\_j,k, with | r\_j,k|<< T\textasciicircum{}c\_j-1 if |y\textasciicircum{}(j-1)\_k|<= T\textasciicircum{}c\_j-1/2 and c\_j-1<= c\_j (true by (2.11) and i...}
\expandafter\gdef\csname CGTip@eq:Y_identity\endcsname{Eq. (8.9), p. 35: y\textasciicircum{}(j)\_k=(1-(L(tau\_k))/(L(x\textasciicircum{}(j-1)\_k)))y\textasciicircum{}(j-1)\_k+(S\_X\_j(x\textasciicircum{}(j-1)\_k)+chi''\_j,k)/(L(x\textasciicircum{}(j-1)\_k)),; henceL(tau\_k)y\textasciicircum{}(j)\_k=S\_X\_j(x\textasciicircum{}(j-1)\_k)+ r\_j,k,}
\expandafter\gdef\csname CGTip@prop:deformation\endcsname{Proposition 8.8 (Local deformation of the second clock) (p. 35): Let X\_1=T\textasciicircum{}c\_1, X\_2=T\textasciicircum{}c\_2 with c\_1=H\_0\textasciicircum{}-1/2, c\_2=H\_1\textasciicircum{}-1/2, and d\_j,k=x\textasciicircum{}(j)\_k+1-x\textasciicircum{}(j)\_k (d\_0,k=tau\_k+1-tau\_k). Then, for T>= T\_0, <|(d\_2,k-d\_0,k)/(d\_0,k)|\textasciicircum{}4>\_I\_T<< c\_2\textasciicircum{}4=(1)/(H\_1\textasciicircum{}2), (1)/(no.I\_T)no.kin I\_T:|(d\_2,k)/(d\_0,k)-1|>alpha <<(1)/(alpha\textasciicircum{}4H\_1\textasciicircum{}2)(0<alpha <=1), in particular no.kin I\_T:d\_2,k<=0<<no.I\_T/(log log log T)\textasciicircum{}2:...}
\expandafter\gdef\csname CGTip@eq:deformation\endcsname{Eq. (8.10), p. 35: <|(d\_2,k-d\_0,k)/(d\_0,k)|\textasciicircum{}4>\_I\_T<< c\_2\textasciicircum{}4=(1)/(H\_1\textasciicircum{}2), (1)/(no.I\_T)no.kin I\_T:|(d\_2,k)/(d\_0,k)-1|>alpha <<(1)/(alpha\textasciicircum{}4H\_1\textasciicircum{}2)(0<alpha <=1),}
\expandafter\gdef\csname CGTip@lem:zero_pairs\endcsname{Lemma 8.9 (Near pairs and triples of ordinates) (p. 36): Let b= b\_T, M=no.I\_T and, for R>=1 and real u, n\_ R(u)=no.lin I\_T:|gamma\_l-u|<= R b. There are absolute constants C\_1,C\_2 such that for T>= T\_0 and every R>=1, no.(k,l)in I\_T\textasciicircum{}2:|gamma\_k-gamma\_l|<= R b<= C\_1 R M,sum\_kin I\_Tn\_ R(gamma\_k)\textasciicircum{}2<= C\_2 R\textasciicircum{}2 M .}
\expandafter\gdef\csname CGTip@eq:zero_pairs\endcsname{Eq. (8.11), p. 36: no.(k,l)in I\_T\textasciicircum{}2:|gamma\_k-gamma\_l|<= R b<= C\_1 R M,sum\_kin I\_Tn\_ R(gamma\_k)\textasciicircum{}2<= C\_2 R\textasciicircum{}2 M .}
\expandafter\gdef\csname CGTip@thm:multiplicity_all\endcsname{Theorem 8.10 (Multiplicity of the corrected clock points) (p. 36): Let j>=0 and let x\textasciicircum{}(j)\_k be given by (8.1) with cutoffs X\_i=T\textasciicircum{}c\_i, c\_i<=(1)/(2), so that x\textasciicircum{}(j)\_kin[T-2T\textasciicircum{}1/2,2T+2T\textasciicircum{}1/2] for kin I\_T and T>= T\_0 (Lemma 8.7). Let tilde n\_j(u)=no.kin I\_T:u- b\_T<= x\textasciicircum{}(j)\_k<= u, and let A>=1, H>=1 be such that ||L(tau\_k)e\_j,k||\_4,I\_T<= Asqrt() H. Then, with an absolute constant C\_3, int\_R t...}
\expandafter\gdef\csname CGTip@eq:multiplicity_all\endcsname{Eq. (8.12), p. 36: int\_R tilde n\_j(u)\textasciicircum{}2 du<= C\_3 Asqrt() H b\_T no.I\_T .}
\expandafter\gdef\csname CGTip@eq:cutoff_law\endcsname{Eq. (8.13), p. 37: H\_j=2log(1)/(c\_j)(j>=1), c\_j=C\_* H\_j-1\textasciicircum{}-9/16(j>=3),}
\expandafter\gdef\csname CGTip@thm:uniform_iteration\endcsname{Theorem 8.11 (All Newton-type steps, uniformly) (p. 37): There are absolute constants C\_*, H\textasciicircum{}*, C and T\_0 such that, with the cutoffs (8.13), for every T>= T\_0 and every 1<= m<= m\textasciicircum{}*, <L(tau\_k)\textasciicircum{}2e\_m,k\textasciicircum{}2>\_I\_T=(H\_m)/(4pi\textasciicircum{}2)+vartheta\_m C(1+sqrt(H\_m)),|vartheta\_m|<=1, and ||L(tau\_k)e\_j,k||\_4,I\_T<= Csqrt(H\_j) for 1<= j<m\textasciicircum{}*. In particular, at m=m\textasciicircum{}*=log\textasciicircum{}*T+O(1), (1)/(no.I\_T)sum\_kin...}
\expandafter\gdef\csname CGTip@eq:uniform_iteration\endcsname{Eq. (8.14), p. 37: <L(tau\_k)\textasciicircum{}2e\_m,k\textasciicircum{}2>\_I\_T=(H\_m)/(4pi\textasciicircum{}2)+vartheta\_m C(1+sqrt(H\_m)),|vartheta\_m|<=1,}
\expandafter\gdef\csname CGTip@eq:final_clock\endcsname{Eq. (8.15), p. 37: (1)/(no.I\_T)sum\_kin I\_T(x\textasciicircum{}(m\textasciicircum{}*)\_k-gamma\_k)\textasciicircum{}2<=(C')/(log\textasciicircum{}2T):}
\expandafter\gdef\csname CGTip@eq:uniform_recursion\endcsname{Eq. (8.16), p. 38: A\_j(r)<= C(sqrt(r)+(r\textasciicircum{}2)/(sqrt(H\_j)))+(C\_5sqrt(r))/(sqrt(H\_j))max (A\_j-1(2r),A\_j-1(2r)\textasciicircum{}9/8),}
\expandafter\gdef\csname CGTip@eq:induction_claim\endcsname{Eq. (8.17), p. 38: A\_j(r)<= B\_* r\textasciicircum{}2(0<= j<= m-1, 1<= r<=2\textasciicircum{}m-j).}
\expandafter\gdef\csname CGTip@cor:third_step\endcsname{Corollary 8.12 (Third step) (p. 39): With c\_3=C\_* H\_2\textasciicircum{}-9/16 and H\_3=(9)/(8)log H\_2+2log(1/C\_*), [ <L(tau\_k)\textasciicircum{}2e\_3,k\textasciicircum{}2>\_I\_T=(H\_3)/(4pi\textasciicircum{}2)+O(1+sqrt(H\_3)),(1)/(no.I\_T)sum\_kin I\_T(x\textasciicircum{}(3)\_k-gamma\_k)\textasciicircum{}2=(H\_3+O(sqrt(H\_3)))/(log\textasciicircum{}2T), ] and the multiplicity of the second clock satisfies int tilde n\_2\textasciicircum{}2llsqrt(H\_2) b\_Tno.I\_T.}
\expandafter\gdef\csname CGTip@cor:clock_tails\endcsname{Corollary 8.13 (Tails of the corrected clocks) (p. 39): With the cutoffs (8.13), for T>= T\_0, every 1<= j<m\textasciicircum{}*, every integer 1<= q<=2\textasciicircum{}m\textasciicircum{}*-j and every lambda>0, no.kin I\_T:L(tau\_k)|e\_j,k|>lambdasqrt(H\_j)<=((B\_* q\textasciicircum{}2)/(lambda))\textasciicircum{}2qno.I\_T . For fixed j and q the right side is <<\_qlambda\textasciicircum{}-2qno.I\_T for T>= T\_0(j,q), since 2\textasciicircum{}m\textasciicircum{}*-j->inf: for every fixed q, and T large in terms of q, th...}
\expandafter\gdef\csname CGTip@eq:clock_tails\endcsname{Eq. (8.18), p. 39: no.kin I\_T:L(tau\_k)|e\_j,k|>lambdasqrt(H\_j)<=((B\_* q\textasciicircum{}2)/(lambda))\textasciicircum{}2qno.I\_T .}
\expandafter\gdef\csname CGTip@lem:sorted\endcsname{Lemma 8.14 (Sorting the final points) (p. 40): Let gamma\_1<=...<=gamma\_ M, let z\_1,...,z\_ M be real, and let z\textasciicircum{}increasing to be the same values in nondecreasing order. Then sum\_k|z\textasciicircum{}increasing to\_k-gamma\_k|\textasciicircum{}q<=sum\_k|z\_k-gamma\_k|\textasciicircum{}q for every q>=1, and z\textasciicircum{}increasing to\_k=z\_k as multisets. In particular, sorting the points x\textasciicircum{}(j)\_k, kin I, of one level does not increa...}
\expandafter\gdef\csname CGTip@cor:sorted_clock\endcsname{Corollary 8.15 (A nondecreasing clock built without ordinates) (p. 40): Let I\_T\textasciicircum{}tau=k:T<=tau\_k<2T, a set of consecutive indices defined by the clock alone. The points x\textasciicircum{}(m\textasciicircum{}*)\_k of Theorem 8.11 are defined by (8.1) for every index k with tau\_k>= T\_0, in particular for kin I\_T\textasciicircum{}tau; let x\textasciicircum{}increasing to\_k, kin I\_T\textasciicircum{}tau, be these values in nondecreasing order. Then x\textasciicircum{}increasing to\_k<= x\textasciicircum{}incre...}
\expandafter\gdef\csname CGTip@sec:envelope\endcsname{Section 9: Theorem E: the pointwise envelope (p. 41)}
\expandafter\gdef\csname CGTip@prop:envelope_equivalence\endcsname{Proposition 9.1 (The envelope of the clock is the envelope of S) (p. 41): (a) limsup\_k->inf|tau\_k-gamma\_k|=2pi a. (b) If X\_k>=2 and sigma\_X\_k=o(log tau\_k), in particular if X\_k<=(log tau\_k)\textasciicircum{}2-varepsilon, then limsup\_k->inf|tau\_k\textasciicircum{}X\_k-gamma\_k|=2pi a.}
\expandafter\gdef\csname CGTip@thm:envelope\endcsname{Theorem 9.2 (Theorem E) (p. 41): Let e be continuous on [(1)/(2),1] with e(1)=0, such that for every varepsilon>0, zeta(sigma+it)<<\_varepsilon|t|\textasciicircum{} e(sigma)+varepsilon uniformly for (1)/(2)<=sigma <=1 and |t|>=2, and put e(sigma)=(1)/(2)-sigma+ e(1-sigma) for 0<=sigma<(1)/(2). Then a<=(1)/(2pi log 2)int\_0\textasciicircum{}1(e(sigma))/(sqrt(sigma(2-sigma))) dsigma ....}
\expandafter\gdef\csname CGTip@eq:envelope_formula\endcsname{Eq. (9.1), p. 41: a<=(1)/(2pi log 2)int\_0\textasciicircum{}1(e(sigma))/(sqrt(sigma(2-sigma))) dsigma .}
\expandafter\gdef\csname CGTip@eq:envelope_value\endcsname{Eq. (9.2), p. 41: C\_inf=0.0968151142790...,limsup\_k->inf|tau\_k-gamma\_k|=2pi a<=2pi C\_inf=0.608307303... ;}
\expandafter\gdef\csname CGTip@rem:envelope_inputs\endcsname{Remark 9.3 (p. 43): (a) Trudgian [ref, S 5 of the preprint] obtained a<=0.1013 from Huxleytextquoteright s exponent (32)/(205) at sigma=(1)/(2) by the same method; Theorem 9.2 replaces the exponents and lets the parameters of the disc tend to their limits. With the explicit exponents used by Bellotti and Wong ((27)/(164) at sigma=(1)/(...}
\expandafter\gdef\csname CGTip@sec:conclusions\endcsname{Section 10: Conclusions (p. 43)}
\expandafter\gdef\csname CGTip@q:constant\endcsname{Question 10.1 (p. 44): For fixed cin(0,(1)/(2)], do the bounded means of u\_k(T\textasciicircum{}c)\textasciicircum{}2 converge, and to (2pi\textasciicircum{}2)\textasciicircum{}-1log(1/c) plus which constant?}
\expandafter\gdef\csname CGTip@q:fixed_c\endcsname{Question 10.2 (p. 44): Does < r\_k(T\textasciicircum{}c)\textasciicircum{}2>\_I\_T stay bounded for fixed c, that is (Remark 7.7), is < t\_k(T\textasciicircum{}c)\textasciicircum{}2>\_I\_T bounded? Is the error of the corrected clock after log\textasciicircum{}*T+O(1) steps bounded in units of zeros for almost every index, and not only in mean square? Can Theorems 5.1 and 7.3 be extended to c>(1)/(2)?}
\expandafter\gdef\csname CGTip@q:explicit\endcsname{Question 10.3 (p. 44): Can the absolute constants C\_*, H\textasciicircum{}* and B\_* of Theorem 8.11, and those of Fujii's and Montgomery--Vaughan's bounds on which they rest, be made explicit, so that the cutoff law (8.13) becomes a formula that can be evaluated at a given height?}
\makeatother

\hypersetup{
  colorlinks=false,
  pdfborder={0 0 1},
  linkbordercolor={1 0 0},
  citebordercolor={0 0.3 1},
  urlbordercolor={0 0.6 0.6},
  filebordercolor={0 0.6 0.6},
  linktoc=all,
  bookmarksnumbered=true,
  bookmarksopen=true,
  bookmarksopenlevel=1,
  pdfstartview=FitH,
  pdftitle={A computable prime-corrected Riemann-von Mangoldt clock: an exact defect identity, the prime-information threshold, and the location of the zeros of the Riemann zeta function in mean square},
  pdfauthor={Rotkiv Ayn\"asac},
  pdfsubject={Zeros of the Riemann zeta function, Riemann-von Mangoldt clock, ordinal defect, truncated Euler products, Newton-type iteration},
  pdfkeywords={Riemann zeta function, zeros of the Riemann zeta function, Riemann-von Mangoldt formula, argument of the Riemann zeta function, Euler product, prime cutoff, Newton-type iteration}
}
% Back-references in the bibliography.
\renewcommand*{\backref}[1]{}
\renewcommand*{\backrefalt}[4]{%
  \ifcase #1 %
    \relax
  \or
    \ (Cited on p.~#2.)%
  \else
    \ (Cited on pp.~#2.)%
  \fi}

\geometry{margin=1in}
\numberwithin{equation}{section}
\setcounter{tocdepth}{2}
\setcounter{secnumdepth}{3}
\setlength{\emergencystretch}{1.5em}

\theoremstyle{plain}
\newtheorem{theorem}{Theorem}[section]
\newtheorem{proposition}[theorem]{Proposition}
\newtheorem{lemma}[theorem]{Lemma}
\newtheorem{corollary}[theorem]{Corollary}
\newtheorem{conjecture}[theorem]{Conjecture}
\theoremstyle{definition}
\newtheorem{definition}[theorem]{Definition}
\newtheorem{hypothesis}[theorem]{Hypothesis}
\newtheorem{heuristic}[theorem]{Heuristic}
\newtheorem{question}[theorem]{Question}
\theoremstyle{remark}
\newtheorem{remark}[theorem]{Remark}

\newcommand{\R}{\mathbb R}
\newcommand{\NO}{N_0}
\newcommand{\lL}{\mathcal L}
\newcommand{\IT}{\mathcal I_T}
\newcommand{\sgn}{\operatorname{sgn}}
\newcommand{\Cr}{\operatorname{Cr}}
\newcommand{\BW}{\operatorname{BW}}
\newcommand{\Rg}{\mathcal R^{\gamma}}
\newcommand{\Rt}{\mathcal R^{\tau}}
\newcommand{\PP}{\mathcal P}
\newcommand{\BWS}{\operatorname{BW}_S}

\begin{document}

\begin{frontmatter}
\title{A computable prime-corrected Riemann--von Mangoldt clock: an exact defect identity, the prime-information threshold, and the location of the zeros of the Riemann zeta function in mean square}
\author{Rotkiv Ayn\"asac}

\begin{abstract}
Let $\gamma_k$ be the ordinates of the nontrivial zeros of the Riemann zeta function, counted with multiplicity, and let $\tau_k$ be the midpoint Riemann--von Mangoldt clock, $\NO(\tau_k)=k-\frac12$. The clock is close to the zeros in height, but the ordinal displacement $s_k=k-\frac12-\NO(\gamma_k)$ has mean square $\frac1{2\pi^2}\log\log T$ (Selberg, Korolev): the clock misplaces almost every zero by an unbounded number of positions. We prove, unconditionally: (A) an exact identity that turns sums of $\mathcal G(u_k+\frac12)-\mathcal G(u_k-\frac12)$, $u_k=k-\frac12-F(\gamma_k)$, over consecutive ordinates into $dF$-integrals of $g(N-F)$, where $\mathcal G'=g$, for every continuous $g$ and every $C^1$ model $F$ of the counting function $N$; (B) for the clock equation corrected by the Euler phase $S_X$ of the primes $p\le X$, the mean square of the ordinal defect at the ordinates is $\frac1{2\pi^2}\log\frac{\log T}{\log X}+O(1+\sqrt{\log\frac{\log T}{\log X}})$ uniformly for $2\le X\le T^{1/2}$, so that it is bounded if and only if $\log X\asymp\log T$; (C) an exact finite-height accounting of the gain of the one-step clock $\tau_k-S_X(\tau_k)/\lL(\tau_k)$, $\lL(t)=\frac1{2\pi}\log\frac t{2\pi}$; (D) with $X=T^{1/\sqrt{\log\log T}}$ the one-step clock, computed from $\tau_k$ and the primes alone, has mean-square error $(\log\log\log T+O(\sqrt{\log\log\log T}))/\log^2T$ in height, against $(2+o(1))\log\log T/\log^2T$ for the clock, and the Newton-type step iterated $\log^*T+O(1)$ times, with cutoffs depending only on $T$ through absolute constants that are not computed, locates the ordinates with a mean-square error of $O(1)$ mean spacings; the proof rests on an exact moment recursion, on Fujii's moments of $S(t+h)-S(t)$, and on a sampling inequality for maximal functions at points of bounded multiplicity; sorting the final points gives a nondecreasing sequence with no larger mean-square error; the ordinal defect has exponential tails, the corrected count with primes up to $T^{1/(C\sqrt\lambda)}$ being off by more than $\lambda$ positions on at most $2^{2-\sqrt{\lambda/C}}$ of the zeros, and the errors of the corrected clocks have all moments bounded in units of $\sqrt{H_m}$; (E) $\limsup_k|\tau_k-\gamma_k|=2\pi\limsup_t|S(t)|/\log t$, and Backlund's method with the exponent profile generated from Bourgain's exponent pair by van der Corput's process gives $\limsup_k|\tau_k-\gamma_k|\le0.6083074$, compared with $2\pi\cdot0.10076=0.6331\ldots$ from the explicit bound of Bellotti and Wong; the new envelope is asymptotic and is not an effective finite-height estimate. No zero of the Riemann zeta function is used in the construction of the corrected clocks.
\end{abstract}

\begin{keyword}
Riemann zeta function \sep zeros of the Riemann zeta function \sep Riemann--von Mangoldt formula \sep argument of the Riemann zeta function \sep truncated Euler product \sep Newton-type iteration
\MSC[2020] 11M06 \sep 11M26
\end{keyword}
\end{frontmatter}

\tableofcontents

%======================================================================
\section{Introduction}
\label{sec:introduction}

\subsection{The Riemann--von Mangoldt clock and the ordinates}
\label{subsec:clock_intro}
Let $\zeta(s)$ be the Riemann zeta function, let $N(T)$ be the number of its nontrivial zeros $\rho=\beta+i\gamma$ with $0<\gamma\le T$, counted with multiplicity, and list their ordinates as $\gamma_1\le\gamma_2\le\cdots$, with multiplicity. Away from ordinates the Riemann--von Mangoldt formula reads \cite[\S\S9.3--9.4]{Titchmarsh1986}
\begin{equation}
 N(T)=\NO(T)+\frac1\pi\arg\zeta\Bigl(\frac12+iT\Bigr)+O\Bigl(\frac1T\Bigr),\qquad \NO(T)=\frac{T}{2\pi}\log\frac{T}{2\pi e}+\frac78,
 \label{eq:intro_rvm}
\end{equation}
the argument of the Riemann zeta function being obtained by continuous variation along $2\to2+iT\to\frac12+iT$ from $\arg\zeta(2)=0$. The \emph{midpoint Riemann--von Mangoldt clock} (RvM clock) is the sequence $\tau_k$ defined by
\begin{equation}
 \NO(\tau_k)=k-\frac12,\qquad \tau_k=\frac{2\pi\,(k-\frac{11}8)}{W_0\bigl((k-\frac{11}8)/e\bigr)}\qquad(k\ge1),
 \label{eq:intro_clock}
\end{equation}
where $W_0$ is the principal branch of the Lambert function \cite{Corless1996}; substituting $w=W_0((k-11/8)/e)$ directly in $we^w=(k-11/8)/e$ verifies the inversion. The other real branch at the negative argument for $k=1$ gives a value below $2\pi$ and is excluded by the defining condition $\tau_k>2\pi$. The clock is explicit and lies, in phase, half way between consecutive Gram points. This paper studies its relation to the ordinates $\gamma_k$: in which sense the clock describes them, and how much prime information has to be added to it. The relation is measured in units of zeros by the ordinal displacement and by its prime-corrected version, the \emph{ordinal defect},
\begin{equation}
 s_k=k-\frac12-\NO(\gamma_k),\qquad u_k(X)=s_k-S_X(\gamma_k),\qquad S_X(t)=-\frac1\pi\sum_{p\le X}\arg\bigl(1-p^{-1/2-it}\bigr).
 \label{eq:intro_defect}
\end{equation}
$u_k(X)$ measures how well the true ordinate satisfies the corrected clock equation $\NO(t)+S_X(t)=k-\frac12$; $|u_k(X)|<\frac12$ means that the corrected count assigns to $\gamma_k$ its correct index.

\subsection{The approximation of the zeros by the clock and the primes}
\label{subsec:program}
The approximation studied here starts from the explicit numbers $\tau_k$ and corrects them with the Euler phase of finitely many primes and prime powers, by one or several Newton-type steps for the corrected equation $\NO(t)+S_X(t)=k-\frac12$ from $t=\tau_k$, with cutoffs that depend only on the height. Every corrected point is a closed expression in $k$, $\tau_k$ and the primes: the construction evaluates neither the Riemann zeta function nor any of its zeros. The ordinates $\gamma_k$ enter only in the statements, to measure the error. Theorems A and C are exact identities; Theorems B and D are asymptotic statements in mean square over the zeros with $T<\gamma_k\le2T$; Theorem E concerns the worst zero.

\subsection{What is known and only used}
\label{subsec:known_intro}
The following facts are quoted with references in Section~\CGref{sec:known_results} and are not reproved.
\begin{itemize}[leftmargin=1.6em]
\item The displacement $\delta_k=\tau_k-\gamma_k$ is bounded, and $\limsup_k|\delta_k|\le2\pi\cdot0.10076=0.6331\ldots$, from the explicit bound of Bellotti and Wong for $N(T)-\frac T{2\pi}\log\frac T{2\pi e}$ \cite{BellottiWong2025}, which improves those of Trudgian \cite{Trudgian2014} and of Hasanalizade, Shen and Wong \cite{HasanalizadeShenWong2022}. For large $T$ only, Trudgian proved $|S(T)|\le0.1013\log T$ with Huxley's exponent \cite[\S5 of the preprint]{Trudgian2010}, which gives $\limsup_k|\delta_k|\le0.6365$.
\item The ordinal displacement $s_k$ differs by less than $1$ from minus the integer displacement of $\gamma_k$ with respect to the Gram intervals, whose moments are those of a normal variable with variance $\frac1{2\pi^2}\log\log T$, by Selberg's formulae as proved by Korolev \cite{KorolevIzv2012,KorolevPIM2012}. Hence $\langle s_k^2\rangle\sim\frac1{2\pi^2}\log\log T$ and $\langle\delta_k^2\rangle\sim2\log\log T/\log^2T$ over $T<\gamma_k\le2T$.
\item In $t$-mean, $S(t)=\frac1\pi\arg\zeta(\frac12+it)$ is asymptotically normal (Selberg \cite{Selberg1946a}, Tsang \cite{Tsang1984}), and Selberg's approximation of $S(t)$ by prime sums up to a small power of $T$ holds in $L^{2k}(T,2T)$ \cite{Selberg1946a,Selberg1946b}, in the form proved by Fujii \cite[\S3]{Fujii1974}, uniformly in $k$ (Lemma~\CGref{lem:selberg_uniform}), together with Fujii's moments of $S(t+h)-S(t)$ \cite{Fujii1975,Titchmarsh1986}. These concern $t$-averages, not the ordinates.
% SOURCECHECK: The former comparison with LeClair and França--LeClair is omitted: the specific original passages were not available for verification. No argument depends on that comparison.
\end{itemize}


\subsection{Comparison with the cited antecedents}
\label{subsec:comparison}
Table~\CGref{tab:antecedents} compares statements about the same quantity wherever possible. A continuous $t$-average, a defect evaluated at a true ordinate, and the error of a point computed from the primes are different objects; a result about one requires a transfer argument before it applies to another. ``Proved here'' identifies the result established by the present proof, without asserting that no earlier equivalent formulation exists. The comparison concerns the cited antecedents, not an exhaustive priority survey.

\begingroup
\footnotesize
\setlength{\tabcolsep}{4pt}
\renewcommand{\arraystretch}{1.13}
\begin{longtable}{@{}p{0.20\textwidth}p{0.36\textwidth}p{0.39\textwidth}@{}}
\caption{Cited antecedents, conclusions proved here, and the scope of the difference.}\label{tab:antecedents}\\
\toprule
Object & Established antecedent & Present result and scope\\
\midrule
\endfirsthead
\multicolumn{3}{l}{\textit{Table \thetable\ continued}}\\
\toprule
Object & Established antecedent & Present result and scope\\
\midrule
\endhead
\bottomrule
\endfoot
Bare midpoint clock $\tau_k$
& The RvM main term and its inversion are classical. Selberg--Korolev displacement moments give $\langle(\tau_k-\gamma_k)^2\rangle_{\IT}\sim2\log\log T/\log^2T$; see Remark~\CGref{rem:korolev} and \cite{KorolevIzv2012,KorolevPIM2012}.
& This is the baseline used here, not a contribution. The Lambert expression is not claimed as new; neither is the tendency of its mean-square height error to zero.\\
Prime approximation of $S(t)$
& Selberg's approximation, the Fujii moment estimates, and the form recorded in \cite[(3.6)]{FazzariGerspachMinelli2025} bound the residual of a short prime sum in continuous $t$-mean. Montgomery--Vaughan bounds its polynomial moments.
& These are external inputs. They do not on their own state an error estimate for the specific output $\tau_k-S_X(\tau_k)/\lL(\tau_k)$; the sampling and transport estimates supply that step.\\
Exact conversion of moments and other test functions
& The counting staircase and its jumps are classical. Moment and distribution formulae for Gram displacement are already available in \cite{KorolevIzv2012,KorolevPIM2012}.
& Theorems~\CGref{thm:identity_general}, \CGref{thm:identity_g} give the exact formula for any real $C^1$ model $F$, with both endpoint terms; no monotonicity of $F$ is required. This is a general identity, not a claim to have discovered the displacement law.\\
Defect at the true ordinate, $u_k(X)$
& Prime-sum approximation controls $S(t)-S_X(t)$ in $t$-mean. This is an input about the counting error, not an already constructed approximation to $\gamma_k$.
& Theorems~\CGref{thm:main}, \CGref{thm:variable_cutoff} give $\langle u_k(X)^2\rangle_{\IT}=(2\pi^2)^{-1}\log(1/c)+O(1+\sqrt{\log(1/c)})$, $c=\log X/\log T$, uniformly for $2\le X\le T^{1/2}$, including an index-dependent cutoff. The difference is the quantified statement at the ordinates and its cutoff range.\\
One explicit prime correction
& The bare clock has the preceding $2\log\log T/\log^2T$ baseline. Approximating the phase by primes is an established idea, not a new ingredient.
& Theorem~\CGref{thm:computable_clock} controls the transport to $\tau_k$. At the analytically specified cutoff, the height MSE is $(\log\log\log T+O(\sqrt{\log\log\log T}))/\log^2T$. Its ratio to the bare-clock MSE tends to zero. Constants and onset are not evaluated.\\
Finite-height accounting
& Expanding a squared error is elementary; it does not justify independent prime responses or a fitted gain.
& Theorem~\CGref{thm:accounting} keeps the diagonal, cross term, transport, and exact nonnegative curvature term. It is an identity for a finite set, not a positive-gain guarantee at every height.\\
Further corrections and order
& Newton-type correction and increasing rearrangement are standard operations. Rearrangement cannot increase a convex matching cost for two ordered real lists.
& Theorem~\CGref{thm:uniform_iteration} gives uniform control through $\log^*T+O(1)$ prescribed corrections, with a multiplicity estimate. Corollary~\CGref{cor:sorted_clock} gives a nondecreasing output with height MSE $O(\log^{-2}T)$. The construction and its transport control are the issue; the rate alone is not a priority claim. The cutoff constants are not yet numerical.\\
Tails
& Moment inequalities and the conversion of moments to tail bounds are classical.
& Theorem~\CGref{thm:tails} gives the stated $2^{2-\sqrt{\lambda/C}}$ bound with its explicit admissible range of $c$; Corollary~\CGref{cor:clock_tails} treats the constructed points. These are asymptotic bounds with unspecified constants, not certified error probabilities for the available data.\\
Worst error of the bare clock
& Bellotti--Wong \cite[Thm.~1.1, Cor.~1.5]{BellottiWong2025} give the explicit coefficient $0.10076$ for the logarithmic term of the counting and argument bounds, implying $\limsup|\tau_k-\gamma_k|\le2\pi\cdot0.10076$. Backlund's method and the exponent inputs are established tools.
& Proposition~\CGref{prop:envelope_equivalence} identifies the two limsups exactly. Theorem~\CGref{thm:envelope} evaluates the stated exponent-profile family, giving $2\pi C_\infty=0.6083073\ldots$. This comparison improves that cited asymptotic coefficient; it is not a new explicit bound valid at a given height, nor a claim to optimize over all known profiles.\\
Accuracy at a specified height
& Direct certified zero computations can locate individual zeros using evaluations of $\zeta$; this is a different computational input from the prime-only clocks considered here.
& No certified finite-height precision constant or individual-zero accuracy is obtained from the asymptotic $O$-terms here. No superiority at the available data heights is inferred from those terms.\\
\end{longtable}
\endgroup

\paragraph{Meaning of the mean-square scale.}\label{par:meansquare_scope}
The assertion $\langle|x_k-\gamma_k|^2\rangle\ll\log^{-2}T$ bounds the RMS error by a constant multiple of one mean spacing, since the local spacing is $1/\lL(T)\sim2\pi/\log T$. It does not assert an $o(1)$ fraction of a spacing, a numerical precision at a specified height, or a bound for each ordinate. Classical moment estimates already identify this natural scale: for example, \cite[\S9.25, (9.25.4)]{Titchmarsh1986} bounds the moments of every fixed order of the gaps between consecutive ordinates on the scale $1/\log T$, and \cite[(3.6)]{FazzariGerspachMinelli2025} supplies a fourth-moment prime approximation. A sublevel-integration construction is another route to this scale; its use of classical inputs does not establish priority for a particular formula. For approaching the zeros, the useful comparison is the error of the specified construction against the bare clock, together with its constants and effective range; the latter remain open when they occur only in asymptotic $O$-terms.

\subsection{What this paper adds}
\label{subsec:main_results}
Let $\IT=\{k:T<\gamma_k\le2T\}$ and let $\langle\cdot\rangle_{\IT}$ be the mean over $\IT$. Theorems A--E are unconditional deductions from the external inputs stated in Section~\CGref{sec:known_results}. Their precise statements are proved here; this description does not establish historical priority. In particular, the scale $O(\log^{-2}T)$ alone is not claimed as a new result. Theorem E applies Backlund's method in the form of \cite{Trudgian2010,Trudgian2014,BellottiWong2025} to exponents of \cite{Bourgain2017,Ivic1985}. Table~\CGref{tab:antecedents} separates the classical inputs from the conclusions for the specific corrected points.

\begin{enumerate}[label=\textbf{\Alph*.},leftmargin=2.2em]
\item \textbf{Exact defect identity} (Theorems~\CGref{thm:identity_general} and \CGref{thm:identity_g}). For every real $C^1$ model $F$ of the counting function, every continuous $g$ with primitive $\mathcal G$, and $u_k=k-\frac12-F(\gamma_k)$, the sum of $\mathcal G(u_k+\frac12)-\mathcal G(u_k-\frac12)$ over consecutive ordinates equals $\int g(N-F)\,dF$ up to two boundary terms. For $g(x)=x^{2j}$ this gives $\sum u_k^{2j}$ ($\frac1{12}$ per zero when $j=1$); for $g(x)=e^{i\xi x}$ it identifies the characteristic function of the empirical law of $u_k$ with a $t$-integral. No hypothesis on the zeros and no monotonicity of $F$ is used.
\item \textbf{How much prime information the clock needs} (Theorems~\CGref{thm:main} and \CGref{thm:variable_cutoff}). Uniformly for $2\le X\le T^{1/2}$, with $c=\log X/\log T$,
\begin{equation}
 \frac1{\#\IT}\sum_{k\in\IT}u_k(X)^2=\frac1{2\pi^2}\log\frac1c+O\Bigl(1+\sqrt{\log(1/c)}\Bigr),
 \label{eq:intro_main}
\end{equation}
and the same holds with the cutoff $X_k=\max(2,(\tau_k/2\pi)^c)$ varying with the index. A fixed set of primes leaves a mean-square defect $\sim(2\pi^2)^{-1}\log\log T$; primes up to $T^c$ keep it bounded; boundedness holds if and only if $\liminf\log X/\log T>0$ (Corollary~\CGref{cor:dichotomy}). Thus merely letting the number of primes tend to infinity does not ensure a bounded ordinal defect: the logarithm of the cutoff must remain comparable with the logarithm of the height.
\item \textbf{Exact accounting for the computable clock} (Theorem~\CGref{thm:accounting}). For the one-step clock $\tau_k^X=\tau_k-S_X(\tau_k)/\lL(\tau_k)$, the mean-square gain over any finite set of zeros equals a diagonal prime-response term minus a cross term, with no independence assumption; the residual of $\tau_k^X$ differs from $u_k(X)$ by an explicit transport term and by the exact nonnegative term $\chi_k=\frac{\gamma_k}{2\pi}(\frac{\tau_k}{\gamma_k}-1-\log\frac{\tau_k}{\gamma_k})$.
\item \textbf{A computable clock that approaches the zeros} (Theorems~\CGref{thm:computable_clock}, \CGref{thm:second_step} and \CGref{thm:uniform_iteration}). For $2\le X\le T^{1/2}$ the residual $r_k(X)=\lL(\tau_k)(\tau_k^X-\gamma_k)$ satisfies $\langle r_k(X)^2\rangle_{\IT}=\frac1{2\pi^2}\log\frac1c+O((1+\sqrt{\log(1/c)})(1+c^2\log\log T))$; the term $c^2\log\log T$ comes from the bound of Lemma~\CGref{lem:transport} for the transport between $\tau_k$ and $\gamma_k$, which is an upper bound and not the size of the transport. With $X=T^{1/\sqrt{\log\log T}}$, or with the cutoff $\tau_k^{1/\sqrt{\log\log\tau_k}}$ read off the clock itself,
\begin{equation}
\frac1{\#\IT}\sum_{k\in\IT}\bigl(\tau_k^X-\gamma_k\bigr)^2=\frac{\log\log\log T+O(\sqrt{\log\log\log T})}{\log^2T}.
\label{eq:intro_D}
\end{equation}
A second Newton-type step from $\tau_k^X$, with cutoff $T^{1/\sqrt{\log\log\log T}}$, has mean-square error $(\log\log\log\log T+O(\sqrt{\log\log\log\log T}))/\log^2T$ (Theorem~\CGref{thm:second_step}), and the $m$-th step, with the cutoffs \CGeqref{eq:cutoff_law} depending only on $T$, has error $(H_m+O(\sqrt{H_m}+1))/\log^2T$ with absolute constants, where $H_m$ is an iterated logarithm of $T$ (Theorem~\CGref{thm:uniform_iteration}); after $m^*=\log^*T+O(1)$ steps the error is $O(1)/\log^2T$, that is, $O(1)$ mean spacings in mean square. The proof rests on an exact moment recursion, on a multiplicity bound for the corrected clock points of every level (Theorem~\CGref{thm:multiplicity_all}, from Fujii's moments of $S(t+h)-S(t)$), and on a sampling inequality for maximal functions at points of bounded multiplicity. The second clock preserves the local order and spacing of the RvM clock outside a proportion $O(H_1^{-2})$ of the indices (Proposition~\CGref{prop:deformation}). Sorting the points of the final level does not increase the unweighted mean-square error, so the final clock can be taken nondecreasing (Corollary~\CGref{cor:sorted_clock}); sorting removes inversions, not coincident points. The moments of every order of the errors are bounded in units of $\sqrt{H_j}$, which gives tails $\ll_q\lambda^{-2q}$ for the $j$-th clock (Corollary~\CGref{cor:clock_tails}); and for the defect at the ordinates, with primes up to $T^{c}$ and $\exp(-(\lambda/C)^{3/2})\le c\le1/(C\sqrt\lambda)$, the corrected count is off by more than $\lambda$ positions on at most $2^{2-\sqrt{\lambda/C}}$ of the zeros (Theorem~\CGref{thm:tails}), where Chebyshev's inequality would give $\lambda^{-2}\log\lambda$.
\item \textbf{The pointwise envelope} (Proposition~\CGref{prop:envelope_equivalence} and Theorem~\CGref{thm:envelope}). $\limsup_k|\tau_k-\gamma_k|=2\pi\limsup_{t\to\infty}|S(t)|/\log t$, also for the corrected clocks with $X\le(\log T)^{2-\varepsilon}$; and for every exponent profile $\mathfrak e$ of the Riemann zeta function on $[\frac12,1]$,
\begin{equation}
\limsup_{k\to\infty}|\tau_k-\gamma_k|\le\frac1{\log2}\int_0^1\frac{\mathfrak e(\sigma)\,d\sigma}{\sqrt{\sigma(2-\sigma)}},
\label{eq:intro_E}
\end{equation}
$\mathfrak e$ being extended to $[0,\frac12)$ by the functional equation. The profiles $\mathfrak e_J$ generated from Bourgain's exponent pair by $J$ applications of van der Corput's $\mathscr A$-process give a decreasing sequence of bounds with limit $C_\infty=0.0968151142790\ldots$, so $\limsup_k|\tau_k-\gamma_k|\le2\pi C_\infty=0.6083073\ldots$ and $\limsup_t|S(t)|/\log t\le0.09681512$, compared with $0.10076$ in \cite{BellottiWong2025} and $0.1013$ in \cite{Trudgian2010} (Remark~\CGref{rem:envelope_inputs}). This is an asymptotic bound and does not replace the explicit bounds of \cite{BellottiWong2025}, valid for all $T\ge e$.
\end{enumerate}

\paragraph{Main finding} The Euler phase of the primes up to a power of the height turns the clock from an approximation in height into an approximation in position. For the bare clock the error in units of zeros, $s_k=\lL(\xi_k)\delta_k$, has mean square $\sim\frac1{2\pi^2}\log\log T$. With $X=T^c$ the error in units of zeros of the corrected clock equation has bounded mean square (Theorem B), so outside a proportion $\ll(1+\log(1/c))/A^2$ of the zeros it is at most $A$; and the computable clock, which uses only $\tau_k$ and the primes, keeps the gain up to the transport between $\tau_k$ and $\gamma_k$ (Theorem D). With $X=T^{1/\sqrt{\log\log T}}$,
\begin{equation}
\frac{\langle(\tau_k^X-\gamma_k)^2\rangle_{\IT}}{\langle(\tau_k-\gamma_k)^2\rangle_{\IT}}\sim\frac{\log\log\log T}{2\log\log T}\longrightarrow0,
\label{eq:intro_ratio}
\end{equation}
and in units of zeros the mean square drops from $\frac1{2\pi^2}\log\log T$ to $\frac1{4\pi^2}\log\log\log T$; iterating the Newton-type step brings it down to $O(1)$. The gain is a statement in mean square; for the worst zero, Proposition~\CGref{prop:envelope_equivalence} shows that no cutoff $X\le(\log T)^{2-\varepsilon}$ changes $\limsup|\tau_k-\gamma_k|$, and the bound \CGeqref{eq:intro_E} comes from the exponents of the Riemann zeta function, not from the primes.

\subsection{Degree of proof}
\label{subsec:status}
\emph{Theorem}, \emph{Proposition}, \emph{Lemma} and \emph{Corollary} are proved here unconditionally, from the external results of Section~\CGref{sec:known_results}; \emph{Question} names an open problem. No statement of this paper rests on a computation with the zeros. The only numbers that are not closed expressions are the values of the integral \CGeqref{eq:envelope_formula} for piecewise linear profiles with rational vertices, computed in closed form by the two programs named in the Data availability statement; the constants in the $O$-terms of Theorems B--D are absolute but not computed. In Theorem~\CGref{thm:uniform_iteration} the cutoffs and the number of steps depend on two such constants, $C_\star$ and $H^*$; the corrected points are explicit expressions in $k$, $\tau_k$ and the primes once these constants are fixed.

\begin{center}
\footnotesize
\begin{longtable}{p{0.40\textwidth}p{0.17\textwidth}p{0.35\textwidth}}
\caption{Main results: degree of proof and role.}\label{tab:status}\\
\toprule
result & status & role\\
\midrule
\endfirsthead
\toprule
result & status & role\\
\midrule
\endhead
\bottomrule
\endfoot
Theorems~\CGref{thm:identity_general}, \CGref{thm:identity_g} (A) & unconditional & statistics of the errors at the zeros as $t$-integrals\\
Theorems~\CGref{thm:main}, \CGref{thm:variable_cutoff} (B) & unconditional & bounded error in units of zeros iff $\log X\asymp\log T$\\
Theorem~\CGref{thm:accounting} (C) & unconditional, exact & finite-height balance of the gain of one correction\\
Theorems~\CGref{thm:computable_clock}, \CGref{thm:second_step}, \CGref{thm:uniform_iteration} (D) & unconditional & zero-free corrected points with mean-square error $O(1)$ mean spacings\\
Theorem~\CGref{thm:tails}, Corollaries~\CGref{cor:clock_tails}, \CGref{cor:sorted_clock} & unconditional & tails of the defect and of the corrected clocks; a nondecreasing clock built without ordinates\\
Proposition~\CGref{prop:envelope_equivalence}, Theorem~\CGref{thm:envelope} (E) & unconditional, asymptotic & the worst error of the clock is the extreme size of $S(t)$\\
Questions~\CGref{q:constant}--\CGref{q:explicit} & open & what remains to be proved\\
\end{longtable}
\end{center}

\subsection{Organization}
Section~\CGref{sec:notation} lists the notation and Section~\CGref{sec:known_results} the external results. Sections~\CGref{sec:identities}--\CGref{sec:computable} prove Theorems A--D, Section~\CGref{sec:iteration} iterates the Newton-type step uniformly in the number of steps, and Section~\CGref{sec:envelope} proves Theorem E. Section~\CGref{sec:conclusions} lists what is open.

%======================================================================

\section{Notation}
\label{sec:notation}

\subsection{List of symbols}
\label{subsec:symbols}
The notation below fixes the meaning of the recurring symbols. Coefficients and dummy indices explicitly introduced inside a proof are local. The letter $\mu$ is not used (it is reserved for the M\"obius function), $\pi$ is only the number $3.14159\ldots$, and prime counts are written $\#\{p\le y\}$. Dummy variables of integration ($t$, $y$) are local.

\begin{center}
\small
\begin{longtable}{p{0.25\textwidth}p{0.67\textwidth}}
\toprule
symbol & meaning\\
\midrule
\endfirsthead
\toprule
symbol & meaning\\
\midrule
\endhead
\bottomrule
\endfoot
$\zeta(s)$ & the Riemann zeta function\\
$\rho_k=\beta_k+i\gamma_k$ & nontrivial zeros of the Riemann zeta function, $0<\gamma_1\le\gamma_2\le\cdots$ with multiplicity\\
$\nu(\gamma)$ & number of zeros of the Riemann zeta function with ordinate $\gamma$ (multiplicities included)\\
$N(t)$, $N(t^\pm)$ & number of zeros with $0<\gamma\le t$; one-sided limits\\
$\vartheta$, $\Theta_0$, $\varpi$ & Riemann--Siegel theta function; $\Theta_0(t)=\frac t2\log\frac t{2\pi e}-\frac\pi8$; $\varpi=\vartheta-\Theta_0$\\
$g_m$ & Gram points, $\vartheta(g_m)=m\pi$\\
$\NO$, $\lL$ & $\NO=\Theta_0/\pi+1$; $\lL=\NO'$, $\lL(t)=\frac1{2\pi}\log\frac t{2\pi}$\\
$S$, $S^*$, $G$ & $S(t)=\frac1\pi\arg\zeta(\frac12+it)$; midpoint value at ordinates; $G=N-\NO=S+\varpi/\pi$\\
$\IT$, $\langle\,\cdot\,\rangle_I$ & $\{k:T<\gamma_k\le2T\}$; mean over $k$ in a finite index set $I$\\
$\tau_k$, $\tau_x$, $W_0$ & RvM clock \CGeqref{eq:clock}; the same for real $x\ge1$; principal Lambert function\\
$s_k$, $\delta_k$, $\xi_k$ & $s_k=k-\frac12-\NO(\gamma_k)$; $\delta_k=\tau_k-\gamma_k$; mean-value point in $s_k=\lL(\xi_k)\delta_k$\\
$\Delta_k$ & integer displacement of $\gamma_k$ with respect to the Gram intervals (Theorem~\CGref{thm:known_korolev})\\
$\nabla_k$ & forward difference $\nabla_ka=a_{k+1}-a_k$ of a sequence (Proposition~\CGref{prop:deformation})\\
$\mathfrak d_k$ & $s_k+\Delta_k$, $|\mathfrak d_k|<1$ (Remark~\CGref{rem:korolev})\\
$\mathfrak s_T$ & $(\frac1{2\pi^2}\log\log T)^{1/2}$ (Remark~\CGref{rem:korolev})\\
$\Lambda$, $a_n$ & von Mangoldt function; $a_n=-\Lambda(n)/(\pi\sqrt n\log n)$\\
$A$, $A_X$ & set of prime powers with $\sum_{n\in A}|a_n|<\infty$; $A_X=\{p^m:p\le X\}$\\
$S_A$, $S_X$ & $\sum_{n\in A}a_n\sin(t\log n)$; $S_{A_X}$\\
$v_A$, $v_X$ & $\frac12\sum_{n\in A}a_n^2$; $v_{A_X}$\\
$\sigma_X$, $M_X$ & bounds for $|S_X|$ and $|S_X'|$, \CGeqref{eq:sigma_M_def}\\
$F$, $F_X$ & a real $C^1$ model of $N$; $F_X=\NO+S_X$\\
$u_k$, $u_k(X)$ & $k-\frac12-F(\gamma_k)$; $k-\frac12-F_X(\gamma_k)$\\
$\Psi_j$, $\kappa_j$ & sawtooth polynomial \CGeqref{eq:Psi_def}; $\kappa_j=\frac{(2j)!}{j!(2\pi)^{2j}}$\\
$X$, $c$, $C$, $T_0$ & prime cutoff; $c=\log X/\log T$; absolute constants\\
$Y$, $Y_0$, $Y_k$, $S^{\rm pr}_Y$ & polynomial length; $Y_0=T^{1/160}$; $Y_k=T^{1/(80k)}$; $S^{\rm pr}_Y(t)=-\frac1\pi\sum_{p<Y}p^{-1/2}\sin(t\log p)$\\
$x_{\rm F}$, $a_2$, $\kappa_{\rm F}$ & Fujii's length $T^{a_2/(60k)}$ and parameter (Theorem~\CGref{thm:known_fujii_selberg}); $\kappa_{\rm F}=1/(300\log2)$ (Lemma~\CGref{lem:selberg_uniform})\\
$Y_{\rm K}$, $\varepsilon_{\rm K}$, $c_{\rm K}$ & length and constants in Korolev's lemma (Remark~\CGref{rem:korolev_uniform})\\
$\mathcal J$, $V$, $\mathcal W$, $\mathcal W_1$ & primes between $X$ and $Y_0$; $V=\sum_{p\in\mathcal J}p^{-1}$; $\mathcal W=S^{\rm pr}_{Y_0}-S_X$; its part over $\mathcal J$ \CGeqref{eq:D_def}\\
$E_0$, $E_X$, $\mathcal E_2$ & $S-S^{\rm pr}_{Y_0}$; $S-S_X$; $\int_T^{2T}E_X^2\,dt$\\
$z_n$, $\mathcal D$, $\mathcal D_1$, $\mathcal D_2$, $\eta_n$, $\eta'_n$ & coefficients and Dirichlet polynomials in Lemmas~\CGref{lem:trig_MV}--\CGref{lem:D_moments}\\
$\tau_k^X$, $\delta_k^X$, $r_k(X)$, $\chi_k$ & one-step clock; $\tau_k^X-\gamma_k$; $\lL(\tau_k)\delta_k^X$; second-order error \CGeqref{eq:one_step_decomposition}\\
$\Rg_n(I)$, $\Rt_n(I)$ & prime-power responses at the zeros and at the clock points \CGeqref{eq:two_responses}\\
$\Cr_A(I)$, $D_X(I)$ & cross term \CGeqref{eq:Cr_def}; diagonal term \CGeqref{eq:full_factor_gain}\\
$\lambda$ & real multiplier of an Euler phase\\
$\varsigma_n$, $\mathbf Q$, $\mathbf b$, $\boldsymbol\lambda$, $\Upsilon_{\boldsymbol\lambda}$ & $\varsigma_n(k)=\sin(\tau_k\log n)$; Gram matrix; projection vector; coefficients; trigonometric correction (Remark~\CGref{rem:matrix_projection})\\
$t_k^X$ & root of $F_X(t)=k-\frac12$ (Remark~\CGref{rem:localization})\\
$\ell_n$ & $\log n$\\
$\psi$, $t_1$, $t_2$ & test function and endpoints in Lemma~\CGref{lem:first_moment_identity}\\
$y_0$, $w_T$ & fixed threshold; threshold tending to infinity (Corollary~\CGref{cor:dichotomy})\\
$\mathfrak g$, $\mathfrak G$, $\omega$ & continuous test function and a primitive (Theorem~\CGref{thm:identity_g}); real frequency variable\\
$\mathfrak t_k(X)$, $y_k$ & transport term $S_X(\gamma_k)-S_X(\tau_k)$; one-step correction $S_X(\tau_k)/\lL(\tau_k)$\\
$\varphi$, $\mathfrak h$, $\mathcal M$, $\mathfrak b_T$, $\mathcal U_k$ & cutoff function; $\varphi S_X'$; centred Hardy--Littlewood maximal operator; $2\pi/\log(3T)$; sampling intervals (Lemma~\CGref{lem:transport})\\
$X_k$, $\tilde u_k$, $X_\pm$, $V_\pm$ & variable cutoff; its defect; extreme cutoffs in a window; $\sum_{X_-<p\le X_+}p^{-1}$ (Theorem~\CGref{thm:variable_cutoff})\\
$\mathcal Q$, $\mathfrak w$ & maximal partial prime sum; smooth window (Lemma~\CGref{lem:maximal_primes})\\
$\mathcal I_T^\tau$, $\tau_k^{\rm loc}$ & $\{k:T\le\tau_k<2T\}$; one-step clock with cutoff $\tau_k^{1/\sqrt{\log\log\tau_k}}$ (Corollary~\CGref{cor:local_clock})\\
$\mathfrak a$, $\BW$, $\BWS$ & $\limsup_{t\to\infty}|S(t)|/\log t$; Bellotti--Wong bounds for $|N(t)-\frac t{2\pi}\log\frac t{2\pi e}|$ and for $|S(t)|$ \CGeqref{eq:BWS_def}\\
$\mathfrak e$, $(\mathsf k,\mathsf l)$ & exponent profile of the Riemann zeta function; exponent pair (Section~\CGref{sec:envelope})\\
$\sigma_\circ$, $\varrho$, $\varkappa$, $\sigma^\sharp$, $\delta^\sharp$ & centre and radius of the Jensen disc; power in Backlund's function; parameters of Backlund's trick (Theorem~\CGref{thm:known_backlund})\\
$x^{(j)}_k$, $e_{j,k}$, $\chi_{j,k}$, $\|\cdot\|_{r,I}$ & iterated Newton points \CGeqref{eq:newton_iteration}; their errors; curvature terms; discrete norms (Section~\CGref{sec:iteration})\\
$H_j$, $c_j$, $y_k$, $\tilde n_j$ & $H_0=\log\log T$, $H_j=2\log(1/c_j)$ and the cutoff law \CGeqref{eq:cutoff_law}; $y_k=S_{X_1}(\tau_k)/\lL(\tau_k)$; multiplicity functions (Lemma~\CGref{lem:multiplicity}, Theorem~\CGref{thm:multiplicity_all})\\
$A_j(r)$, $m^*$, $C_\star$, $H^*$, $B_\star$, $C_5$, $C_6$ & normalized moments of the errors, number of steps and constants of Theorem~\CGref{thm:uniform_iteration}\\
$n_{\mathsf R}(u)$, $C_1,C_2,C_3$ & number of ordinates of the window within $\mathsf R\mathfrak b_T$ of $u$; constants of Lemma~\CGref{lem:zero_pairs} and Theorem~\CGref{thm:multiplicity_all}\\
$\mathsf M$, $\mathsf P$, $\mathsf F$, $\mathsf s$, $\mathsf W$, $\mathsf y$ & local names in proofs: $\#\IT$; number of close pairs; $g^{4r}$ (Lemma~\CGref{lem:sampling_rho}); window and scale (Lemma~\CGref{lem:maximal_primes})\\
$\mathfrak k$, $\mathfrak k_j$ & multiplicity constants: least $\mathfrak k\ge1$ with $\int\tilde n^2\le\mathfrak k\,\mathfrak b_T\#\IT$ (Lemma~\CGref{lem:sampling_rho}, Theorem~\CGref{thm:uniform_iteration})\\
$y^{(j)}_k$, $\mathfrak r_{j,k}$, $d_{j,k}$ & $\tau_k-x^{(j)}_k$ and the remainder of Lemma~\CGref{lem:second_identity}; gaps $x^{(j)}_{k+1}-x^{(j)}_k$ (Proposition~\CGref{prop:deformation})\\
$\mathsf g_j$ & $\mathcal M(\varphi_jS_{X_j}')$, maximal function of the derivative of the $j$-th Euler phase (Proposition~\CGref{prop:recursion})\\
$\mathcal A_p$ & norm of the maximal operator on $L^p$ (Theorem~\CGref{thm:known_maximal})\\
$\mathscr A$ & van der Corput\textquoteright s $\mathscr A$-process on exponent pairs (Theorem~\CGref{thm:known_exponents})\\
$\varepsilon_\circ$, $\mathsf q$, $C_{\rm b}$, $C_{\rm d}$ & parameter of the Jensen disc (proof of Theorem~\CGref{thm:envelope}); fixed exponent $\log n=\mathsf q\log T$ (Remark~\CGref{rem:transport}); constants in the proof of Lemma~\CGref{lem:higher_moments}\\
$\mathfrak f_\varkappa$, $\mathcal Z_\varkappa$, $F_\circ$ & Backlund's auxiliary function, its zero sum and the majorant (Theorem~\CGref{thm:known_backlund})\\
\end{longtable}
\end{center}

\subsection{Counting functions}
\label{subsec:counting}
Let $\vartheta(t)=\Im\log\Gamma(\frac14+\frac{it}2)-\frac t2\log\pi$ for $t>0$. For $t>0$ not an ordinate \cite[Thms.~9.3--9.4]{Titchmarsh1986},
\begin{equation}
N(t)=\frac1\pi\vartheta(t)+1+S(t),\qquad S(t)=\frac1\pi\arg\zeta\bigl(\tfrac12+it\bigr)=O(\log t),
\label{eq:RvM}
\end{equation}
with $\arg\zeta$ defined by continuous variation from $2$ to $2+it$ and then horizontally to $\frac12+it$. At an ordinate we keep the right-continuous $N$ and define $S$ by \CGeqref{eq:RvM}; $S^*(\gamma)=\frac12(S(\gamma^-)+S(\gamma^+))$. Put
\begin{equation}
\Theta_0(t)=\frac t2\log\frac t{2\pi e}-\frac\pi8,\qquad \NO(t)=\frac1\pi\Theta_0(t)+1,\qquad \varpi(t)=\vartheta(t)-\Theta_0(t).
\label{eq:N0_def}
\end{equation}
By Stirling's formula $\varpi(t)=\frac1{48t}+O(t^{-3})$ and $\varpi'(t)=O(t^{-2})$; in particular $|\varpi(t)|/\pi<\frac1{100}$ for $t\ge14$. For all $t>0$,
\begin{equation}
N(t)=\NO(t)+S(t)+\frac{\varpi(t)}\pi,\qquad \lL(t)=\NO'(t)=\frac1{2\pi}\log\frac t{2\pi}.
\label{eq:N_decomposition}
\end{equation}
Since $N(2T)-N(T)=\NO(2T)-\NO(T)+O(\log T)$,
\begin{equation}
\#\IT=T\lL(T)+O(T).
\label{eq:IT_size}
\end{equation}

\subsection{The RvM clock and the displacements}
\label{subsec:clock}
$\NO$ is $C^\infty$ and strictly increasing on $(2\pi,\infty)$, with $\NO(2\pi)=-\frac18$ and $\NO(t)\to\infty$. Hence for every real $x\ge1$ the equation $\NO(t)=x-\frac12$ has exactly one solution $\tau_x>2\pi$. Writing $t=2\pi y$, the equation is $y\log(y/e)=x-\frac{11}8\ge-\frac38$; its root $y>1$ is $y=(x-\frac{11}8)/W_0((x-\frac{11}8)/e)$, with $y=e$ if $x=\frac{11}8$ \cite{Corless1996}. Thus
\begin{equation}
\tau_x=\frac{2\pi\,(x-\frac{11}8)}{W_0\bigl((x-\frac{11}8)/e\bigr)},\qquad
\frac{d\tau_x}{dx}=\frac1{\lL(\tau_x)}=\frac{2\pi}{\log(\tau_x/2\pi)},
\label{eq:clock}
\end{equation}
the derivative by implicit differentiation. Since $\Theta_0(\tau_k)=\pi(k-\frac32)$ and $\varpi=O(1/t)$, $\tau_k$ lies in phase half way between $g_{k-2}$ and $g_{k-1}$, up to $O(\tau_k^{-1}(\log\tau_k)^{-1})$.

\begin{definition}\label{def:displacements}
For $k\ge1$, $s_k=k-\frac12-\NO(\gamma_k)$ and $\delta_k=\tau_k-\gamma_k$.
\end{definition}

Since $\gamma_1>14>2\pi$ and $\tau_k>2\pi$, the mean value theorem applied to $s_k=\NO(\tau_k)-\NO(\gamma_k)$ gives
\begin{equation}
s_k=\lL(\xi_k)\,\delta_k\qquad\text{for some }\xi_k\text{ between }\gamma_k\text{ and }\tau_k .
\label{eq:s_delta}
\end{equation}

\begin{lemma}\label{lem:s_bounds}
For every $k\ge1$,
\begin{equation}
N(\gamma_k^-)-\NO(\gamma_k)+\frac12\le s_k\le N(\gamma_k^+)-\NO(\gamma_k)-\frac12,
\label{eq:s_bounds}
\end{equation}
and
\begin{equation}
\Bigl|s_k-\frac{\varpi(\gamma_k)}\pi\Bigr|\le\max\bigl(|S(\gamma_k^-)|,|S(\gamma_k^+)|\bigr)-\frac12 .
\label{eq:s_S_bound}
\end{equation}
\end{lemma}

\begin{proof}
The indices $j$ with $\gamma_j=\gamma_k$ are exactly $N(\gamma_k^-)+1,\dots,N(\gamma_k^+)$, so $N(\gamma_k^-)+1\le k\le N(\gamma_k^+)$; subtract $\frac12+\NO(\gamma_k)$. By \CGeqref{eq:N_decomposition} at $\gamma_k^\pm$ and continuity of $\NO$ and $\varpi$, \CGeqref{eq:s_bounds} reads $S^-+\frac12\le x\le S^+-\frac12$ with $x=s_k-\varpi(\gamma_k)/\pi$, $S^-=S(\gamma_k^-)$, $S^+=S(\gamma_k^+)$, and $S^+-S^-=\nu(\gamma_k)\ge1$. Put $\mathsf m=\max(|S^-|,|S^+|)$. Then $x\le S^+-\frac12\le\mathsf m-\frac12$ and $x\ge S^-+\frac12\ge-\mathsf m+\frac12$, which is \CGeqref{eq:s_S_bound}.
\end{proof}

\subsection{Euler phases and the corrected clock}
\label{subsec:euler_phase}
For $n\ge2$ let $a_n=-\Lambda(n)/(\pi\sqrt n\log n)$. For a set $A$ of prime powers with $\sum_{n\in A}|a_n|<\infty$,
\begin{equation}
S_A(t)=\sum_{n\in A}a_n\sin(t\log n),\qquad v_A=\frac12\sum_{n\in A}a_n^2 ;
\label{eq:SA_def}
\end{equation}
the series converges absolutely and uniformly in $t$, so $S_A$ is continuous and bounded. For $X\ge2$, $\sum_{n\in A_X}|a_n|\le\frac1\pi\sum_{p\le X}\sum_{m\ge1}\frac1{mp^{m/2}}<\infty$ and, since $\arg(1-z)=-\Im\sum_{m\ge1}z^m/m$ for $|z|<1$ with the principal argument,
\begin{equation}
S_X(t)=-\frac1\pi\sum_{p\le X}\arg\bigl(1-p^{-1/2-it}\bigr)=-\frac1\pi\sum_{p\le X}\sum_{m\ge1}\frac{\sin(mt\log p)}{m\,p^{m/2}}=S_{A_X}(t).
\label{eq:SX_def}
\end{equation}
Each factor has positive real part: $1-p^{-1/2}\cos(t\log p)>0$. Its argument is therefore the single continuous value $\arctan(p^{-1/2}\sin(t\log p)/(1-p^{-1/2}\cos(t\log p)))$ in $(-\pi/2,\pi/2)$. The phase in \CGeqref{eq:SX_def} is the sum of these values; reducing the argument of the entire product modulo $2\pi$ would change the function. Thus no branch is selected from the zeros or by fitting. The series of derivatives converges uniformly, because $\sum_{p\le X}\log p\sum_{m\ge1}p^{-m/2}<\infty$; hence $S_X\in C^\infty(\R)$, its derivatives may be taken termwise, and
\begin{equation}
|S_X(t)|\le\sigma_X:=\frac1\pi\sum_{p\le X}\sum_{m\ge1}\frac1{m\,p^{m/2}},\qquad |S_X'(t)|\le M_X:=\frac1\pi\sum_{p\le X}\frac{\log p}{\sqrt p-1},
\label{eq:sigma_M_def}
\end{equation}
with $\sigma_X\ll\sqrt X/\log X$ and $M_X\ll\sqrt X$ by the prime number theorem.

\begin{definition}\label{def:defect}
For $X\ge2$ let $F_X=\NO+S_X$ and $u_k(X)=k-\frac12-F_X(\gamma_k)=s_k-S_X(\gamma_k)$. The one-step clock and its residual are
\begin{equation}
\tau_k^X=\tau_k-\frac{S_X(\tau_k)}{\lL(\tau_k)},\qquad \delta_k^X=\tau_k^X-\gamma_k,\qquad r_k(X)=\lL(\tau_k)\,\delta_k^X .
\label{eq:one_step_clock}
\end{equation}
\end{definition}
Since $F_X(\tau_k)-(k-\frac12)=S_X(\tau_k)$, $\tau_k^X$ is the first step of $t\mapsto t-(F_X(t)-k+\frac12)/\lL(t)$ from $t=\tau_k$. At a simple ordinate, \CGeqref{eq:N_decomposition} gives $\NO(\gamma_k)+S^*(\gamma_k)+\varpi(\gamma_k)/\pi=k-\frac12$, so $u_k(X)=S^*(\gamma_k)-S_X(\gamma_k)+\varpi(\gamma_k)/\pi$ is the error of the prime approximation of the argument of the Riemann zeta function at the zero.

%======================================================================

\section{External results}
\label{sec:known_results}

The results of this section are known; they are used, not reproved. Where a statement is adapted (a parameter fixed, an interval split, a prime sum lengthened), the adaptation is proved after it.
% SOURCECHECK (v7): read in the original at the cited pages: Fujii 1974 (pp. 225, 226, 232--233),
% Fujii 1975 (p. 140), Titchmarsh--Heath-Brown (pp. 246, Thms 9.3, 9.4, 9.9, 9.12, 3.5, (4.12.3)),
% Korolev Publ. Inst. Math. 2012 (Lemma 2.4 p. 60, Thm 3.2 p. 67, Thm 4.2 p. 73), Bellotti--Wong,
% Montgomery--Vaughan (Cor. 3), Kinnunen ((1.2), Thm 1.4), Ivic (Lemma 2.8, (7.57)), Bourgain (Thm 6),
% Trudgian 2014 (Lemma 3), Fazzari--Gerspach--Minelli ((1.1), (3.6)), Tenenbaum (Thms 1.8, 1.10),
% Di Plinio--Lerner (arXiv version, p. 1). Not available: Fujii 1981 (Part II); Korolev 2016,
% whose Lemma 3.13 is cited through Ivic--Korolev 2018, Lemma 12 (p. 16).
% SOURCECHECK: No claim based on Turing, Selberg's proportion on the line,
% or Pratt--Robles--Zaharescu--Zeindler is used in this analytic article.


\begin{theorem}[{\cite[Thms.~9.4, 9.9 and \S9.12]{Titchmarsh1986}}, {\cite{Littlewood1924}}]\label{thm:known_basic}
$S(t)=O(\log t)$; $\gamma_{k+1}-\gamma_k$ is bounded; and $S_1(t)=\int_0^tS(y)\,dy=O(\log t)$.
\end{theorem}

\begin{theorem}[Bellotti--Wong {\cite[Thm.~1.1, Cor.~1.5 and (1.2)]{BellottiWong2025}}]\label{thm:known_BW}
For $T\ge e$, $\bigl|N(T)-\frac T{2\pi}\log\frac T{2\pi e}\bigr|\le\BW(T)$ and $|S(T)|\le\BWS(T)$, where
\begin{equation}
\begin{split}
\BW(t)&=\min\{0.10076\log t+0.24460\log\log t+8.08344,\\
&\qquad\qquad0.11200\log t+0.12567\log\log t+3.77417\},\\
\BWS(t)&=0.10076\log t+\min\{0.24460\log\log t+7.20844,\ 1.68845\log\log t+1.50956\}.
\end{split}
\label{eq:BWS_def}
\end{equation}
Moreover $|S(T)|\le2.5167$ for $0\le T\le30\,610\,046\,000$.
\end{theorem}
Earlier explicit versions are due to Trudgian \cite{Trudgian2014} and to Hasanalizade, Shen and Wong \cite{HasanalizadeShenWong2022}. For $T$ sufficiently large, Trudgian \cite[\S5 of the preprint]{Trudgian2010} proved $|S(T)|\le0.1013\log T$ with Huxley\textquoteright s exponent $\frac{32}{205}$ for $\zeta(\frac12+it)$. The bound $|S(T)|\le2.5167$ rests on Platt's rigorous isolation of the zeros up to that height \cite[(1.2)]{BellottiWong2025}.

\begin{theorem}[Selberg {\cite{Selberg1946a}}; see {\cite[p.~225]{Fujii1974}}]\label{thm:known_selberg_moments}
For every integer $j\ge1$, $\int_T^{2T}S(t)^{2j}\,dt=\kappa_j\,T(\log\log T)^j+O_j\bigl(T(\log\log T)^{j-1/2}\bigr)$.
\end{theorem}
Fujii \cite[p.~225]{Fujii1974} quotes Selberg's formula for $\int_T^{T+H}$ with $T^{1/2+a}<H<T$, $0<a<\frac14$; the form above is the sum of the cases $[T,\frac32T]$ and $[\frac32T,2T]$ with $H=\frac T2$, since $(\log\log\frac32T)^j=(\log\log T)^j+O_j((\log\log T)^{j-1}/\log T)$.

\begin{theorem}[Selberg; Korolev {\cite{KorolevIzv2012}}, {\cite[Thm.~4.2, p.~73]{KorolevPIM2012}}]\label{thm:known_korolev}
Let $\Delta_k$ be the integer such that $g_{k+\Delta_k-2}<\gamma_k\le g_{k+\Delta_k-1}$ (in \cite{KorolevPIM2012} the Gram points are $t_m=g_{m-1}$ and $t_{m-1}<\gamma_k\le t_m$ with $m=k+\Delta_k$), and let $0<\varepsilon_0<10^{-3}$. For every fixed $j\ge1$, as $n_0\to\infty$, uniformly for $n_0^{27/82+0.9\varepsilon_0}\le h_0\le n_0^{27/82+\varepsilon_0}$,
\begin{equation}
\sum_{n_0<k\le n_0+h_0}\Delta_k^{2j}=\kappa_jh_0(\log\log n_0)^j\bigl(1+o(1)\bigr),\qquad
\Bigl|\sum_{n_0<k\le n_0+h_0}\Delta_k^{2j-1}\Bigr|\ll_j h_0(\log\log n_0)^{j-1}.
\label{eq:korolev_moments}
\end{equation}
\end{theorem}

\begin{theorem}[Selberg's formula with Fujii's bounds {\cite[\S3, pp.~232--233]{Fujii1974}}]\label{thm:known_fujii_selberg}
Let $0<a_2<\frac14$ be fixed. There are $T_0$ and $C$, depending only on $a_2$, such that for $T\ge T_0$, $T^{1/2+a_2}<H<T$ and every integer $k\ge1$ with $x_{\rm F}=T^{a_2/(60k)}\ge2$,
\begin{equation}
\int_T^{T+H}\bigl(S(t)-S^{\rm pr}_{x_{\rm F}^3}(t)\bigr)^{2k}dt\le(Ck^2)^{2k}H .
\label{eq:fujii_selberg}
\end{equation}
\end{theorem}
For the Riemann zeta function (the principal character, $q=1$), \cite[p.~232]{Fujii1974} writes Selberg's formula \cite{Selberg1946b} for $x\ge2$ as $S(t)=S^{\rm pr}_{x^3}(t)+\sum_{i=1}^5O(B_i(t))$, with five explicit functions $B_i$ and absolute implied constants, and uses it with $x=T^{a_2/(60k)}$. On p.~233 it bounds the norms $(\int_T^{T+H}|B_i|^{2k})^{1/2k}$ by $k^{1/2}H^{1/2k}$ for $i=1,4$, by $k^2H^{1/2k}$ for $i=2,3$ and by at most $k^{3/2}H^{1/2k}$ for $i=5$, under the hypotheses of its Main Theorem, which for $q=1$ are $0<a_2<\frac14$ and $T^{1/2+a_2}<H<T$. By the remark on p.~226, the implied constants in that paper do not depend on $q$, $k$ and $h$. Fujii states no range for $k$; $x_{\rm F}\ge2$ is the condition under which he quotes Selberg's formula. Minkowski's inequality gives \CGeqref{eq:fujii_selberg}. Part~II \cite{Fujii1981} contains corrections to Part~I which we have not been able to consult; Remark~\CGref{rem:korolev_uniform} gives an independent route to the uniformity in $k$ used below, and Theorem~\CGref{thm:known_fujii_difference} is taken in a form stated in 1986.

\begin{lemma}[Selberg's approximation uniformly in the order]\label{lem:selberg_uniform}
Let $\kappa_{\rm F}=1/(300\log2)$ and $Y_k=T^{1/(80k)}$. There are absolute constants $C$ and $T_0$ such that for $T\ge T_0$ and every integer $k$ with $1\le k\le\kappa_{\rm F}\log T$,
\begin{equation}
\int_T^{2T}\bigl(S(t)-S^{\rm pr}_{Y_k}(t)\bigr)^{2k}dt\le(Ck^2)^{2k}\,T .
\label{eq:selberg_uniform}
\end{equation}
In particular $\int_T^{2T}(S-S^{\rm pr}_{Y_0})^4\ll T$ with $Y_0=Y_2=T^{1/160}$.
\end{lemma}

\begin{proof}
Take $a_2=\frac15$ in Theorem~\CGref{thm:known_fujii_selberg}, so that its constants are absolute, and apply it on $[T_1,T_1+H]$ with $H=\frac T2$ and $T_1\in\{T,\frac32T\}$. The hypothesis $T_1^{7/10}<\frac T2<T_1$ holds for $T\ge T_0$, and $x_{\rm F}=T_1^{1/(300k)}\ge T^{1/(300k)}\ge2$ because $k\le\kappa_{\rm F}\log T$. It remains to lengthen the prime sum from $x_{\rm F}^3$ to $Y_k$, with the same polynomial $S^{\rm pr}_{Y_k}$ on both halves. Since $T\le T_1\le2T$ and $T\ge16$, $8\le x_{\rm F}^3=T_1^{1/(100k)}\le(2T)^{1/(100k)}\le T^{1/(80k)}=Y_k$, and $S^{\rm pr}_{x_{\rm F}^3}-S^{\rm pr}_{Y_k}=\frac1\pi\Im\mathcal D$ with $\mathcal D(t)=\sum_{p\in\mathcal J'}p^{-1/2-it}$, $\mathcal J'$ the primes in $[x_{\rm F}^3,Y_k)$. As $\log Y_k/\log x_{\rm F}^3\le\frac54$, Theorem~\CGref{thm:known_mertens} gives $\sum_{p\in\mathcal J'}p^{-1}\le\log\frac54+O(1)\le V'$ with an absolute $V'$. Next, $\mathcal D^k=\sum_n\eta_nn^{-it}$ over products of $k$ primes of $\mathcal J'$, so $n\le Y_k^k=T^{1/80}$, and $\sum_n|\eta_n|^2\le k!\,V'^k$, since each $n$ has at most $k!$ ordered factorizations; Theorem~\CGref{thm:known_MV} on $[T_1,T_1+H]$ gives $\int_{T_1}^{T_1+H}|\mathcal D|^{2k}\le k!\,V'^k(H+O(T^{1/80}))\le(C'k)^kT$. Minkowski's inequality in $L^{2k}(T_1,T_1+H)$ gives $\|S-S^{\rm pr}_{Y_k}\|_{2k}\le(Ck^2+\sqrt{C'k})T^{1/2k}$, and the two halves add up to \CGeqref{eq:selberg_uniform}.
\end{proof}

\begin{remark}[An independent route to the uniformity in the order]\label{rem:korolev_uniform}
Korolev \cite[Lemma~3.13]{Korolev2016}, as quoted in \cite[Lemma~12, p.~16]{IvicKorolev2018} (we have not seen the original), proves: for fixed $0<\varepsilon_{\rm K}<10^{-3}$, $T_1\ge T_0(\varepsilon_{\rm K})$, $x_{\rm K}=T_1^{0.1\varepsilon_{\rm K}}$, $T_1^{27/82+0.9\varepsilon_{\rm K}}\le H\le T_1^{27/82+\varepsilon_{\rm K}}$, integers $1\le m\le c_{\rm K}\log x_{\rm K}$ with an absolute constant $c_{\rm K}>0$ not given numerically, and $Y_{\rm K}=x_{\rm K}^{1/(8m+3)}$,
\begin{equation}
\int_{T_1}^{T_1+H}\Bigl|\pi S(t)+\sum_{p\le Y_{\rm K}}\frac{\sin(t\log p)}{\sqrt p}\Bigr|^{2m}dt\le\frac{c_0H}{56\,\varepsilon_{\rm K}}\Bigl(\frac{\pi c_0}{\varepsilon_{\rm K}}\Bigr)^{2m},\qquad c_0=2880 .
\label{eq:korolev_lemma}
\end{equation}
For a fixed $\varepsilon_{\rm K}$ this gives \CGeqref{eq:selberg_uniform}, even with $(C\sqrt k)^{2k}$ in place of $(Ck^2)^{2k}$, for $1\le k\le0.1c_{\rm K}\varepsilon_{\rm K}\log T$, after three adaptations. (i) \emph{Cutoff.} With $m=k$, $\pi S+\sum_{p\le Y_{\rm K}}p^{-1/2}\sin(t\log p)=\pi(S-S^{\rm pr}_{Y_k})-\sum_{Y_{\rm K}<p<Y_k}p^{-1/2}\sin(t\log p)$, and $1<\log Y_k/\log Y_{\rm K}=\frac{(8k+3)\log T}{8k\,\varepsilon_{\rm K}\log T_1}\le\frac{11}{8\varepsilon_{\rm K}}$ for $T\le T_1\le3T$; so the primes between $Y_{\rm K}$ and $Y_k$ have $\sum p^{-1}\le\log\frac{11}{8\varepsilon_{\rm K}}+O(1)$, and Theorem~\CGref{thm:known_MV} on $[T_1,T_1+H]$, applicable since $Y_k^k=T^{1/80}\le H$, bounds their sum in $L^{2k}(T_1,T_1+H)$ by $C_{\varepsilon_{\rm K}}\sqrt k\,H^{1/2k}$, as in the proof of Lemma~\CGref{lem:selberg_uniform}. (ii) \emph{Short intervals.} Cover $[T,2T]$ by consecutive intervals $[T_1,T_1+H]$ with the common length $H=T^{27/82+0.95\varepsilon_{\rm K}}$, which is admissible for every $T_1\in[T,3T]$ once $T\ge T_0(\varepsilon_{\rm K})$. Each interval has its own $Y_{\rm K}=Y_{\rm K}(T_1)$, and (i) is applied on each interval before summing, so that the sum is over the same polynomial $S^{\rm pr}_{Y_k}$. (iii) \emph{Range of orders.} The condition $k\le c_{\rm K}\log x_{\rm K}=0.1c_{\rm K}\varepsilon_{\rm K}\log T_1$ follows from $k\le0.1c_{\rm K}\varepsilon_{\rm K}\log T$. Along this route $\kappa_{\rm F}$ is replaced by the unspecified absolute constant $0.1c_{\rm K}\varepsilon_{\rm K}$; Theorems~\CGref{thm:tails} and \CGref{thm:uniform_iteration} below need only some absolute constant in that place, with $300$ replaced by $1/(0.1c_{\rm K}\varepsilon_{\rm K}\log2)$ in the choice of $C_8$. Uniform bounds of the same shape at the Gram points $t_n$, for $\sum_n(S(t_n+0)+\frac1\pi\sum_{p<y}p^{-1/2}\sin(t_n\log p))^{2k}$ and for the moments of $S(t_n+0)-S(t_{n-1}+0)$, are \cite[Lemma~2.4, p.~60, and Thm.~3.2, p.~67]{KorolevPIM2012}; they concern sums over Gram points and are not used here.
\end{remark}

\begin{theorem}[Fujii {\cite[Lemma, p.~140]{Fujii1975}}, in the form of {\cite[(9.25.3), p.~246]{Titchmarsh1986}}]\label{thm:known_fujii_difference}
There are absolute constants $A$ and $T_0$ such that for every integer $k\ge1$, $T\ge T_0$ and $0\le h\le\frac12T$,
\[
\int_0^T\bigl|S(t+h)-S(t)\bigr|^{2k}dt\le T\bigl\{Ak^4\log(3+h\log T)\bigr\}^k .
\]
\end{theorem}
Fujii \cite[p.~140]{Fujii1975} announces an asymptotic formula for $\int_0^T(S(t+h)-S(t))^{2k}dt$ with main term $\frac{(2k)!}{k!(2\pi)^{2k}}T(2\log(3+h\log T))^k$ and an error term uniform in $k$; the corresponding formula on $[T,T+H]$, $T^{1/2+a_2}<H<T$, is proved in \cite[Main Theorem, p.~225]{Fujii1974}, and Titchmarsh--Heath-Brown state the upper bound above as a consequence. Only $k=1,2$ and $0<h\le1$ are used (Lemma~\CGref{lem:zero_pairs}), in the form
\begin{equation}
\int_{T-1}^{2T}\bigl(S(t+h)-S(t)\bigr)^{2k}dt\le C\,T\bigl(\log(2+h\log T)\bigr)^k\qquad(k=1,2,\ 0<h\le1,\ T\ge T_0),
\label{eq:fujii_difference}
\end{equation}
which follows from the theorem with $2T$ in place of $T$, since $\log(3+h\log2T)\le\log(4+h\log T)\le2\log(2+h\log T)$.

\begin{theorem}[Montgomery--Vaughan {\cite[Cor.~3]{MontgomeryVaughan1974}}]\label{thm:known_MV}
For complex $z_n$, real $t_0$ and $H>0$, $\int_{t_0}^{t_0+H}\bigl|\sum_{n\le Y}z_nn^{-it}\bigr|^2dt=\sum_{n\le Y}|z_n|^2\bigl(H+O(n)\bigr)$ with an absolute implied constant.
\end{theorem}
(Replace $z_n$ by $z_nn^{-it_0}$ to pass from $[0,H]$ to $[t_0,t_0+H]$.)

\begin{theorem}[Mertens {\cite{Mertens1874}}; see {\cite[Ch.~I.1, Thms.~1.8 and 1.10]{Tenenbaum2015}}]\label{thm:known_mertens}
For $y\ge2$, $\sum_{p\le y}\frac{\log p}p=\log y+O(1)$ and $\sum_{p\le y}\frac1p=\log\log y+O(1)$.
\end{theorem}

\begin{theorem}[Selberg {\cite{Selberg1946a}}; Tsang {\cite{Tsang1984}}; see {\cite[(1.1)]{FazzariGerspachMinelli2025}}]\label{thm:known_selberg_clt}
For fixed real $y_1<y_2$, $\frac1T\operatorname{meas}\{t\in[T,2T]:y_1<S(t)/\mathfrak s_T\le y_2\}\to\frac1{\sqrt{2\pi}}\int_{y_1}^{y_2}e^{-y^2/2}dy$, where $\mathfrak s_T=(\frac1{2\pi^2}\log\log T)^{1/2}$.
\end{theorem}
In \cite{FazzariGerspachMinelli2025} the statement is that $\log\zeta(\frac12+it)/\sqrt{\log\log T}$ tends to a standard complex Gaussian for $t$ uniform in $[T,2T]$; its imaginary part is $\pi S(t)/\sqrt{\log\log T}$, of variance $\frac12$.

\begin{theorem}[Hardy--Littlewood {\cite[Thm.~I.1]{Stein1970}}, as quoted in {\cite[(1.2)]{Kinnunen1997}}; Kinnunen {\cite[Thm.~1.4]{Kinnunen1997}}]\label{thm:known_maximal}
Let $\mathcal Mf(x)=\sup_{y>0}\frac1{2y}\int_{x-y}^{x+y}|f|$ be the centred maximal function of $f\in L^1_{\rm loc}(\R)$. For $1<p<\infty$ and $f\in L^p(\R)$, $\|\mathcal Mf\|_p\le\mathcal A_p\|f\|_p$, where $\mathcal A_p$ is bounded for $p\ge2$. If $f$ and $f'$ belong to $L^p(\R)$, $1<p<\infty$, then $\mathcal Mf$ has a locally absolutely continuous representative, $(\mathcal Mf)'\in L^p(\R)$, and $|(\mathcal Mf)'|\le\mathcal M(f')$ almost everywhere.
\end{theorem}

\begin{theorem}[Carleson {\cite{Carleson1966}}, Hunt {\cite{Hunt1968}}; on $\R$ see {\cite[(1.1) and p.~1]{DiPlinioLerner2014}}]\label{thm:known_carleson}
For $f\in L^2(\R)\cap L^4(\R)$ with Fourier transform $\mathcal Ff(\omega)=\int_\R f(t)e^{-i\omega t}dt$, let $\mathcal Cf(x)=\sup_{\omega_1\in\R}\bigl|\frac1{2\pi}\int_{-\infty}^{\omega_1}\mathcal Ff(\omega)e^{i\omega x}d\omega\bigr|$. Then $\|\mathcal Cf\|_4\ll\|f\|_4$.
\end{theorem}
In \cite[p.~1]{DiPlinioLerner2014} the Carleson--Hunt theorem is stated on $\R$ in the form that $f\mapsto\sup_{\xi\in\R}|\mathcal H(e^{i\xi\cdot}f)|$, $\mathcal H$ the Hilbert transform, is bounded on $L^p(\R)$ for $1<p<\infty$; since the partial Fourier integral up to $\omega_1$ is $e^{i\omega_1x}\cdot\frac12(\mathrm{Id}\pm i\mathcal H)(e^{-i\omega_1\cdot}f)(x)$, the operator $\mathcal Cf$ is at most $\frac12(|f|+\sup_\xi|\mathcal H(e^{i\xi\cdot}f)|)$.

\begin{theorem}[{\cite[Lemma~2.8 and (7.57)]{Ivic1985}}; Bourgain {\cite[\S5, Thm.~6]{Bourgain2017}}]\label{thm:known_exponents}
(a) If $(\mathsf k,\mathsf l)$ is an exponent pair, so is $\mathscr A(\mathsf k,\mathsf l)=\bigl(\frac{\mathsf k}{2\mathsf k+2},\frac12+\frac{\mathsf l}{2\mathsf k+2}\bigr)$. (b) If $(\mathsf k,\mathsf l)$ is an exponent pair with $\mathsf l-\mathsf k\ge\frac12$, then $\zeta(\mathsf l-\mathsf k+it)\ll|t|^{\mathsf k}\log|t|$ for $|t|\ge2$. (c) For every $\varepsilon>0$, $(\frac{13}{84}+\varepsilon,\frac{55}{84}+\varepsilon)$ is an exponent pair.
\end{theorem}
For (b): \cite[(7.57)]{Ivic1985} gives $\zeta(\sigma+it)\ll|t|^{(\mathsf k+\mathsf l-\sigma)/2}\log|t|$ for $\frac12\le\sigma\le\mathsf l-\mathsf k$; take $\sigma=\mathsf l-\mathsf k$.

\begin{theorem}[Phragm\'en--Lindel\"of principle in a strip {\cite[Lemma~2.5]{BellottiWong2025}}, {\cite[Lemma~3]{Trudgian2014}}]\label{thm:known_PL}
Let $\sigma_1<\sigma_2$ and $e_1\ge e_2\ge0$. If $|\zeta(\sigma_i+it)|\le C_i|t|^{e_i}\log|t|$ for $|t|\ge2$ and $i=1,2$, then $|\zeta(\sigma+it)|\ll|t|^{e(\sigma)}\log|t|$ uniformly for $\sigma_1\le\sigma\le\sigma_2$ and $|t|\ge2$, where $e$ is the linear function with $e(\sigma_i)=e_i$.
\end{theorem}
(Apply \cite[Lemma~2.5]{BellottiWong2025} to the entire function $(s-1)\zeta(s)$ of order $1$, with exponents $1+e_i$ and one power of the logarithm on the two lines.)

\begin{theorem}[Backlund's method {\cite[Lemmas~2.1, 2.3 and 2.4]{BellottiWong2025}}, after {\cite{Trudgian2014,HasanalizadeShenWong2022}}]\label{thm:known_backlund}
Let $T$ be large, $\varkappa\in\mathbb N$, $\mathfrak f_\varkappa(s)=\frac12\bigl[((s+iT-1)\zeta(s+iT))^\varkappa+((s-iT-1)\zeta(s-iT))^\varkappa\bigr]$, and for a disc $D(\sigma_\circ,\varrho)$ let $\mathcal Z_\varkappa(\sigma_\circ,\varrho)=\frac1\varkappa\sum_z\log\frac{\varrho}{|z-\sigma_\circ|}$ over the zeros $z$ of $\mathfrak f_\varkappa$ in $D(\sigma_\circ,\varrho)$. Let $1<\sigma_\circ<\varrho$, $\sigma^\sharp=\sigma_\circ+(\sigma_\circ-\frac12)^2/\varrho$, $\delta^\sharp=2\sigma_\circ-\sigma^\sharp-\frac12$, and assume $0<\delta^\sharp<\frac12$.
\begin{enumerate}[label=(\alph*),leftmargin=2em]
\item If $F_\circ:[-\pi,\pi]\to\R$ is even and $F_\circ(\theta)\ge\frac1\varkappa\log|\mathfrak f_\varkappa(\sigma_\circ+\varrho e^{i\theta})|$ for all $\varkappa$ and $\theta$, then there are $\varkappa_m\to\infty$ with $\limsup_m\mathcal Z_{\varkappa_m}(\sigma_\circ,\varrho)\le\log\frac{\zeta(\sigma_\circ)}{\sqrt{(\sigma_\circ-1)^2+T^2}\,\zeta(2\sigma_\circ)}+\frac1\pi\int_0^\pi F_\circ(\theta)\,d\theta$.
\item For every $\varkappa$, the change of $\arg\bigl((\sigma+iT-1)\zeta(\sigma+iT)\bigr)$ by continuous variation from $\sigma=\sigma^\sharp$ to $\sigma=\frac12$ is at most, in absolute value,
$\frac{\pi\mathcal Z_\varkappa(\sigma_\circ,\varrho)}{2\log(\varrho/(\sigma_\circ-\frac12))}+\frac{\pi}{100}+\frac{3\pi}{2\varkappa}+\frac\pi4$.
\end{enumerate}
\end{theorem}
In (b) the term $\frac\pi{100}$ bounds $\frac12E(T,\delta^\sharp)$ of \cite[Lemma~2.3]{BellottiWong2025}, by the monotonicity in $\delta^\sharp$ and the bound for $E(T,\frac12)$ in \cite[Lemma~2.4]{BellottiWong2025}; (a) is \cite[Lemma~2.1]{BellottiWong2025}, which also requires $\sigma_\circ-\varrho<\frac12<1<\sigma_\circ<\sigma^\sharp<\sigma_\circ+\varrho$, a consequence of $1<\sigma_\circ<\varrho$.

\begin{remark}[Known consequences for the clock]\label{rem:korolev}
(a) \emph{Moments and law of $s_k$.} Since $\vartheta(\gamma_k)/\pi\in(k+\Delta_k-2,k+\Delta_k-1]$ and $\vartheta/\pi=\NO-1+\varpi/\pi$, one has $s_k+\Delta_k=k+\Delta_k-\frac12-\NO(\gamma_k)\in[-\frac12+\frac{\varpi(\gamma_k)}\pi,\frac12+\frac{\varpi(\gamma_k)}\pi)$, so $s_k=-\Delta_k+\mathfrak d_k$ with $|\mathfrak d_k|<1$. The index set $\IT$ is the integer interval $(N(T),N(2T)]$, of length $\asymp T\log T$. Cover it by consecutive windows $(n_0,n_0+h_0]$ with $h_0=\lfloor n_0^{27/82+0.95\varepsilon_0}\rfloor$, the last one shifted to end at $N(2T)$, so that it overlaps the previous one in at most $h_0=O(T^{-1/2}\#\IT)$ indices; in each window $\log\log n_0=\log\log T+o(1)$. Summing \CGeqref{eq:korolev_moments} over the windows, and bounding the doubly counted indices by one window sum (for odd powers, by H\"older's inequality and the even moments), gives
$\langle\Delta_k^{2j}\rangle_{\IT}=\kappa_j(\log\log T)^j(1+o(1))$ and $\langle\Delta_k^{2j-1}\rangle_{\IT}=O_j((\log\log T)^{j-1})$.
Expanding $s_k^{r}=(-\Delta_k+\mathfrak d_k)^{r}$ and bounding $\langle|\Delta_k|^{r-i}\rangle_{\IT}\le\langle\Delta_k^{2\lceil r/2\rceil}\rangle_{\IT}^{(r-i)/2\lceil r/2\rceil}\ll (\log\log T)^{(r-i)/2}$ for $1\le i\le r$, one obtains for every fixed $j\ge1$
\begin{equation}
\langle s_k^{2j}\rangle_{\IT}=\kappa_j(\log\log T)^j\bigl(1+o(1)\bigr),\qquad \langle s_k^{2j-1}\rangle_{\IT}=O_j\bigl((\log\log T)^{j-1}\bigr).
\label{eq:s_moments}
\end{equation}
With $\mathfrak s_T^2=\frac{\log\log T}{2\pi^2}$, so that $\kappa_j(\log\log T)^j=\frac{(2j)!}{j!2^j}\mathfrak s_T^{2j}$, the normalized moments $\langle(s_k/\mathfrak s_T)^r\rangle_{\IT}$ therefore tend to the moments of a standard normal variable $Z$ for every $r\ge1$: the even ones by the first relation, the odd ones to $0$ because $(\log\log T)^{j-1}/\mathfrak s_T^{2j-1}\ll_j(\log\log T)^{-1/2}$. The even moments alone determine only the law of $|s_k|$; it is the odd-moment bound in \CGeqref{eq:korolev_moments} that gives the sign. The normal law is determined by its moments, so the method of moments \cite[Thm.~30.2]{Billingsley1995} shows that the empirical law of $s_k/\mathfrak s_T$ on $\IT$ tends to that of $Z$. Only consequences for $|s_k|$ are used in this paper: $\langle s_k^2\rangle_{\IT}\sim\frac1{2\pi^2}\log\log T$, $\limsup_k|s_k|=\infty$, and $\#\{k\in\IT:|s_k|\le y_0\}=o(\#\IT)$ for every fixed $y_0>0$; for these the even moments suffice.

(b) \emph{Bound for $\delta_k$.} Theorem~\CGref{thm:known_BW}, applied as $T\to\gamma_k^\pm$ (both sides have one-sided limits), gives $|N(\gamma_k^\pm)-\NO(\gamma_k)+\frac78|\le\BW(\gamma_k)$, and \CGeqref{eq:s_bounds} gives $-\BW(\gamma_k)-\frac38\le s_k\le\BW(\gamma_k)-\frac{11}8$. As $\lL$ is positive and increasing on $(2\pi,\infty)$ and $s_k=\int_{\gamma_k}^{\tau_k}\lL$,
\begin{equation}
-\frac{\BW(\gamma_k)+3/8}{\lL(\tau_k)}\le\delta_k\le\frac{\BW(\gamma_k)-11/8}{\lL(\gamma_k)}\qquad(k\ge1).
\label{eq:BW_pointwise}
\end{equation}
Hence $\delta_k=O(1)$, $\lL(\tau_k)/\lL(\gamma_k)\to1$, and $\limsup_k|\delta_k|\le2\pi\cdot0.10076=0.6331\ldots$, an immediate consequence of \cite{BellottiWong2025}.
\end{remark}

%======================================================================

\section{Theorem A: an exact defect identity at the zeros}
\label{sec:identities}

\begin{theorem}\label{thm:identity_general}
Let $1\le k_1<k_2$ be indices, $F$ a real $C^1$ function on an open interval containing $[\gamma_{k_1},\gamma_{k_2}]$, and $u_k=k-\frac12-F(\gamma_k)$ for $k_1\le k\le k_2$. For every integer $j\ge1$,
\begin{equation}
\int_{\gamma_{k_1}}^{\gamma_{k_2}}\bigl(N(t)-F(t)\bigr)^{2j}F'(t)\,dt=\sum_{k=k_1+1}^{k_2-1}\Psi_j(u_k)+\frac{(\frac12+u_{k_1})^{2j+1}+(\frac12-u_{k_2})^{2j+1}}{2j+1},
\label{eq:identity_general}
\end{equation}
where
\begin{equation}
\Psi_j(u)=\frac{(\frac12+u)^{2j+1}+(\frac12-u)^{2j+1}}{2j+1}=\sum_{l=0}^{j}\binom{2j+1}{2l}\frac{4^{l-j}}{2j+1}\,u^{2l}=u^{2j}+\sum_{l<j}\binom{2j+1}{2l}\frac{4^{l-j}}{2j+1}u^{2l}.
\label{eq:Psi_def}
\end{equation}
In particular, since $\Psi_1(u)=u^2+\frac1{12}$,
\begin{equation}
\sum_{k=k_1+1}^{k_2-1}u_k^2=\int_{\gamma_{k_1}}^{\gamma_{k_2}}\bigl(N(t)-F(t)\bigr)^2F'(t)\,dt-\frac{k_2-k_1-1}{12}-\frac{(\frac12+u_{k_1})^3+(\frac12-u_{k_2})^3}3 .
\label{eq:identity_quadratic}
\end{equation}
\end{theorem}

\begin{proof}
\emph{Existence.} On $[\gamma_{k_1},\gamma_{k_2}]$ the function $N$ is a step function with finitely many jumps (Theorem~\CGref{thm:known_basic}) and $F$, $F'$ are continuous; the integrand is bounded and piecewise continuous, hence Riemann integrable.

\emph{One gap.} Fix $k$ with $k_1\le k<k_2$. If $\gamma_k<\gamma_{k+1}$, then $N(t)=k$ for $t\in(\gamma_k,\gamma_{k+1})$, because exactly $k$ ordinates counted with multiplicity are $\le t$. The function $t\mapsto-(k-F(t))^{2j+1}/(2j+1)$ is $C^1$ with derivative $(k-F(t))^{2j}F'(t)$ (chain rule; no monotonicity of $F$ is needed), so
\begin{equation}
\begin{split}
\int_{\gamma_k}^{\gamma_{k+1}}\bigl(k-F(t)\bigr)^{2j}F'(t)\,dt&=\frac{(k-F(\gamma_k))^{2j+1}-(k-F(\gamma_{k+1}))^{2j+1}}{2j+1}\\
&=\frac{(\frac12+u_k)^{2j+1}+(\frac12-u_{k+1})^{2j+1}}{2j+1},
\end{split}
\label{eq:identity_one_gap}
\end{equation}
because $k-F(\gamma_k)=\frac12+u_k$ and $k-F(\gamma_{k+1})=-\frac12+u_{k+1}$. If $\gamma_k=\gamma_{k+1}$, the left side of \CGeqref{eq:identity_one_gap} is $0$; then $u_{k+1}=u_k+1$, and the right side equals $((\frac12+u_k)^{2j+1}+(-\frac12-u_k)^{2j+1})/(2j+1)=0$. So \CGeqref{eq:identity_one_gap} holds for every $k$, and the values of $N$ at the finitely many ordinates do not affect the integrals.

\emph{Summation and boundary terms.} Summing \CGeqref{eq:identity_one_gap} over $k_1\le k<k_2$ (a finite sum), each $k$ with $k_1<k<k_2$ contributes $(\frac12+u_k)^{2j+1}$ from the gap starting at $\gamma_k$ and $(\frac12-u_k)^{2j+1}$ from the gap ending at $\gamma_k$, that is $\Psi_j(u_k)$; the unmatched terms are $(\frac12+u_{k_1})^{2j+1}$ and $(\frac12-u_{k_2})^{2j+1}$. In the binomial expansion of $(\frac12+u)^{2j+1}+(\frac12-u)^{2j+1}$ the odd powers of $u$ cancel and $u^{2l}$ has coefficient $2\binom{2j+1}{2l}2^{2l-2j-1}$, which gives \CGeqref{eq:Psi_def}; the leading coefficient is $\binom{2j+1}{2j}/(2j+1)=1$.
\end{proof}

\begin{theorem}\label{thm:identity_g}
Let $k_1<k_2$, $F$ and $u_k$ be as in Theorem~\CGref{thm:identity_general}, and let $\mathfrak g:\R\to\mathbb C$ be continuous with a primitive $\mathfrak G$. Then
\begin{equation}
\int_{\gamma_{k_1}}^{\gamma_{k_2}}\mathfrak g\bigl(N(t)-F(t)\bigr)F'(t)\,dt=\sum_{k=k_1+1}^{k_2-1}\Bigl[\mathfrak G\bigl(u_k+\tfrac12\bigr)-\mathfrak G\bigl(u_k-\tfrac12\bigr)\Bigr]+\mathfrak G\bigl(u_{k_1}+\tfrac12\bigr)-\mathfrak G\bigl(u_{k_2}-\tfrac12\bigr).
\label{eq:identity_g}
\end{equation}
\end{theorem}

\begin{proof}
The integrand is bounded and piecewise continuous, as in the proof of Theorem~\CGref{thm:identity_general}. If $\gamma_k<\gamma_{k+1}$, then $N=k$ on $(\gamma_k,\gamma_{k+1})$ and $t\mapsto-\mathfrak G(k-F(t))$ is $C^1$ with derivative $\mathfrak g(k-F(t))F'(t)$, so the integral over $[\gamma_k,\gamma_{k+1}]$ is $\mathfrak G(k-F(\gamma_k))-\mathfrak G(k-F(\gamma_{k+1}))=\mathfrak G(u_k+\frac12)-\mathfrak G(u_{k+1}-\frac12)$. If $\gamma_k=\gamma_{k+1}$, the integral is $0$ and $u_{k+1}=u_k+1$, so the same formula holds. Summing over $k_1\le k<k_2$, each $k$ with $k_1<k<k_2$ contributes $\mathfrak G(u_k+\frac12)$ from the gap that starts at $\gamma_k$ and $-\mathfrak G(u_k-\frac12)$ from the gap that ends there.
\end{proof}

For $\mathfrak g(x)=x^{2j}$ this is \CGeqref{eq:identity_general}, since $\mathfrak G(u+\frac12)-\mathfrak G(u-\frac12)=\Psi_j(u)$.

\begin{corollary}\label{cor:identity_exp}
For real $\omega\ne0$,
\begin{equation}
\frac{2\sin(\omega/2)}{\omega}\sum_{k=k_1+1}^{k_2-1}e^{i\omega u_k}=\int_{\gamma_{k_1}}^{\gamma_{k_2}}e^{i\omega(N(t)-F(t))}F'(t)\,dt-\frac{e^{i\omega(u_{k_1}+1/2)}-e^{i\omega(u_{k_2}-1/2)}}{i\omega}.
\label{eq:identity_exp}
\end{equation}
\end{corollary}

\begin{proof}
Theorem~\CGref{thm:identity_g} with $\mathfrak g(x)=e^{i\omega x}$ and $\mathfrak G(x)=e^{i\omega x}/(i\omega)$, for which $\mathfrak G(u+\frac12)-\mathfrak G(u-\frac12)=e^{i\omega u}\cdot2\sin(\omega/2)/\omega$.
\end{proof}

\begin{corollary}[Signed law of $s_k$ from the $t$-law of $S$]\label{cor:signed_law}
For fixed real $\xi$, $\langle e^{i\xi s_k/\mathfrak s_T}\rangle_{\IT}\to e^{-\xi^2/2}$. Hence the empirical law of $s_k/\mathfrak s_T$ on $\IT$ tends to the standard normal law.
\end{corollary}

\begin{proof}
Let $\xi\ne0$, $\omega=\xi/\mathfrak s_T$, $k_1=N(T)$ and $k_2=N(2T)+1$, so that $\{k_1+1,\dots,k_2-1\}=\IT$. Apply \CGeqref{eq:identity_exp} with $F=\NO$, so $u_k=s_k$ and $N-F=G$. By Theorem~\CGref{thm:known_basic} the intervals $[\gamma_{k_1},T]$ and $[2T,\gamma_{k_2}]$ have bounded length, so the integral equals $\int_T^{2T}e^{i\omega G}\lL\,dt+O(\log T)$; the boundary term is $O(1/|\omega|)=O(\mathfrak s_T)$; and $2\sin(\omega/2)/\omega=1+O(\omega^2)$. Divide by $\#\IT=\int_T^{2T}\lL\,dt+O(\log T)$. On $[T,2T]$, $|e^{i\omega G}-e^{i\omega S}|\le|\omega||\varpi|/\pi\ll(\mathfrak s_TT)^{-1}$ and $\lL(t)=\lL(T)(1+O(1/\log T))$, whence
\[
\langle e^{i\omega s_k}\rangle_{\IT}=\frac1T\int_T^{2T}e^{i\xi S(t)/\mathfrak s_T}dt+O\Bigl(\omega^2+\frac1{\log T}\Bigr).
\]
By Theorem~\CGref{thm:known_selberg_clt} and the definition of convergence in distribution, applied to the bounded continuous function $y\mapsto e^{i\xi y}$, the integral tends to $e^{-\xi^2/2}$, and $\omega\to0$. The last statement is L\'evy's continuity theorem \cite[Thm.~26.3]{Billingsley1995}.
\end{proof}

\begin{remark}\label{rem:twelfth}
$\Psi_j(u)=\int_{u-1/2}^{u+1/2}x^{2j}dx$ is the mean of $x^{2j}$ over the unit interval centred at $u$; in particular $\Psi_1(u)=u^2+\frac1{12}$ and $\frac1{12}=\int_{-1/2}^{1/2}x^2dx$. The constant comes from the unit jump of $N-F$ at each zero, through the primitive in Theorem~\CGref{thm:identity_g}; when $F$ is not monotone, $N-F$ need not run through an interval of length one between consecutive zeros, and the identities remain exact with the integral counted with sign. With $F=\NO$, Theorem~\CGref{thm:known_selberg_moments} and induction on $j$ through the lower-degree terms of $\Psi_j$ recover the even moments in \CGeqref{eq:s_moments} without Gram intervals; Corollary~\CGref{cor:signed_law} gives the sign without Korolev's odd moments. The main purpose of Theorem~\CGref{thm:identity_general} is Theorem~\CGref{thm:main}, where $F=F_X$ need not be monotone and no Gram-interval structure is available.
\end{remark}

\begin{corollary}\label{cor:identity_clock}
For $X\ge2$ and $1\le k_1<k_2$,
\begin{equation}
\begin{split}
\sum_{k=k_1+1}^{k_2-1}u_k(X)^2={}&\int_{\gamma_{k_1}}^{\gamma_{k_2}}\Bigl(S-S_X+\frac{\varpi}\pi\Bigr)^2\bigl(\lL+S_X'\bigr)dt\\
&-\frac{k_2-k_1-1}{12}-\frac{(\frac12+u_{k_1}(X))^3+(\frac12-u_{k_2}(X))^3}3 .
\end{split}
\label{eq:identity_corrected}
\end{equation}
\end{corollary}

\begin{proof}
$F_X$ is $C^\infty$ on $(0,\infty)$ (Section~\CGref{subsec:euler_phase}), $F_X'=\lL+S_X'$, and $N-F_X=S-S_X+\varpi/\pi$ by \CGeqref{eq:N_decomposition}; apply Theorem~\CGref{thm:identity_general} with $j=1$.
\end{proof}

\begin{lemma}[First-moment identity]\label{lem:first_moment_identity}
Let $2\pi<t_1<t_2$ be points which are not ordinates and $\psi\in C^1[t_1,t_2]$, real or complex. Then
\begin{equation}
\sum_{t_1<\gamma_k\le t_2}s_k\psi(\gamma_k)=\int_{t_1}^{t_2}G(t)\psi(t)\lL(t)\,dt+\frac12\Bigl[\psi\,G^2\Bigr]_{t_1}^{t_2}-\frac12\int_{t_1}^{t_2}G(t)^2\psi'(t)\,dt .
\label{eq:first_moment_identity}
\end{equation}
\end{lemma}

\begin{proof}
On $[t_1,t_2]$ there are finitely many ordinates, none at $t_1$ or $t_2$. The function $G=N-\NO$ is right-continuous and of bounded variation, $C^\infty$ with $G'=-\lL$ between ordinates, and it jumps by $\nu(\gamma)$ at an ordinate $\gamma$. Let $G^*(\gamma)=\frac12(G(\gamma^-)+G(\gamma^+))$. The indices with $\gamma_k=\gamma$ are $N(\gamma^-)+1,\dots,N(\gamma^-)+\nu(\gamma)$, so $\sum_{\gamma_k=\gamma}(k-\frac12)=\nu(\gamma)N(\gamma^-)+\frac{\nu(\gamma)^2}2=\nu(\gamma)\frac{N(\gamma^-)+N(\gamma^+)}2$ and, $\NO$ being continuous, $\sum_{\gamma_k=\gamma}s_k=\nu(\gamma)G^*(\gamma)$. The left side of \CGeqref{eq:first_moment_identity} therefore equals the finite sum $\sum_{t_1<\gamma<t_2}\nu(\gamma)G^*(\gamma)\psi(\gamma)$ over distinct ordinates. As Lebesgue--Stieltjes measures on $(t_1,t_2]$,
\begin{equation}
d(G^2)=-2G\lL\,dt+\sum_{t_1<\gamma<t_2}2G^*(\gamma)\nu(\gamma)\,\delta_\gamma ,
\label{eq:first_moment_measures}
\end{equation}
because $(G^2)'=2GG'=-2G\lL$ between ordinates and $G(\gamma^+)^2-G(\gamma^-)^2=(G(\gamma^+)+G(\gamma^-))\nu(\gamma)=2G^*(\gamma)\nu(\gamma)$. Integrating the bounded Borel function $\psi$ against \CGeqref{eq:first_moment_measures} gives $\sum_{t_1<\gamma<t_2}\nu(\gamma)G^*(\gamma)\psi(\gamma)=\frac12\int_{(t_1,t_2]}\psi\,d(G^2)+\int_{t_1}^{t_2}G\psi\lL\,dt$. Since $\psi$ is $C^1$ and $G^2$ is right-continuous of bounded variation and continuous at $t_1$ and $t_2$, integration by parts for Lebesgue--Stieltjes integrals \cite[\S3.5]{Folland1999} gives $\int_{(t_1,t_2]}\psi\,d(G^2)=[\psi G^2]_{t_1}^{t_2}-\int_{t_1}^{t_2}G^2\psi'\,dt$. All integrands are bounded and measurable, so all integrals exist.
\end{proof}

%======================================================================
\section{Theorem B: how much prime information the clock needs}
\label{sec:main_theorem}

\begin{theorem}\label{thm:main}
There are absolute constants $T_0$ and $C$ such that for $T\ge T_0$ and every $X$ with $2\le X\le T^{1/2}$, writing $c=\log X/\log T$,
\begin{equation}
\Bigl|\frac1{\#\IT}\sum_{k\in\IT}u_k(X)^2-\frac1{2\pi^2}\log\frac1c\Bigr|\le C\Bigl(1+\sqrt{\log(1/c)}\Bigr).
\label{eq:main_theorem}
\end{equation}
\end{theorem}

Here $\log(1/c)=\log\log T-\log\log X\in[\log2,\log\log T-\log\log2]$. For fixed $c$ both the main term and the error are bounded; the theorem does not identify a limit at fixed $c$. It concerns the defect at the true ordinate and does not select a root of $F_X(t)=k-\frac12$.

\begin{corollary}[fixed cutoff]\label{cor:fixed_cutoff}
For fixed $X\ge2$, $\frac1{\#\IT}\sum_{k\in\IT}u_k(X)^2=\frac1{2\pi^2}\log\log T+O_X(\sqrt{\log\log T})$, and $\#\{k\in\IT:|u_k(X)|\le y_0\}=o(\#\IT)$ for every fixed $y_0>0$.
\end{corollary}

\begin{proof}
The first statement is \CGeqref{eq:main_theorem}. Since $|u_k(X)-s_k|\le\sigma_X$, $|u_k(X)|\le y_0$ implies $|s_k|\le y_0+\sigma_X$; apply Remark~\CGref{rem:korolev}(a).
\end{proof}

\begin{corollary}[power cutoff]\label{cor:power_cutoff}
For fixed $0<c\le\frac12$ and $X=T^c$,
\begin{equation}
\limsup_{T\to\infty}\frac1{\#\IT}\sum_{k\in\IT}u_k(T^c)^2\le\frac1{2\pi^2}\log\frac1c+C\bigl(1+\sqrt{\log(1/c)}\bigr),
\label{eq:power_cutoff}
\end{equation}
and $\#\{k\in\IT:|u_k(T^c)|>w_T\}=o(\#\IT)$ whenever $w_T\to\infty$. The same holds for $X=N(T)^c$ with fixed $0<c<\frac12$, since then $X=T^{c'}$ with $c'=c+O(\log\log T/\log T)$.
\end{corollary}

\begin{proof}
\CGeqref{eq:main_theorem} and Chebyshev's inequality.
\end{proof}

\begin{corollary}[dichotomy]\label{cor:dichotomy}
Let $X=X(T)$ with $2\le X\le T^{1/2}$ for $T\ge T_0$. Then
\[
\sup_{T\ge T_0}\frac1{\#\IT}\sum_{k\in\IT}u_k(X(T))^2<\infty\quad\text{if and only if}\quad\liminf_{T\to\infty}\frac{\log X}{\log T}>0 .
\] For fixed $X_0\ge2$ and $c\in(0,\frac12]$,
\begin{equation}
\frac1{\#\IT}\sum_{k\in\IT}\Bigl(u_k(X_0)^2-u_k(T^c)^2\Bigr)=\frac1{2\pi^2}\log\log T+O\bigl(\sqrt{\log\log T}\bigr) .
\label{eq:dichotomy}
\end{equation}
\end{corollary}

\begin{proof}
Both directions follow from \CGeqref{eq:main_theorem} (on a bounded range of $T$ the means are bounded, since $|u_k(X)|\ll\log T+\sqrt X$), because $y\mapsto\frac{y}{2\pi^2}\pm C(1+\sqrt y)$ is bounded on bounded sets and tends to infinity; \CGeqref{eq:dichotomy} is the difference of the two cases of \CGeqref{eq:main_theorem}.
\end{proof}

\begin{proposition}[Rounding the corrected count misassigns a positive proportion]\label{prop:rounding}
There are absolute constants $c_*>0$, $\kappa_0>0$ and $T_0$ such that for $T\ge T_0$ and $2\le X\le T^{c_*}$,
\[
\frac1T\operatorname{meas}\Bigl\{t\in[T,2T]:\Bigl\lfloor F_X(t)+\tfrac12\Bigr\rfloor\ne N(t)\Bigr\}\ge\kappa_0 .
\]
\end{proposition}

\begin{proof}
If $\lfloor F_X(t)+\frac12\rfloor=N(t)$ then $|N(t)-F_X(t)|\le\frac12$, i.e.\ $|E|\le\frac12$ with $E=E_X+\varpi/\pi$. Let $\mathfrak m$ be the measure of $\{t\in[T,2T]:|E|>\frac12\}$. Then $\int_T^{2T}E^2\le\frac T4+(\int_T^{2T}E^4)^{1/2}\mathfrak m^{1/2}$. By \CGeqref{eq:J_asymptotic} and \CGeqref{eq:R_fourth} (with $|\varpi|\ll1/T$), $\int_T^{2T}E^2\ge\frac T{2\pi^2}V-CT(1+\sqrt V)$ and $\int_T^{2T}E^4\le CT(1+V)^2$, so $\mathfrak m/T\ge(\frac V{2\pi^2}-C(1+\sqrt V)-\frac14)_+^2/(C(1+V)^2)$, which is $\ge\kappa_0>0$ once $V\ge V_0$, that is $c\le c_*$ by Lemma~\CGref{lem:V_size}.
\end{proof}

Rounding the corrected count to the nearest integer therefore fails to reproduce $N(t)$, for small cutoffs, on a set of $t$ of positive density: the obstruction is the size of $S-S_X$, not the sawtooth $N-\NO$. The proposition concerns the measure of a set of $t$; it does not give, without a further argument on the values at the $\gamma_k$, a positive proportion of ordinates to which rounding assigns a wrong index. For $c$ near $\frac12$ the question depends on the constant in \CGeqref{eq:main_theorem} (Question~\CGref{q:constant}), and for $c>\frac12$ nothing is proved.

\begin{remark}\label{rem:comparison_cutoff}
Since $\log\#\{p\le X\}=\log X+O(\log\log X)$, the condition of Corollary~\CGref{cor:dichotomy} says that the logarithm of the number of primes must grow proportionally to $\log T$. A fixed number of primes leaves almost all zeros violating the corrected equation by more than any fixed amount; growth of the number of primes alone is necessary but not sufficient. The range $c>\frac12$ is not covered. In Theorem~\CGref{thm:main} the cutoff is constant on each window $\IT$; Theorem~\CGref{thm:variable_cutoff} below allows the cutoff $X_k=(\tau_k/2\pi)^c$, which varies at every ordinate.
\end{remark}

\begin{remark}[Localization of the ordinates]\label{rem:localization}
If $M_X\le\frac12\lL(T)$, then $F_X'\ge\lL-M_X\ge\frac12\lL(T)$ on $[T,2T]$, so $F_X$ is strictly increasing there. For $k\in\IT$ with $F_X(T)\le k-\frac12\le F_X(2T)$ the equation $F_X(t)=k-\frac12$ has exactly one root $t_k^X\in[T,2T]$ (intermediate value theorem), and $u_k(X)=F_X(t_k^X)-F_X(\gamma_k)$ gives $|\gamma_k-t_k^X|\le2|u_k(X)|/\lL(T)$ (mean value theorem). As $M_X\asymp\sqrt X$, this applies only for $X\ll\log^2T$, and in that range the localization gains nothing in order (Proposition~\CGref{prop:monotone}). For $X=T^c$, $F_X$ need not be monotone, and Theorem~\CGref{thm:main} controls the equation at the true ordinates without selecting a root; Theorem~\CGref{thm:computable_clock} replaces the root by one Newton-type step.
\end{remark}

\subsection{Lemmas on prime polynomials}
\label{subsec:main_lemmas}
Fix $T\ge T_0$ and $2\le X\le T^{1/2}$, let $Y_0=T^{1/160}$, and split $S_X=S_X^{(1)}+S_X^{(2)}+S_X^{(\ge3)}$ according to $m=1$, $m=2$, $m\ge3$ in \CGeqref{eq:SX_def}; the three series converge absolutely and uniformly. Let $\mathcal J$ be the set of primes in $(X,Y_0)$ if $X<Y_0$ and in $[Y_0,X]$ if $X\ge Y_0$, and put
\begin{equation}
\begin{split}
&\mathcal W_1(t)=-\frac1\pi\sum_{p\in\mathcal J}\frac{\sin(t\log p)}{\sqrt p},\qquad V=\sum_{p\in\mathcal J}\frac1p,\\
&\mathcal W=S^{\rm pr}_{Y_0}-S_X=\pm\mathcal W_1-S_X^{(2)}-S_X^{(\ge3)},
\end{split}
\label{eq:D_def}
\end{equation}
with $+$ if $X<Y_0$ and $-$ if $X\ge Y_0$. The last identity holds because $S_X^{(1)}$ is the sum of $-\frac1\pi p^{-1/2}\sin(t\log p)$ over the primes $p\le X$, while $S^{\rm pr}_{Y_0}$ is the same sum over $p<Y_0$.

\begin{lemma}\label{lem:V_size}
$V=\log(1/c)+O(1)$ uniformly for $\log2/\log T\le c\le\frac12$.
\end{lemma}

\begin{proof}
By Theorem~\CGref{thm:known_mertens}, $V=|\log\log Y_0-\log\log X|+O(1)=|\log\frac1{160c}|+O(1)$. For $c\le\frac1{160}$ this is $\log\frac1c-\log160+O(1)$; for $\frac1{160}<c\le\frac12$ both $V$ and $\log\frac1c$ are $O(1)$.
\end{proof}

\begin{lemma}\label{lem:trig_MV}
For real $z_n$, $2\le n\le Y$,
\begin{equation}
\int_T^{2T}\Bigl(\sum_{n\le Y}z_n\sin(t\log n)\Bigr)^2dt=\frac T2\sum_{n\le Y}z_n^2+O\Bigl(\sum_{n\le Y}nz_n^2+\Bigl(\sum_{n\le Y}|z_n|\Bigr)^2\Bigr),
\label{eq:trig_MV}
\end{equation}
and the same with $\cos$ in place of $\sin$.
\end{lemma}

\begin{proof}
With the finite sum $\mathcal D(t)=\sum z_nn^{-it}$, $\sum z_n\sin(t\log n)=-\Im\mathcal D$ and $(\Im\mathcal D)^2=\frac12(|\mathcal D|^2-\Re\mathcal D^2)$. Theorem~\CGref{thm:known_MV} evaluates $\int|\mathcal D|^2$, and $|\int_T^{2T}(n_1n_2)^{-it}dt|\le2/\log(n_1n_2)\le2/\log4$ gives $|\int\mathcal D^2|\le(2/\log4)(\sum|z_n|)^2$. For $\cos$ use $\Re\mathcal D$ and $(\Re\mathcal D)^2=\frac12(|\mathcal D|^2+\Re\mathcal D^2)$.
\end{proof}

\begin{lemma}\label{lem:D_moments}
With absolute implied constants,
\begin{align}
\int_T^{2T}\mathcal W^2\,dt&=\frac T{2\pi^2}V+O\bigl(T(1+\sqrt V)\bigr),
\label{eq:D_second}\\
\int_T^{2T}\mathcal W^4\,dt&\ll T(1+V)^2,
\label{eq:D_fourth}\\
\int_T^{2T}S_X'(t)^2\,dt&\ll T(1+\log X)^2 .
\label{eq:SX_derivative}
\end{align}
\end{lemma}

\begin{proof}
All primes involved are $\le\max(X,Y_0)\le T^{1/2}$. Lemma~\CGref{lem:trig_MV} with $z_p=\frac1{\pi\sqrt p}$ ($p\in\mathcal J$) gives $\int\mathcal W_1^2=\frac T{2\pi^2}V+O(T^{1/2}+(\sum_{p\le T^{1/2}}p^{-1/2})^2)=\frac T{2\pi^2}V+O(T^{1/2})$. For $S_X^{(2)}(t)=-\frac1{2\pi}\sum_{p\le X}p^{-1}\sin(t\log p^2)$ the same lemma gives $\int(S_X^{(2)})^2\ll T\sum_pp^{-2}+\#\{p\le X\}+(\sum_{p\le X}p^{-1})^2\ll T$. Finally $|S_X^{(\ge3)}|\le\frac1\pi\sum_p\sum_{m\ge3}\frac1{mp^{m/2}}\le\frac1\pi\sum_p\frac{p^{-3/2}}{1-p^{-1/2}}<\infty$, uniformly in $X$. Expanding $\mathcal W^2$ and bounding the cross terms by Cauchy--Schwarz gives \CGeqref{eq:D_second}.

For \CGeqref{eq:D_fourth} use $(y_1+y_2+y_3)^4\le27(y_1^4+y_2^4+y_3^4)$. Let $\mathcal D_1(t)=\sum_{p\in\mathcal J}p^{-1/2}p^{-it}$, so that $|\mathcal W_1|\le|\mathcal D_1|/\pi$ and $\mathcal D_1^2=\sum_n\eta_nn^{-it}$ with $\eta_{p^2}=p^{-1}$, $\eta_{p_1p_2}=2(p_1p_2)^{-1/2}$ for $p_1<p_2$, all $n\le T$. Theorem~\CGref{thm:known_MV} gives $\int|\mathcal D_1|^4=\sum|\eta_n|^2(T+O(n))\ll T\sum|\eta_n|^2\le T(1+2V^2)$. For $S_X^{(2)}$ let $\mathcal D_2(t)=\sum_{p\le X}p^{-1}(p^2)^{-it}$, so that $|S_X^{(2)}|\le|\mathcal D_2|/2\pi$ and $\mathcal D_2^2=\sum_n\eta'_nn^{-it}$ over $n=(p_1p_2)^2$ with $|\eta'_n|\le2(p_1p_2)^{-1}$; then $\sum|\eta'_n|^2\ll1$ and $\sum n|\eta'_n|^2\le4\#\{p\le X\}^2\le4X^2\le4T$. $S_X^{(\ge3)}$ is bounded.

For \CGeqref{eq:SX_derivative}, $S_X'(t)=\sum_{n\in A_X}a_n\log n\cos(t\log n)$ (termwise differentiation, Section~\CGref{subsec:euler_phase}). The terms $m=1$ have $|z_p|=\frac{\log p}{\pi\sqrt p}$; Lemma~\CGref{lem:trig_MV} and Theorem~\CGref{thm:known_mertens} give $\ll T\log X\sum_{p\le X}\frac{\log p}p+\sum_{p\le X}\log^2p+(\sum_{p\le X}\frac{\log p}{\sqrt p})^2\ll T(1+\log X)^2+X\log X$, and $X\le T^{1/2}$. The terms $m=2$ have $|z_{p^2}|=\frac{\log p}{\pi p}$ and contribute $\ll T+X\log X+\log^2X$. The terms $m\ge3$ are bounded by $\frac1\pi\sum_p\log p\sum_{m\ge3}p^{-m/2}<\infty$. Combine with $(y_1+y_2+y_3)^2\le3(y_1^2+y_2^2+y_3^2)$.
\end{proof}

\subsection{Proof of Theorem~\ref{thm:main}}
\label{subsec:main_proof}

\begin{proof}
\emph{Integrability.} On $[T,2T]$, $S$ is bounded (Theorem~\CGref{thm:known_basic}) and piecewise continuous, and $S_X$, $S_X'$, $\varpi$ are continuous; all integrals below exist.

\emph{Second and fourth moments of $E_X$.} Let $E_0=S-S^{\rm pr}_{Y_0}$ and $E_X=S-S_X=E_0+\mathcal W$. Lemma~\CGref{lem:selberg_uniform} with $k=2$ and H\"older's inequality give
\begin{equation}
\int_T^{2T}E_0^4\,dt\ll T,\qquad\int_T^{2T}E_0^2\,dt\le T^{1/2}\Bigl(\int_T^{2T}E_0^4\,dt\Bigr)^{1/2}\ll T .
\label{eq:E0_bounds}
\end{equation}
Let $\mathcal E_2=\int_T^{2T}E_X^2\,dt$. From \CGeqref{eq:D_second}, \CGeqref{eq:E0_bounds} and Cauchy--Schwarz,
\begin{equation}
\mathcal E_2=\frac T{2\pi^2}V+O\bigl(T(1+\sqrt V)\bigr),
\label{eq:J_asymptotic}
\end{equation}
and from \CGeqref{eq:D_fourth}, \CGeqref{eq:E0_bounds} and $(y_1+y_2)^4\le8(y_1^4+y_2^4)$,
\begin{equation}
\int_T^{2T}E_X^4\,dt\ll T(1+V)^2 .
\label{eq:R_fourth}
\end{equation}

\emph{Transfer to the ordinates.} Let $k_1,k_2$ be the smallest and largest indices in $\IT$ ($\IT\ne\emptyset$ for $T\ge T_0$) and apply \CGeqref{eq:identity_corrected}. By Lemma~\CGref{lem:s_bounds} and \CGeqref{eq:sigma_M_def}, $|u_k(X)|\ll\log T+\sqrt X\ll T^{1/4}$, so the boundary terms and $u_{k_1}(X)^2+u_{k_2}(X)^2$ are $O(T^{3/4})$. The intervals $[T,\gamma_{k_1}]$ and $[\gamma_{k_2},2T]$ have bounded length (Theorem~\CGref{thm:known_basic}), and on them $E_X^2(\lL+|S_X'|)\ll T^{1/2}\cdot T^{1/4}$. Since $|\varpi|\ll1/T$, replacing $(E_X+\varpi/\pi)^2$ by $E_X^2$ costs $\ll T^{-1}\int_T^{2T}(|E_X|+1)(\lL+|S_X'|)\,dt\ll(1+V)^{1/2}\log T$ by Cauchy--Schwarz, \CGeqref{eq:J_asymptotic} and \CGeqref{eq:SX_derivative}. Hence
\begin{equation}
\sum_{k\in\IT}u_k(X)^2=\int_T^{2T}E_X^2\lL\,dt+\int_T^{2T}E_X^2S_X'\,dt-\frac{\#\IT}{12}+O(T^{3/4}).
\label{eq:main_step1}
\end{equation}
\emph{Main term and derivative term.} Since $0\le\lL(t)-\lL(T)\le\frac{\log2}{2\pi}$ on $[T,2T]$, the first integral is $\lL(T)\mathcal E_2+O(\mathcal E_2)$. By Cauchy--Schwarz, \CGeqref{eq:R_fourth} and \CGeqref{eq:SX_derivative},
\begin{equation}
\Bigl|\int_T^{2T}E_X^2S_X'\,dt\Bigr|\le\Bigl(\int_T^{2T}E_X^4\Bigr)^{1/2}\Bigl(\int_T^{2T}S_X'^2\Bigr)^{1/2}\ll T(1+V)(1+\log X).
\label{eq:main_step2}
\end{equation}
Dividing by $\#\IT=T\lL(T)+O(T)$ and using \CGeqref{eq:J_asymptotic},
\begin{equation}
\frac1{\#\IT}\sum_{k\in\IT}u_k(X)^2=\frac V{2\pi^2}+O\Bigl(1+\sqrt V+\frac{1+V}{\log T}+(1+V)\,c\Bigr),
\label{eq:main_step3}
\end{equation}
because $\log X=c\log T$. By Lemma~\CGref{lem:V_size}, $(1+V)c\ll c(1+\log\frac1c)\ll1$ on $(0,\frac12]$; $V\ll\log\log T$ since $c\ge\log2/\log T$, so $(1+V)/\log T=o(1)$. A second use of Lemma~\CGref{lem:V_size} replaces $\frac V{2\pi^2}$ by $\frac1{2\pi^2}\log\frac1c+O(1)$ and $\sqrt V$ by $\sqrt{\log(1/c)}+O(1)$, which proves \CGeqref{eq:main_theorem}.
\end{proof}

\subsection{A cutoff varying with the index}
\label{subsec:variable_cutoff}
Let $\mathcal I_T^\tau=\{k:T\le\tau_k<2T\}$; by \CGeqref{eq:clock}, $\#\mathcal I_T^\tau=\NO(2T)-\NO(T)+O(1)\asymp T\log T$.

\begin{lemma}\label{lem:maximal_primes}
Let $T\ge T_0$ and $2\le X_-<X_+\le(T/3)^{1/2}$ with $X_+\le4X_-$; let $V_\pm=\sum_{X_-<p\le X_+}p^{-1}$ and
\[
\mathcal Q(t)=\max_{X_-\le y\le X_+}\Bigl|\sum_{X_-<p\le y}p^{-1/2-it}\Bigr| .
\]
Then $\mathcal Q$ is Lipschitz, $|\mathcal Q'(t)|\le\max_{X_-\le y\le X_+}|\sum_{X_-<p\le y}\log p\;p^{-1/2-it}|$ for almost every $t$, and
\begin{align}
&\int_{T-3}^{2T+3}\mathcal Q^4\,dt\ll TV_\pm^2,\qquad\int_{T-3}^{2T+3}|\mathcal Q'|^4\,dt\ll T(\log X_+)^4V_\pm^2,
\label{eq:maximal_L4}\\
&\langle\mathcal Q(\gamma_k)^2\rangle_{\IT}\ll V_\pm\Bigl(1+\frac{\log X_+}{\log T}\sqrt{\log\log T}\Bigr)+\frac{X_+}T,\qquad\langle\mathcal Q(\tau_k)^2\rangle_{\mathcal I^\tau_T}\ll V_\pm .
\label{eq:maximal_points}
\end{align}
\end{lemma}

\begin{proof}
If $(X_-,X_+]$ contains no prime, $\mathcal Q=0$. Otherwise let $p_1<\dots<p_J$ be its primes and $P_j(t)=\sum_{i\le j}p_i^{-1/2-it}$ ($P_0=0$), so that $\mathcal Q=\max_{0\le j\le J}|P_j|$, a maximum of finitely many Lipschitz functions; hence $\mathcal Q$ is Lipschitz and differentiable almost everywhere. At a point $t$ where $\mathcal Q$ and all $|P_j|$ are differentiable and $\mathcal Q(t)=|P_j(t)|$, the function $\mathcal Q-|P_j|\ge0$ has a minimum, so $\mathcal Q'(t)=|P_j|'(t)$, and $\bigl||P_j|'\bigr|\le|P_j'|$ wherever $|P_j|$ is differentiable. This gives the bound for $\mathcal Q'$.

\emph{Window.} Fix $\mathsf s\in C_c^\infty(-\frac12,\frac12)$, real, even, nonnegative and not identically $0$, let $\check{\mathsf s}(x)=\int\mathsf s(\omega)e^{i\omega x}d\omega$ and $\mathsf W=|\check{\mathsf s}|^2/\check{\mathsf s}(0)^2$. Then $\mathsf W\ge0$ is a Schwartz function with $\mathsf W(0)=1$, and $\mathcal FW$ is supported in $[-1,1]$. Choose $0<\mathsf y\le\frac14$ with $\mathsf W\ge\frac12$ on $[-\mathsf y,\mathsf y]$ and put $\mathfrak w(t)=\mathsf W(\mathsf y(t-\frac32T)/(\frac T2+3))$. Then $\mathfrak w\ge\frac12$ on $[T-3,2T+3]$, $\|\mathfrak w\|_4^4\asymp T$, and $\mathcal F\mathfrak w$ and $\mathcal F(\mathfrak w^4)$ are supported in $|\omega|\le2\mathsf y/T$ and $|\omega|\le8\mathsf y/T$.

\emph{Proof of \CGeqref{eq:maximal_L4}.} Let $f_j=\mathfrak wP_j$. Then $\mathcal Ff_j(\omega)=\sum_{i\le j}p_i^{-1/2}\mathcal F\mathfrak w(\omega+\log p_i)$, a sum of packets supported in $|\omega+\log p_i|\le2\mathsf y/T$. Since $\log p_{i+1}-\log p_i\ge\log(1+1/X_+)\ge(3/T)^{1/2}/2>4\mathsf y/T$ for $T\ge T_0$, the packets are disjoint, and $f_j$ is obtained from $f_J$ by removing the packets of $p_{j+1},\dots,p_J$, which lie below a point $\omega_j$ separating them from the others: $f_j(t)=f_J(t)-\frac1{2\pi}\int_{-\infty}^{\omega_j}\mathcal Ff_J(\omega)e^{i\omega t}d\omega$. Hence $\max_j|f_j|\le|f_J|+\mathcal Cf_J$, and on $[T-3,2T+3]$, $\mathcal Q\le2\max_j|f_j|$. By Theorem~\CGref{thm:known_carleson},
\[
\int_{T-3}^{2T+3}\mathcal Q^4\le16\int_\R\bigl(|f_J|+\mathcal Cf_J\bigr)^4\ll\int_\R\mathfrak w^4|P_J|^4\,dt .
\]
Write $P_J^2=\sum_n\eta_nn^{-it}$ with $\eta_{p^2}=p^{-1}$ and $\eta_{p_ip_{i'}}=2(p_ip_{i'})^{-1/2}$ for $i<i'$, all $n\le X_+^2\le T/3$ (unique factorization makes the representation unique). Expanding, $\int\mathfrak w^4|P_J^2|^2=\sum_{n,n'}\eta_n\eta_{n'}\,\mathcal F(\mathfrak w^4)(\log(n/n'))$, and $|\log(n/n')|\ge1/\max(n,n')\ge3/T>8\mathsf y/T$ for $n\ne n'$, so only the diagonal remains: $\int\mathfrak w^4|P_J|^4=\|\mathfrak w\|_4^4\sum_n\eta_n^2\ll T\bigl(\sum_pp^{-2}+4V_\pm^2\bigr)\ll TV_\pm^2$, because $\sum_{X_-<p\le X_+}p^{-2}\le V_\pm/X_-\le4V_\pm/X_+\le4V_\pm^2$. The same argument with the polynomials $\sum_{i\le j}\log p_i\,p_i^{-1/2-it}$, whose maximum dominates $|\mathcal Q'|$, gives the second bound, with $\eta_n$ multiplied by $\log p_i\log p_{i'}\le(\log X_+)^2$.

\emph{At the ordinates.} Let $T,2T$ be non-ordinates (otherwise move them by less than the distance to the nearest ordinate, which changes neither side). As Lebesgue--Stieltjes measures on $(T,2T]$, $dN=\lL\,dt+dG$, and $\mathcal Q^2$ is Lipschitz, so integration by parts \cite[\S3.5]{Folland1999} gives
\begin{equation}
\sum_{k\in\IT}\mathcal Q(\gamma_k)^2=\int_T^{2T}\mathcal Q^2\lL\,dt+\bigl[\mathcal Q^2G\bigr]_T^{2T}-2\int_T^{2T}G\,\mathcal Q\mathcal Q'\,dt,
\label{eq:maximal_stieltjes}
\end{equation}
multiplicities included. By Cauchy--Schwarz and \CGeqref{eq:maximal_L4} the first term is $\ll\log T\cdot T^{1/2}(TV_\pm^2)^{1/2}=TV_\pm\log T$. The boundary term is $\ll X_+\log T$, since $\mathcal Q\le\sum_{p\le X_+}p^{-1/2}\ll X_+^{1/2}$ and $G\ll\log T$. By H\"older's inequality with exponents $2,4,4$, Theorem~\CGref{thm:known_selberg_moments} ($\int_T^{2T}G^2\ll T\log\log T$) and \CGeqref{eq:maximal_L4}, the last term is $\ll(T\log\log T)^{1/2}(TV_\pm^2)^{1/4}(T(\log X_+)^4V_\pm^2)^{1/4}=T\sqrt{\log\log T}\,(\log X_+)V_\pm$. Divide by $\#\IT\gg T\log T$.

\emph{At the clock points.} For consecutive indices, $\tau_{k+1}-\tau_k=\int_k^{k+1}dx/\lL(\tau_x)\ge2\pi/\log(3T)=\mathfrak b_T$ while $\tau_{k+1}\le3T$, so the intervals $\mathcal U_k=[\tau_k,\tau_k+\mathfrak b_T]$, $k\in\mathcal I_T^\tau$, are disjoint and contained in $[T,2T+1]$. For an absolutely continuous $g$ on an interval $\mathcal U$ of length $\mathfrak b$ and $x\in\mathcal U$, $|g(x)|\le\frac1{\mathfrak b}\int_{\mathcal U}|g|+\int_{\mathcal U}|g'|$, hence $g(x)^2\le\frac2{\mathfrak b}\int_{\mathcal U}g^2+2\mathfrak b\int_{\mathcal U}g'^2$. With $g=\mathcal Q$ and summing,
$\sum_{k\in\mathcal I_T^\tau}\mathcal Q(\tau_k)^2\le\frac2{\mathfrak b_T}\int_T^{2T+1}\mathcal Q^2+2\mathfrak b_T\int_T^{2T+1}\mathcal Q'^2\ll TV_\pm\log T+\frac{T(\log X_+)^2V_\pm}{\log T}$, by Cauchy--Schwarz and \CGeqref{eq:maximal_L4}. Divide by $\#\mathcal I^\tau_T\gg T\log T$ and use $\log X_+\le\log T$.
\end{proof}

\begin{theorem}\label{thm:variable_cutoff}
For $\frac{\log2}{\log T}\le c\le\frac12$ let $X_k=\max\{2,(\tau_k/2\pi)^c\}$ and $\tilde u_k=s_k-S_{X_k}(\gamma_k)$. Uniformly in $c$,
\begin{equation}
\Bigl|\langle\tilde u_k^2\rangle_{\IT}-\frac1{2\pi^2}\log\frac1c\Bigr|\ll1+\sqrt{\log(1/c)} .
\label{eq:variable_cutoff}
\end{equation}
For fixed $c$, $\langle(\tilde u_k-u_k(X_-))^2\rangle_{\IT}\to0$, with $X_-$ as in the proof.
\end{theorem}

\begin{proof}
For $T\ge T_0$ and $k\in\IT$, $|\delta_k|\le1$ by \CGeqref{eq:BW_pointwise}, so $\tau_k\in[T-1,2T+1]$ and $X_-\le X_k\le X_+$ with $X_\pm$ the values of $\max\{2,(y/2\pi)^c\}$ at $y=T-1$ and $y=2T+1$. Then $X_+\le((2T+1)/2\pi)^{1/2}\le(T/3)^{1/2}$ and $X_+/X_-\le((2T+1)/(T-1))^{1/2}\le2$. The prime powers $p^m$ with $X_-<p\le X_k$ are those in $S_{X_k}-S_{X_-}$, so by \CGeqref{eq:SX_def}
\[
|S_{X_k}(\gamma_k)-S_{X_-}(\gamma_k)|\le\frac1\pi\mathcal Q(\gamma_k)+\frac1\pi\sum_{X_-<p\le X_+}\sum_{m\ge2}\frac1{mp^{m/2}}\le\frac1\pi\mathcal Q(\gamma_k)+1.1\,V_\pm ,
\]
since the terms $m=1$ are $\frac1\pi\Im\sum_{X_-<p\le X_k}p^{-1/2-i\gamma_k}$ and $\sum_{m\ge2}p^{-m/2}/m\le p^{-1}(1-2^{-1/2})^{-1}$. By Theorem~\CGref{thm:known_mertens}, $V_\pm=\log\frac{\log X_+}{\log X_-}+O(\frac1{\log X_-})\ll\frac1{\log X_-}\ll\frac1{1+c\log T}$ (and $V_\pm=0$ if $X_-=X_+=2$). Lemma~\CGref{lem:maximal_primes} gives
\[
\langle(\tilde u_k-u_k(X_-))^2\rangle_{\IT}\ll V_\pm+\frac{(1+c\log T)\sqrt{\log\log T}}{(1+c\log T)\log T}+T^{-1/2}+V_\pm^2\ll1,
\]
which tends to $0$ for fixed $c$. Theorem~\CGref{thm:main} with $X=X_-$, for which $\log(1/c_-)=\log(1/c)+O(1)$ with $c_-=\log X_-/\log T$, gives \CGeqref{eq:main_theorem} for $u_k(X_-)$; \CGeqref{eq:variable_cutoff} follows from $|\langle\tilde u^2\rangle-\langle u^2\rangle|\le2\langle u^2\rangle^{1/2}\langle(\tilde u-u)^2\rangle^{1/2}+\langle(\tilde u-u)^2\rangle$.
\end{proof}

%======================================================================
\section{Theorem C: exact accounting for the one-step clock}
\label{sec:accounting}

For a finite nonempty set $I$ of indices and a prime power $n$, define the response statistics at the zeros and at the clock points,
\begin{equation}
\Rg_n(I)=\frac{2}{a_n}\langle s_k\sin(\gamma_k\log n)\rangle_I,\qquad
\Rt_n(I)=\frac{2}{a_n}\langle s_k\sin(\tau_k\log n)\rangle_I,
\label{eq:two_responses}
\end{equation}
and, for a set $A$ of prime powers with $\sum_{n\in A}|a_n|<\infty$, the cross term
\begin{equation}
\Cr_A(I)=\langle S_A(\tau_k)^2\rangle_I-v_A .
\label{eq:Cr_def}
\end{equation}
For a prime $p$, $2/a_p=-2\pi\sqrt p$. These are finite-window statistics of the zeros of the Riemann zeta function, defined without any hypothesis; $\Cr_A$ contains the finite-window diagonal discrepancies and all off-diagonal products and has no fixed sign.

\begin{theorem}[Finite-window accounting]\label{thm:accounting}
Let $I$ be a finite nonempty set of indices and $A$ a set of prime powers with $\sum_{n\in A}|a_n|<\infty$.
\begin{enumerate}[label=(\alph*),leftmargin=2em]
\item For every $k$ and $X\ge2$,
\begin{equation}
\begin{split}
&r_k(X)=u_k(X)+\bigl(S_X(\gamma_k)-S_X(\tau_k)\bigr)+\chi_k,\\
&\chi_k=\frac{\gamma_k}{2\pi}\Bigl(\frac{\tau_k}{\gamma_k}-1-\log\frac{\tau_k}{\gamma_k}\Bigr),\qquad\frac{\delta_k^2}{4\pi\max(\gamma_k,\tau_k)}\le\chi_k\le\frac{\delta_k^2}{4\pi\min(\gamma_k,\tau_k)},\\
&|S_X(\gamma_k)-S_X(\tau_k)|\le M_X|\delta_k|.
\end{split}
\label{eq:one_step_decomposition}
\end{equation}
\item For every real $\lambda$,
\begin{equation}
\left\langle s_k^2-(s_k-\lambda S_A(\tau_k))^2\right\rangle_I
=\sum_{n\in A}a_n^2\Bigl(\lambda\,\Rt_n(I)-\frac{\lambda^2}{2}\Bigr)-\lambda^2\Cr_A(I).
\label{eq:exact_gain}
\end{equation}
In particular, for the full local factors,
\begin{equation}
\begin{split}
&\langle s_k^2-(s_k-S_X(\tau_k))^2\rangle_I=D_X(I)-\Cr_X(I),\\
&D_X(I)=\sum_{n\in A_X}a_n^2\Bigl(\Rt_n(I)-\frac12\Bigr)=2\langle s_kS_X(\tau_k)\rangle_I-v_X,
\end{split}
\label{eq:full_factor_gain}
\end{equation}
with $\Cr_X=\Cr_{A_X}$.
\item In units of zeros, the gain of the one-step clock over the bare clock is
\begin{equation}
\begin{split}
&\bigl\langle(\lL(\tau_k)\delta_k)^2-r_k(X)^2\bigr\rangle_I=D_X(I)-\Cr_X(I)+2\langle\chi_kS_X(\tau_k)\rangle_I,\\
&|\langle\chi_kS_X(\tau_k)\rangle_I|\le\sigma_X\max_{k\in I}\chi_k .
\end{split}
\label{eq:physical_gain}
\end{equation}
\end{enumerate}
\end{theorem}

\begin{proof}
(a) By \CGeqref{eq:one_step_clock}, $r_k(X)=\lL(\tau_k)\delta_k-S_X(\tau_k)$. Put $\chi_k=\lL(\tau_k)\delta_k-s_k$; since $s_k=\NO(\tau_k)-\NO(\gamma_k)=\int_{\gamma_k}^{\tau_k}\lL$,
\[
\chi_k=\int_{\gamma_k}^{\tau_k}\bigl(\lL(\tau_k)-\lL(y)\bigr)dy=\frac1{2\pi}\int_{\gamma_k}^{\tau_k}\log\frac{\tau_k}y\,dy=\frac{\gamma_k}{2\pi}\Bigl(\frac{\tau_k}{\gamma_k}-1-\log\frac{\tau_k}{\gamma_k}\Bigr),
\]
the last step by $\int\log y\,dy=y\log y-y$. If $\gamma_k<\tau_k$, then $\frac{\tau_k-y}{\tau_k}\le\log\frac{\tau_k}y\le\frac{\tau_k-y}{y}\le\frac{\tau_k-y}{\gamma_k}$ for $y\in[\gamma_k,\tau_k]$, and integration gives $\frac{\delta_k^2}{4\pi\tau_k}\le\chi_k\le\frac{\delta_k^2}{4\pi\gamma_k}$; if $\tau_k<\gamma_k$, write $\chi_k=\frac1{2\pi}\int_{\tau_k}^{\gamma_k}\log\frac y{\tau_k}\,dy$ and use $\frac{y-\tau_k}{\gamma_k}\le\frac{y-\tau_k}y\le\log\frac y{\tau_k}\le\frac{y-\tau_k}{\tau_k}$; if $\tau_k=\gamma_k$, $\chi_k=0$. Insert $s_k=u_k(X)+S_X(\gamma_k)$. The last bound is the mean value theorem with \CGeqref{eq:sigma_M_def}.

(b) The series $S_A(\tau_k)=\sum_na_n\sin(\tau_k\log n)$ converges absolutely, and so does the double series for $S_A(\tau_k)^2$; since $I$ is finite, the mean over $I$ may be exchanged with these sums. Expanding, the left side is $2\lambda\langle s_kS_A(\tau_k)\rangle_I-\lambda^2\langle S_A(\tau_k)^2\rangle_I$, with $\langle s_kS_A(\tau_k)\rangle_I=\sum_na_n\langle s_k\sin(\tau_k\log n)\rangle_I=\frac12\sum_na_n^2\Rt_n(I)$ and $\langle S_A(\tau_k)^2\rangle_I=v_A+\Cr_A(I)$. Taking $\lambda=1$ and $A=A_X$ gives \CGeqref{eq:full_factor_gain}.

(c) By (a), $\lL(\tau_k)\delta_k=s_k+\chi_k$ and $r_k(X)=s_k+\chi_k-S_X(\tau_k)$, so $(\lL(\tau_k)\delta_k)^2-r_k(X)^2=2s_kS_X(\tau_k)-S_X(\tau_k)^2+2\chi_kS_X(\tau_k)$; average and use \CGeqref{eq:full_factor_gain} and $|S_X|\le\sigma_X$.
\end{proof}

\begin{remark}\label{rem:accounting_meaning}
\CGeqref{eq:full_factor_gain} is an exact sign criterion for the gain of the full Euler phase at the clock points: the gain is positive exactly when the diagonal responses outweigh the cross term. No independence of the phases $\tau_k\log p$ is assumed. A threshold $\Rt_p=\frac12$ describes a single diagonal prime contribution and does not determine the optimal cutoff when $\Cr_X$ and the prime powers matter.
\end{remark}

\begin{remark}[Optimal trigonometric correction]\label{rem:matrix_projection}
For a finite set $A$ of frequencies let $\varsigma_n(k)=\sin(\tau_k\log n)$, $\mathbf Q=(\langle\varsigma_{n_1}\varsigma_{n_2}\rangle_I)_{n_1,n_2\in A}$, $\mathbf b=(\langle s_k\varsigma_n(k)\rangle_I)_{n\in A}$ and $\Upsilon_{\boldsymbol\lambda}=\sum_n\lambda_n\varsigma_n$. Ordinary least squares gives $\langle(s_k-\Upsilon_{\boldsymbol\lambda}(k))^2\rangle_I=\langle s_k^2\rangle_I-2\boldsymbol\lambda^{\mathsf T}\mathbf b+\boldsymbol\lambda^{\mathsf T}\mathbf Q\boldsymbol\lambda$, minimized by $\boldsymbol\lambda=\mathbf Q^\dagger\mathbf b$ with gain $\mathbf b^{\mathsf T}\mathbf Q^\dagger\mathbf b$; for prime powers $b_n=\frac12a_n\Rt_n(I)$. The Euler phase is the choice $\lambda_n=a_n$, and \CGeqref{eq:exact_gain} measures what it loses through $\Rt_n\ne1$ and $\Cr_A\ne0$.
\end{remark}

\begin{remark}[Clock point against zero]\label{rem:transport}
For a prime power $n$, $\sin(\gamma_k\ell_n)-\sin(\tau_k\ell_n)=\int_0^1\frac{d}{dy}\sin((\tau_k-y\delta_k)\ell_n)\,dy$, whence, exchanging the finite mean and the integral,
\begin{equation}
\Rg_n(I)-\Rt_n(I)=-\frac{2\ell_n}{a_n}\int_0^1\bigl\langle s_k\delta_k\cos\bigl((\tau_k-y\delta_k)\ell_n\bigr)\bigr\rangle_I\,dy .
\label{eq:response_transport}
\end{equation}
For fixed $n$ and $I=\IT$ the right side is $O(\langle|s_k\delta_k|\rangle_{\IT})=O(\log\log T/\log T)$ by \CGeqref{eq:s_delta} and Remark~\CGref{rem:korolev}; for $\log n=\mathsf q\log T$ with fixed $0<\mathsf q<1$ the phase change $\delta_k\ell_n$ has mean square $\mathsf q^2\log^2T\,\langle\delta_k^2\rangle_{\IT}\sim2\mathsf q^2\log\log T$, so it is not small, and the transport integral need not be small; no asymptotic statement about this difference is proved.
\end{remark}

\begin{remark}[Fixed prime sets]\label{rem:fixed_sup}
For fixed $X$, $|\delta_k^X-\delta_k|\le\sigma_X/\lL(\tau_k)\to0$ and $|u_k(X)-s_k|\le\sigma_X$. A fixed set of primes changes neither $\limsup_k|\delta_k|$ nor the unboundedness of the defect.
\end{remark}

%======================================================================
\section{Theorem D: a computable clock that approaches the zeros}
\label{sec:computable}

The one-step clock $\tau_k^X=\tau_k-y_k$, $y_k=S_X(\tau_k)/\lL(\tau_k)$, uses only $\tau_k$ and the primes $p\le X$. By Theorem~\CGref{thm:accounting}(a),
\begin{equation}
r_k(X)=u_k(X)+\mathfrak t_k(X)+\chi_k,\qquad\mathfrak t_k(X)=S_X(\gamma_k)-S_X(\tau_k),
\label{eq:r_decomposition}
\end{equation}
where Theorem~\CGref{thm:main} controls $u_k(X)$ and $0\le\chi_k\ll1/T$ on $\IT$. This section controls the transport term $\mathfrak t_k(X)$ in mean square, with no assumption on the joint behaviour of $\delta_k$ and the phases $\tau_k\log p$.

\begin{lemma}\label{lem:SX_fourth}
For $T\ge T_0$ and $2\le X\le T^{1/2}$,
\[
\int_{T-3}^{2T+3}S_X'(t)^4\,dt\ll T(1+\log X)^4,\qquad\int_{T-3}^{2T+3}S_X''(t)^4\,dt\ll T(1+\log X)^8 .
\]
\end{lemma}

\begin{proof}
By Section~\CGref{subsec:euler_phase}, the series
\[
S_X'(t)=\sum_{n\in A_X}a_n\log n\cos(t\log n),\qquad S_X''(t)=-\sum_{n\in A_X}a_n\log^2n\sin(t\log n)
\]
converge absolutely and uniformly. Split according to $m=1$, $m=2$, $m\ge3$ in $n=p^m$ and use $(y_1+y_2+y_3)^4\le27(y_1^4+y_2^4+y_3^4)$. For $\nu\in\{1,2\}$ the terms $m\ge3$ of the $\nu$-th derivative are bounded by $\frac1\pi\sum_p\sum_{m\ge3}\frac{(m\log p)^\nu}{mp^{m/2}}<\infty$, uniformly in $X$ and $t$. The terms $m=1$ are bounded by $\frac1\pi|\mathcal D_\nu(t)|$ with $\mathcal D_\nu(t)=\sum_{p\le X}\log^\nu p\;p^{-1/2-it}$; $\mathcal D_\nu^2=\sum_n\eta_nn^{-it}$ over $n=pp'\le X^2\le T$ with $\sum_n\eta_n^2\le2(\sum_{p\le X}\log^{2\nu}p/p)^2\ll(1+\log X)^{4\nu}$ (Theorem~\CGref{thm:known_mertens} and partial summation), so Theorem~\CGref{thm:known_MV} on an interval of length $T+6$ gives $\int|\mathcal D_\nu|^4=\sum\eta_n^2(T+6+O(n))\ll T(1+\log X)^{4\nu}$. The terms $m=2$ are bounded by $\frac1{2\pi}|\mathcal D'_\nu(t)|$ with $\mathcal D'_\nu(t)=\sum_{p\le X}(2\log p)^\nu p^{-1-2it}$; here $\mathcal D'^2_\nu=\sum_n\eta'_nn^{-it}$ over $n=(pp')^2\le X^4$, with $\sum_n\eta'^2_n\ll1$ and $\sum_nn\eta'^2_n\ll(\sum_{p\le X}\log^{2\nu}p)^2\ll X^2(1+\log X)^{4\nu}\le T(1+\log X)^{4\nu}$, so again $\int|\mathcal D'_\nu|^4\ll T(1+\log X)^{4\nu}$.
\end{proof}

\begin{lemma}[Transport]\label{lem:transport}
For $T\ge T_0$ and $2\le X\le T^{1/2}$,
\begin{equation}
\bigl\langle\mathfrak t_k(X)^2\bigr\rangle_{\IT}\ll\frac{\log\log T\,(1+\log X)^2}{\log^2T}.
\label{eq:transport}
\end{equation}
\end{lemma}

\begin{proof}
\emph{Fourth moment of $\delta_k$.} For $T\ge T_0$ and $k\in\IT$, $|\delta_k|\le1$ by \CGeqref{eq:BW_pointwise}; so the point $\xi_k$ of \CGeqref{eq:s_delta} is $\ge T-1$ and $\lL(\xi_k)\gg\log T$. By \CGeqref{eq:s_moments} with $j=2$,
\begin{equation}
\langle\delta_k^4\rangle_{\IT}=\langle s_k^4\lL(\xi_k)^{-4}\rangle_{\IT}\ll\frac{(\log\log T)^2}{\log^4T}.
\label{eq:delta_fourth}
\end{equation}
\emph{Maximal function.} Let $\varphi\in C^\infty(\R)$ with $0\le\varphi\le1$, $\varphi=1$ on $[T-2,2T+2]$, $\varphi=0$ outside $[T-3,2T+3]$ and $|\varphi'|\le2$, and let $\mathfrak h=\varphi S_X'$. For $k\in\IT$ the segment between $\gamma_k$ and $\tau_k$ lies in $[T-1,2T+1]$, where $\mathfrak h=S_X'$, so
\begin{equation}
|\mathfrak t_k(X)|=\Bigl|\int_{\tau_k}^{\gamma_k}\mathfrak h(y)\,dy\Bigr|\le\int_{\tau_k-|\delta_k|}^{\tau_k+|\delta_k|}|\mathfrak h|\le2|\delta_k|\,\mathcal M\mathfrak h(\tau_k)
\label{eq:transport_maximal}
\end{equation}
(trivially if $\delta_k=0$). The function $\mathfrak h$ is $C^\infty$ with compact support. By Theorem~\CGref{thm:known_maximal}, $g=\mathcal M\mathfrak h$ is locally absolutely continuous with $\|g\|_4\ll\|\mathfrak h\|_4$ and $\|g'\|_4\le\|\mathcal M(\mathfrak h')\|_4\ll\|\mathfrak h'\|_4$; by Lemma~\CGref{lem:SX_fourth}, $\|\mathfrak h\|_4^4\ll T(1+\log X)^4$ and, since $\mathfrak h'=\varphi'S_X'+\varphi S_X''$, $\|\mathfrak h'\|_4^4\ll T(1+\log X)^8$.

\emph{Sampling at the clock points.} As in the proof of Lemma~\CGref{lem:maximal_primes}, the intervals $\mathcal U_k=[\tau_k,\tau_k+\mathfrak b_T]$, $k\in\IT$, are disjoint and contained in $[T-1,2T+2]$. For $x\in\mathcal U$, an interval of length $\mathfrak b$, $|g(x)|\le\frac1{\mathfrak b}\int_{\mathcal U}|g|+\int_{\mathcal U}|g'|$; by H\"older's inequality and $(y_1+y_2)^4\le8(y_1^4+y_2^4)$, $g(x)^4\le\frac8{\mathfrak b}\int_{\mathcal U}g^4+8\mathfrak b^3\int_{\mathcal U}|g'|^4$. Summing over $k\in\IT$ and dividing by $\#\IT\gg T\log T$,
\[
\langle g(\tau_k)^4\rangle_{\IT}\ll\frac1T\int_\R g^4+\frac1{T\log^4T}\int_\R|g'|^4\ll(1+\log X)^4+\frac{(1+\log X)^8}{\log^4T}\ll(1+\log X)^4 .
\]
\emph{Conclusion.} By \CGeqref{eq:transport_maximal} and Cauchy--Schwarz, $\langle\mathfrak t_k^2\rangle\le4\langle\delta_k^2g(\tau_k)^2\rangle\le4\langle\delta_k^4\rangle^{1/2}\langle g(\tau_k)^4\rangle^{1/2}$; insert \CGeqref{eq:delta_fourth}.
\end{proof}

\begin{theorem}[Theorem D]\label{thm:computable_clock}
There are absolute constants $T_0$ and $C$ such that for $T\ge T_0$ and $2\le X\le T^{1/2}$, with $c=\log X/\log T$,
\begin{equation}
\Bigl|\bigl\langle r_k(X)^2\bigr\rangle_{\IT}-\frac1{2\pi^2}\log\frac1c\Bigr|\le C\Bigl(1+\sqrt{\log(1/c)}\Bigr)\bigl(1+c^2\log\log T\bigr).
\label{eq:computable_clock}
\end{equation}
\end{theorem}

\begin{proof}
By Theorem~\CGref{thm:accounting}(a), $0\le\chi_k\le\delta_k^2/(4\pi(T-1))\le1/T$ on $\IT$. Since $1+\log X\le(1+\frac1{\log2})c\log T$, Lemma~\CGref{lem:transport} gives $\langle\mathfrak t_k^2\rangle\ll c^2\log\log T$, so $d_k=r_k(X)-u_k(X)=\mathfrak t_k+\chi_k$ has $\langle d_k^2\rangle\ll c^2\log\log T+T^{-2}$. Theorem~\CGref{thm:main} gives $\langle u_k(X)^2\rangle\ll1+\log(1/c)$. Then $|\langle r_k^2\rangle-\langle u_k^2\rangle|\le2\langle u_k^2\rangle^{1/2}\langle d_k^2\rangle^{1/2}+\langle d_k^2\rangle\ll(1+\sqrt{\log(1/c)})(c\sqrt{\log\log T}+T^{-1})+c^2\log\log T$, which is bounded by the right side of \CGeqref{eq:computable_clock} because $y\le1+y^2$. Combine with \CGeqref{eq:main_theorem}.
\end{proof}

\begin{corollary}\label{cor:computable_clock}
Let $X=T^{1/\sqrt{\log\log T}}$. Then
\begin{align}
\langle r_k(X)^2\rangle_{\IT}&=\frac1{4\pi^2}\log\log\log T+O\bigl(\sqrt{\log\log\log T}\bigr),
\label{eq:computable_units}\\
\frac1{\#\IT}\sum_{k\in\IT}\bigl(\tau_k^X-\gamma_k\bigr)^2&=\frac{\log\log\log T+O(\sqrt{\log\log\log T})}{\log^2T},
\label{eq:computable_height}
\end{align}
while $\frac1{\#\IT}\sum_{k\in\IT}(\tau_k-\gamma_k)^2=(2+o(1))\log\log T/\log^2T$; the ratio is asymptotic to $\log\log\log T/(2\log\log T)$.
\end{corollary}

\begin{proof}
Here $c=(\log\log T)^{-1/2}$, $2\le X\le T^{1/2}$ for $T\ge T_0$, $\log(1/c)=\frac12\log\log\log T$ and $c^2\log\log T=1$, so \CGeqref{eq:computable_clock} gives \CGeqref{eq:computable_units}. On $\IT$, $\tau_k\in[T-1,2T+1]$ and $\lL(\tau_k)=\frac{\log T}{2\pi}(1+O(\frac1{\log T}))$; divide $r_k^2$ by $\lL(\tau_k)^2$. For the bare clock, $\delta_k=s_k/\lL(\xi_k)$ with $\lL(\xi_k)=\frac{\log T}{2\pi}(1+O(\frac1{\log T}))$, and $\langle s_k^2\rangle\sim\frac1{2\pi^2}\log\log T$ by \CGeqref{eq:s_moments}.
\end{proof}

\begin{corollary}[A cutoff read off the clock]\label{cor:local_clock}
For $\tau_k>e^e$ let $X_k^{\rm loc}=\tau_k^{1/\sqrt{\log\log\tau_k}}$ and $\tau_k^{\rm loc}=\tau_k-S_{X_k^{\rm loc}}(\tau_k)/\lL(\tau_k)$. Then
\[
\frac1{\#\mathcal I_T^\tau}\sum_{k\in\mathcal I_T^\tau}\bigl(\tau_k^{\rm loc}-\gamma_k\bigr)^2=\frac{\log\log\log T+O(\sqrt{\log\log\log T})}{\log^2T}.
\]
\end{corollary}

\begin{proof}
Let $X_0=T^{1/\sqrt{\log\log T}}$ and $X_0'=(2T)^{1/\sqrt{\log\log2T}}$. The function $y\mapsto y/\sqrt{\log y}$ is increasing for $y>e^{1/2}$, so $X_0\le X_k^{\rm loc}\le X_0'$ for $k\in\mathcal I_T^\tau$, and $\log X_0'-\log X_0\le\log2/\sqrt{\log\log T}=o(1)$; hence $X_0'\le2X_0$ and $X_0'\le(T/3)^{1/2}$ for $T\ge T_0$. As in the proof of Theorem~\CGref{thm:variable_cutoff}, $|S_{X_k^{\rm loc}}(\tau_k)-S_{X_0}(\tau_k)|\le\frac1\pi\mathcal Q(\tau_k)+1.1V_\pm$ with $X_\pm$ replaced by $X_0,X_0'$, and Lemma~\CGref{lem:maximal_primes} gives a mean square over $\mathcal I_T^\tau$ that is $\ll V_\pm\ll1/\log X_0=\sqrt{\log\log T}/\log T$. The sets $\IT$ and $\mathcal I_T^\tau$ differ by $O(\log T)$ indices, those with $\gamma_k$ or $\tau_k$ within distance $1$ of $T$ or $2T$, and for them $|r_k(X_0)|\le|s_k|+\chi_k+\sigma_{X_0}\ll\log T+X_0^{1/2}$ (Lemma~\CGref{lem:s_bounds}); so the means of $r_k(X_0)^2$ over the two sets differ by $O(T^{-1/2})$, and \CGeqref{eq:computable_units} holds over $\mathcal I_T^\tau$. Finally $\lL(\tau_k)(\tau_k^{\rm loc}-\gamma_k)=r_k(X_0)-(S_{X_k^{\rm loc}}(\tau_k)-S_{X_0}(\tau_k))$; conclude with Cauchy--Schwarz and $\lL(\tau_k)=\frac{\log T}{2\pi}(1+O(\frac1{\log T}))$.
\end{proof}

\begin{proposition}[The root of the corrected equation in the monotone regime]\label{prop:monotone}
Let $X=X(T)\ge2$ with $M_X\le\frac12\lL(T)$, and let $\mathcal I'_T$ be the set of $k\in\IT$ with $F_X(T)\le k-\frac12\le F_X(2T)$, so that the root $t_k^X\in[T,2T]$ of Remark~\CGref{rem:localization} exists. Then $X\ll\log^2T$, $\#(\IT\setminus\mathcal I'_T)\ll\log T$, and
\[
\frac1{\#\IT}\sum_{k\in\mathcal I'_T}\bigl(t_k^X-\gamma_k\bigr)^2\ge\frac4{9\lL(2T)^2}\Bigl(\frac1{2\pi^2}\log\log T-O\bigl(\sqrt{\log\log T}\bigr)\Bigr)\gg\frac{\log\log T}{\log^2T}.
\]
\end{proposition}

\begin{proof}
$M_X\ge\frac1\pi\sum_{p\le X}\log p/\sqrt p\gg\sqrt X$ by Chebyshev's estimate, so $X\ll\lL(T)^2\ll\log^2T$ and $c=\log X/\log T\le(2\log\log T+O(1))/\log T$. On $[T,2T]$, $0<F_X'\le\lL+M_X\le\frac32\lL(2T)$, so $|u_k(X)|=|F_X(t_k^X)-F_X(\gamma_k)|\le\frac32\lL(2T)|t_k^X-\gamma_k|$ for $k\in\mathcal I'_T$. An index $k\in\IT$ is missing from $\mathcal I'_T$ only if $k-\frac12<F_X(T)$ or $k-\frac12>F_X(2T)$; since $N(T)<k\le N(2T)$ and $|F_X-N|\le|G|+\sigma_X\ll\log T$ at $T$ and $2T$, there are $O(\log T)$ such $k$, each with $|u_k(X)|\ll\log T$ (Lemma~\CGref{lem:s_bounds}). Hence $\frac1{\#\IT}\sum_{k\in\mathcal I'_T}u_k(X)^2=\langle u_k(X)^2\rangle_{\IT}+O(\frac{\log^3T}{T\log T})$, and Theorem~\CGref{thm:main} with $\log(1/c)\ge\log\log T-\log(2\log\log T+O(1))$ bounds it from below.
\end{proof}

Thus, in the regime $M_X\le\frac12\lL(T)$ where the corrected equation is monotone on the whole window, its root gains nothing in order: the error in height of $t_k^X$ has mean square of the same order $\log\log T/\log^2T$ as that of $\tau_k$. The proposition says nothing about roots at larger cutoffs selected by a local criterion. Theorem~\CGref{thm:computable_clock} avoids roots and takes one Newton-type step instead.

\begin{remark}[The Newton term and the restriction on $c$]\label{rem:newton}
Let $\xi^X_k$ be a point between $\gamma_k$ and $\tau_k$ with $\mathfrak t_k(X)=-S_X'(\xi_k^X)\delta_k$ (mean value theorem). Since $\delta_k=(r_k(X)+S_X(\tau_k))/\lL(\tau_k)$ by \CGeqref{eq:one_step_clock}, \CGeqref{eq:r_decomposition} becomes the exact relation
\begin{equation}
r_k(X)\Bigl(1+\frac{S_X'(\xi_k^X)}{\lL(\tau_k)}\Bigr)=u_k(X)+\chi_k-y_k\,S_X'(\xi_k^X).
\label{eq:newton_exact}
\end{equation}
The term $y_kS_X'(\xi_k^X)$ is the second-order term of Newton's method. For fixed $c$, Theorem~\CGref{thm:main} gives $\|u(X)\|_{2,\IT}=O_c(1)$, and $\|\chi\|_{2,\IT}\ll1/T$; by \CGeqref{eq:r_decomposition} and the triangle inequality, $\|r(X)\|_{2,\IT}=O_c(1)$ holds if and only if $\|\mathfrak t(X)\|_{2,\IT}=O_c(1)$. A transport of unbounded mean square cannot be compensated by the defect, which is bounded in mean square. Lemma~\CGref{lem:transport} gives only the upper bound $\langle\mathfrak t_k(X)^2\rangle_{\IT}\ll c^2\log\log T$, and whether the transport is bounded for fixed $c$ is open (Question~\CGref{q:fixed_c}). A heuristic count, treating $S_X(\tau_k)$ and $S_X'(\xi^X_k)$ as uncorrelated with $\langle S_X^2\rangle\approx\frac1{2\pi^2}\log\log X$ and $\langle S_X'^2\rangle/\lL(T)^2\approx c^2$, would give the Newton term a mean square of order $c^2\log\log T$; this is not proved, and it is not a lower bound for the transport. The denominator in \CGeqref{eq:newton_exact} can vanish only when $S_X'(\xi_k^X)=-\lL(\tau_k)$; its sign cannot be decided from $F_X'(\tau_k)$ alone, because the derivative here is evaluated at the intermediate point $\xi_k^X$. The construction \CGeqref{eq:one_step_clock} does not divide by this denominator or select a root of $F_X$.
\end{remark}

%======================================================================
\section{Iterating the Newton-type step}
\label{sec:iteration}

Theorem~\CGref{thm:computable_clock} takes one Newton-type step from $\tau_k$. Its restriction $c\le(\log\log T)^{-1/2}$ comes from the available bound (Lemma~\CGref{lem:transport}) for the transport over the interval between $\tau_k$ and $\gamma_k$, whose length is that of $\delta_k$. After one step the error is smaller, so a second step can use a larger cutoff. This section proves the second step and then every further step, with constants uniform in the number of steps (Section~\CGref{subsec:further_steps}).

\subsection{The exact recursion}
\label{subsec:recursion}
For a sequence of cutoffs $X_1,X_2,\dots\ge2$ define $x^{(0)}_k=\tau_k$ and
\begin{equation}
x^{(j)}_k=x^{(j-1)}_k-\frac{\NO(x^{(j-1)}_k)+S_{X_j}(x^{(j-1)}_k)-(k-\frac12)}{\lL(x^{(j-1)}_k)},\qquad e_{j,k}=x^{(j)}_k-\gamma_k,
\label{eq:newton_iteration}
\end{equation}
so that $x^{(1)}_k=\tau_k^{X_1}$ and $e_{0,k}=\delta_k$. For a finite index set $I$ and $r\ge1$ let $\|a\|_{r,I}=\langle|a_k|^r\rangle_I^{1/r}$.

\begin{proposition}[Exact recursion]\label{prop:recursion}
Let $j\ge1$ and suppose $x^{(j-1)}_k>2\pi$. Then
\begin{equation}
\lL(x^{(j-1)}_k)\,e_{j,k}=u_k(X_j)+\bigl(S_{X_j}(\gamma_k)-S_{X_j}(x^{(j-1)}_k)\bigr)+\chi_{j,k},\qquad 0\le\chi_{j,k}\le\frac{e_{j-1,k}^2}{4\pi\min(\gamma_k,x^{(j-1)}_k)} .
\label{eq:recursion_identity}
\end{equation}
Let $\varphi_j\in C^\infty(\R)$, $0\le\varphi_j\le1$, be $1$ on an interval containing all segments between $\gamma_k$ and $x^{(j-1)}_k$, $k\in I$, and $\mathsf g_j=\mathcal M(\varphi_jS'_{X_j})$. Then for every $r\ge1$,
\begin{equation}
\begin{split}
&\bigl\|\lL(x^{(j-1)})e_j\bigr\|_{2r,I}\le\|u(X_j)\|_{2r,I}+2\,\|e_{j-1}\|_{4r,I}\,\|\mathsf g_j(x^{(j-1)})\|_{4r,I}+\|\chi_j\|_{2r,I},\\
&\|\chi_j\|_{2r,I}\le\frac{\|e_{j-1}\|_{4r,I}^2}{4\pi\min_{k\in I}\min(\gamma_k,x^{(j-1)}_k)} .
\end{split}
\label{eq:recursion_norms}
\end{equation}
\end{proposition}

\begin{proof}
Write $x=x^{(j-1)}_k$. By \CGeqref{eq:newton_iteration}, $\lL(x)e_{j,k}=\lL(x)(x-\gamma_k)-\NO(x)-S_{X_j}(x)+k-\frac12$. As in the proof of Theorem~\CGref{thm:accounting}(a), $\NO(x)-\NO(\gamma_k)=\int_{\gamma_k}^x\lL=\lL(x)(x-\gamma_k)-\chi_{j,k}$ with $\chi_{j,k}=\int_{\gamma_k}^{x}(\lL(x)-\lL(y))\,dy$, which lies between $(x-\gamma_k)^2/(4\pi\max(\gamma_k,x))$ and $(x-\gamma_k)^2/(4\pi\min(\gamma_k,x))$ by the same estimates of $\log$. Hence $\lL(x)e_{j,k}=k-\frac12-\NO(\gamma_k)-S_{X_j}(x)+\chi_{j,k}=u_k(X_j)+S_{X_j}(\gamma_k)-S_{X_j}(x)+\chi_{j,k}$. For \CGeqref{eq:recursion_norms}, $|S_{X_j}(\gamma_k)-S_{X_j}(x)|\le\int_{x-|e_{j-1,k}|}^{x+|e_{j-1,k}|}|\varphi_jS'_{X_j}|\le2|e_{j-1,k}|\,\mathsf g_j(x)$, and Minkowski's and H\"older's inequalities in $\ell^{2r}(I)$.
\end{proof}

The recursion \CGeqref{eq:recursion_norms} is exact and uses no independence between the errors and the prime phases. It reduces the error of each step to three quantities: the moments of the defect $u_k(X_j)$ at the true ordinates (Theorem~\CGref{thm:main} and its higher-moment versions), the moments of the previous error, and the moments of the maximal function $\mathsf g_j$ sampled at the computable points $x^{(j-1)}_k$. The last is the only new difficulty: the points $x^{(j-1)}_k$ are not separated like the $\tau_k$.

\subsection{Moments of order \texorpdfstring{$2r$}{2r} of the defect}
\begin{lemma}\label{lem:higher_moments}
There is an absolute constant $C$ such that for every integer $r\ge1$, every $T\ge T_0$ with $2r\le\kappa_{\rm F}\log T$ and every $2\le X\le T^{1/(160r)}$, with $V$ as in Lemma~\CGref{lem:V_size} computed with $Y_{2r}=T^{1/(160r)}$ in place of $Y_0$,
\[
\int_T^{2T}E_X^{4r}\,dt\le\bigl(C(r^2+\sqrt{r(1+V)})\bigr)^{4r}T,\qquad
\langle u_k(X)^{2r}\rangle_{\IT}\le\bigl(C(r^2+\sqrt{r(1+V)})\bigr)^{2r},
\]
and for every $r\ge1$ and $2\le X\le T^{1/(2r)}$,
\[
\int_{T-3}^{2T+3}|S_X^{(\nu)}(t)|^{4r}dt\le(Cr)^{2r}\,T(1+\log X)^{4r\nu}\qquad(\nu=1,2).
\]
\end{lemma}

\begin{proof}
The first bound: $E_X=E_0+\mathcal W$ with $E_0=S-S^{\rm pr}_{Y_{2r}}$, $\|E_0\|_{L^{4r}(T,2T)}\le C(2r)^2T^{1/4r}$ by Lemma~\CGref{lem:selberg_uniform} with $k=2r$ (admissible since $2r\le\kappa_{\rm F}\log T$), and $\mathcal W$ as in Lemma~\CGref{lem:D_moments}: $|\mathcal W_1|\le|\mathcal D_1|/\pi$, $\mathcal D_1^{2r}=\sum_n\eta_nn^{-it}$ over products of $2r$ primes $\le\max(X,Y_{2r})$, so $n\le T^{1/80}$, with $\sum|\eta_n|^2\le(2r)!\,(V+1)^{2r}$ (each $n$ has at most $(2r)!$ ordered factorizations, so $\sum_n|\eta_n|^2\le(2r)!\sum_{p_1,\dots,p_{2r}}\prod p_i^{-1}$); Theorem~\CGref{thm:known_MV} gives $\int|\mathcal D_1|^{4r}\le2T\sum|\eta_n|^2$, whence $\|\mathcal W_1\|_{L^{4r}}\le C\sqrt r\sqrt{1+V}\,T^{1/4r}$ by $((2r)!)^{1/4r}\le\sqrt{2r}$. The terms $m=2$ and $m\ge3$ are $O(T^{1/4r})$ in the same norm (products of $2r$ squares of primes, $\sum|\eta'_n|^2\le(2r)!\,(\sum_pp^{-2})^{2r}\le(2r)!$, and a bounded function). Minkowski's inequality in $L^{4r}$ gives the first bound. The second: Theorem~\CGref{thm:identity_general} with $j=r$ and $F=F_X$, $\Psi_r(u)\ge u^{2r}$, gives $\sum_{k\in\IT}u_k^{2r}\le\int_T^{2T}(|E_X|+T^{-1})^{2r}(\lL+|S_X'|)\,dt+C_{\rm b}\bigl(4^rT^{1/100}+(2\log T)^{2r+1}\bigr)$ with an absolute $C_{\rm b}$, the last term bounding the boundary terms and the end pieces, since $|u_k(X)|\le\log T+\sqrt X$ (Lemma~\CGref{lem:s_bounds} and \CGeqref{eq:sigma_M_def}) and $X^{r+1/2}\le T^{1/100}$. Since $r\le\log T/(160\log2)$, $4^rT^{1/100}\le T^{1/80+1/100}\le T^{1/2}$. For $(2\log T)^{2r+1}$ there are two cases: if $r\ge\sqrt{2\log T/C}$ then $2\log T\le Cr^2$ and $(2\log T)^{2r+1}\le(Cr^2)^{2r}\,2\log T\le(Cr^2)^{2r}\,T\log T$; if $r<\sqrt{2\log T/C}$ then $(2\log T)^{2r+1}\le\exp\bigl(3\sqrt{2\log T/C}\,\log(2\log T)\bigr)\le T^{1/2}$ for $T\ge T_0$. In both cases the boundary term is at most $C_{\rm b}(Cr^2)^{2r}(T^{1/2}+T\log T)\le9\pi C_{\rm b}(Cr^2)^{2r}\,\#\IT$ for $T\ge T_0$. With $\mathcal E_{4r}=\int_T^{2T}E_X^{4r}$, H\"older's inequality gives $\int_T^{2T}E_X^{2r}\le T^{1/2}\mathcal E_{4r}^{1/2}$ and, by \CGeqref{eq:SX_derivative}, $\int_T^{2T}E_X^{2r}|S_X'|\le\mathcal E_{4r}^{1/2}C_{\rm d}T^{1/2}(1+\log X)$, $C_{\rm d}$ the constant of \CGeqref{eq:SX_derivative}; so the integral is at most $(\lL(2T)+C_{\rm d}(1+\log X))T^{1/2}(\mathcal E_{4r}^{1/2}+T^{-2r+1/2})$. Dividing by $\#\IT\ge\frac12T\lL(T)$ and using $\lL(2T)\le2\lL(T)$ and $1+\log X\le2\pi\lL(T)$, $\langle u_k(X)^{2r}\rangle_{\IT}\le(4+4\pi C_{\rm d})T^{-1/2}\mathcal E_{4r}^{1/2}+(9\pi C_{\rm b}+1)(Cr^2)^{2r}$, and the first bound gives the second with $(9\pi C_{\rm b}+2)C$ in place of $C$, since $(9\pi C_{\rm b}+2)^{2r}\ge9\pi C_{\rm b}+2$. The third: Lemma~\CGref{lem:SX_fourth} with products of $2r$ primes, $n\le X^{2r}\le T$, and $\sum\eta_n^2\le(2r)!\,(\sum_{p\le X}\log^{2\nu}p/p)^{2r}\le(Cr)^{2r}(1+\log X)^{4r\nu}$, by Theorem~\CGref{thm:known_mertens} and partial summation; the terms $m=2$ and $m\ge3$ as in Lemma~\CGref{lem:SX_fourth}.
\end{proof}

The dependence on $r$ matters in Theorem~\CGref{thm:uniform_iteration}, where $r$ grows with the number of Newton-type steps; for fixed $r$ the three bounds read $\ll_rT(1+V)^{2r}$, $\ll_r(1+V)^r$ and $\ll_rT(1+\log X)^{4r\nu}$.

The moments of all orders give a tail bound for the defect that no second-moment argument can give: the proportion of zeros whose corrected index is off by more than $\lambda$ positions decays exponentially in $\sqrt\lambda$, provided the cutoff is chosen in a range that depends on $\lambda$.

\begin{theorem}[Exponential tails of the ordinal defect]\label{thm:tails}
There are absolute constants $C_8\ge1$, $\lambda_0$ and $T_0$ such that, for $T\ge T_0$, every $\lambda\ge\lambda_0$ and every $X=T^c\ge2$ with
\begin{equation}
\exp\bigl(-(\lambda/C_8)^{3/2}\bigr)\le c\le\frac1{C_8\sqrt\lambda},
\label{eq:tails_range}
\end{equation}
\begin{equation}
\#\{k\in\IT:|u_k(X)|>\lambda\}\le2^{\,2-\sqrt{\lambda/C_8}}\;\#\IT .
\label{eq:tails}
\end{equation}
\end{theorem}

\begin{proof}
Let $C$ be the constant of Lemma~\CGref{lem:higher_moments} and $r=\lfloor\sqrt{\lambda/(4C)}\rfloor$, so $r\ge1$ for $\lambda\ge4C$ and $Cr^2\le\lambda/4$. The lemma applies (we may take $C\ge1$). First, $X\le T^{1/(160r)}$ because $c\le1/(C_8\sqrt\lambda)\le1/(160r)$ once $C_8\ge80$, since $160r\le160\sqrt{\lambda/(4C)}\le80\sqrt\lambda$. Second, $2r\le\kappa_{\rm F}\log T$: indeed $2r\le\sqrt{\lambda/C}\le\sqrt\lambda$, while $X=T^c\ge2$ gives $\log T\ge\log2/c\ge C_8\log2\,\sqrt\lambda\ge\sqrt\lambda/\kappa_{\rm F}$ once $C_8\ge300=1/(\kappa_{\rm F}\log2)$. With $Y_{2r}=T^{1/(160r)}\ge X$, Lemma~\CGref{lem:V_size} gives $V=\log(1/(160rc))+O(1)\le\log(1/c)+C_9\le(\lambda/C_8)^{3/2}+C_9$, with an absolute constant $C_9\ge0$, by \CGeqref{eq:tails_range}, so $1+V\le2(\lambda/C_8)^{3/2}$ for $\lambda\ge C_8(C_9+1)^{2/3}$. Then $r(1+V)\le\sqrt{\lambda/(4C)}\cdot2\lambda^{3/2}C_8^{-3/2}=\lambda^2/(\sqrt C\,C_8^{3/2})$ and $C\sqrt{r(1+V)}\le C^{3/4}C_8^{-3/4}\lambda\le\lambda/4$ once $C_8\ge7C$. Hence $C(r^2+\sqrt{r(1+V)})\le\lambda/2$, and Chebyshev's inequality with the second bound of the lemma gives
\[
\#\{k\in\IT:|u_k(X)|>\lambda\}\le\lambda^{-2r}\sum_{k\in\IT}u_k(X)^{2r}\le2^{-2r}\,\#\IT\le2^{\,2-\sqrt{\lambda/C}}\,\#\IT ,
\]
since $2r\ge2\sqrt{\lambda/(4C)}-2$. Take $C_8=\max(300,7C)$ and $\lambda_0=\max\bigl(4C,C_8(C_9+1)^{2/3}\bigr)$; then every condition used above holds for $\lambda\ge\lambda_0$, and $2^{-\sqrt{\lambda/C}}\le2^{-\sqrt{\lambda/C_8}}$.
\end{proof}

For the cutoff $X=T^{1/(C_8\sqrt\lambda)}$, Theorem~\CGref{thm:main} gives $\langle u_k(X)^2\rangle_{\IT}=\frac1{4\pi^2}\log\lambda+O(\sqrt{\log\lambda})$, so Chebyshev's inequality alone bounds the proportion in \CGeqref{eq:tails} by $\lambda^{-2}\log\lambda$; the moments of order $2r\asymp\sqrt\lambda$ replace this by $2^{2-\sqrt{\lambda/C_8}}$. The bound \CGeqref{eq:tails} is for cutoffs that depend on $\lambda$; for a fixed cutoff exponent $c$ it applies only to $\lambda\le(C_8c)^{-2}$, and no exponential tail is claimed beyond that. The range \CGeqref{eq:tails_range} is the price of the moment order: Lemma~\CGref{lem:higher_moments} needs $X\le T^{1/(160r)}$, and Lemma~\CGref{lem:selberg_uniform} with $2r$ in place of $k$ needs the prime polynomial of length $T^{1/(160r)}$ and the order $2r\le\kappa_{\rm F}\log T$.

\subsection{Multiplicity of the one-step clock points}
Let $X_1=T^{c_1}$ with $c_1=(\log\log T)^{-1/2}$, $x^{(1)}_k=\tau_k^{X_1}=\tau_k-y_k$ with $y_k=S_{X_1}(\tau_k)/\lL(\tau_k)$, $\mathfrak b_T=2\pi/\log(3T)$ as in Lemma~\CGref{lem:maximal_primes}, and
\[
\tilde n(u)=\#\{k\in\IT:\ u-\mathfrak b_T\le x^{(1)}_k\le u\}\qquad(u\in\R).
\]
Note that $|y_k|\le\sigma_{X_1}/\lL(\tau_k)\ll T^{c_1/2}$, so $x^{(1)}_k\in[T-T^{1/2},2T+T^{1/2}]$ for $T\ge T_0$.

\begin{lemma}[Multiplicity]\label{lem:multiplicity}
$\displaystyle\int_\R\tilde n(u)^2\,du=\sum_{j,k\in\IT}\bigl(\mathfrak b_T-|x^{(1)}_j-x^{(1)}_k|\bigr)_+\ll\mathfrak b_T\,\#\IT .$
\end{lemma}

\begin{proof}
The identity: $\tilde n=\sum_k\mathbf 1_{[x^{(1)}_k,x^{(1)}_k+\mathfrak b_T]}$, and the measure of the intersection of two such intervals is $(\mathfrak b_T-|x^{(1)}_j-x^{(1)}_k|)_+$. Write $b=\mathfrak b_T$ and $\mathsf M=\#\IT$. The pairs with $|j-k|\le1$ contribute at most $3b\mathsf M$. Let $j=k+m$ with $m\ge2$ and $|x^{(1)}_{k+m}-x^{(1)}_k|\le b$. Since $\tau_{k+m}-\tau_k\ge mb$ (proof of Lemma~\CGref{lem:maximal_primes}), this forces
\begin{equation}
|y_{k+m}-y_k|\ge(m-1)b .
\label{eq:close_pair}
\end{equation}
Two consequences of \CGeqref{eq:close_pair} are counted separately; a close pair at distance $m$ satisfies both.

\emph{Variation.} $y=S_{X_1}/\lL$ is $C^\infty$ with $y'=S'_{X_1}/\lL-S_{X_1}\lL'/\lL^2$, and on $J_{k,m}=[\tau_k,\tau_{k+m}]$, of length $|J_{k,m}|\le mb(1+o(1))$,
\[
|y_{k+m}-y_k|\le\frac{1+o(1)}{\lL(T)}\int_{J_{k,m}}|S'_{X_1}|+\frac{\sigma_{X_1}|J_{k,m}|}{2\pi(T-1)\lL(T)^2},
\]
and the last term is $\ll T^{c_1/2-1}mb/\log^2T=o((m-1)b)$ uniformly in $m\ge2$. As $b\lL(T)=1-o(1)$, \CGeqref{eq:close_pair} gives $\frac1{|J_{k,m}|}\int_{J_{k,m}}|S'_{X_1}|\ge\frac13\lL(T)$ for $T\ge T_0$. Each $t$ lies in at most $m$ of the intervals $J_{k,m}$ ($k$ varying), and $|J_{k,m}|\ge mb$; by Chebyshev's inequality and Jensen's inequality,
\[
\begin{aligned}
\#\{k:\text{variation}\}
&\le\frac9{\lL(T)^2}\sum_k\frac1{|J_{k,m}|}\int_{J_{k,m}}|S'_{X_1}|^2\\
&\le\frac9{\lL(T)^2\,b}\int_{T-1}^{2T+1}|S'_{X_1}|^2
\ll T\log T\,c_1^2\asymp \mathsf M\,c_1^2 .
\end{aligned}
\]
by \CGeqref{eq:SX_derivative}, uniformly in $m$.

\emph{Size.} \CGeqref{eq:close_pair} gives $\max(|y_k|,|y_{k+m}|)\ge\frac{m-1}2b$, hence $\max(|S_{X_1}(\tau_k)|,|S_{X_1}(\tau_{k+m})|)\ge\frac{m-1}3$ for $T\ge T_0$. Sampling at the clock points as in the proof of Lemma~\CGref{lem:maximal_primes} (Sobolev on the disjoint intervals $[\tau_k,\tau_k+b]$), $\langle S_{X_1}(\tau_k)^4\rangle_{\IT}\ll\frac1T\int S_{X_1}^4+\frac1{T\log^4T}\int|S'_{X_1}|^4\ll(1+\log\log T)^2$, the first integral by Theorem~\CGref{thm:known_MV} with products of two primes ($\sum_{p\le X_1}p^{-1}\ll\log\log T$) and the second by Lemma~\CGref{lem:SX_fourth}. By Chebyshev's inequality of order $4$, $\#\{k:\text{size}\}\ll \mathsf M(\log\log T)^2/(m-1)^4$.

\emph{Sum.} With $H_0=\log\log T$ and $c_1^2=H_0^{-1}$, the number of close ordered pairs at distance $m\ge2$ is $\ll \mathsf M\min(H_0^{-1},H_0^2(m-1)^{-4})$, and
\[
\sum_{m\ge2}\min\Bigl(H_0^{-1},\frac{H_0^2}{(m-1)^4}\Bigr)\le\sum_{m-1\le H_0^{3/4}}H_0^{-1}+\sum_{m-1>H_0^{3/4}}\frac{H_0^2}{(m-1)^4}\ll H_0^{-1/4}.
\]
Hence the number of pairs $(j,k)$ with $|x^{(1)}_j-x^{(1)}_k|\le b$ is $\ll \mathsf M$, and each contributes at most $b$.
\end{proof}

\begin{lemma}[Sampling at points of bounded multiplicity]\label{lem:sampling_rho}
Let $\mathsf x_k$, $k\in\IT$, be real points of $[T-2T^{1/2},2T+2T^{1/2}]$ with multiplicity function $\tilde n(u)=\#\{k\in\IT:u-\mathfrak b_T\le\mathsf x_k\le u\}$ satisfying $\int_\R\tilde n^2\le\mathfrak k\,\mathfrak b_T\#\IT$ for some $\mathfrak k\ge1$. Let $r\ge1$ be an integer, $2\le X\le T^{1/(4r)}$, $\varphi$ a cutoff equal to $1$ on $[T-2T^{1/2},2T+2T^{1/2}]$ as in Lemma~\CGref{lem:transport}, and $g=\mathcal M(\varphi S'_X)$. Then, with an absolute constant $C$,
\[
\|g(\mathsf x_k)\|_{4r,\IT}\le C\,\mathfrak k^{1/(8r)}\sqrt r\,(1+\log X).
\]
In particular $\langle g(x^{(1)}_k)^4\rangle_{\IT}\ll(1+\log X)^4$ at the one-step clock points $x^{(1)}_k$, for $2\le X\le T^{1/4}$.
\end{lemma}

\begin{proof}
Let $\mathsf F=g^{4r}$, locally absolutely continuous with $\mathsf F'=4rg^{4r-1}g'$ (Theorem~\CGref{thm:known_maximal}). For $x\in[u-b,u]$, $\mathsf F(x)\le\frac1b\int_{u-b}^uF+\int_{u-b}^u|\mathsf F'|$. Averaging over $u\in[\mathsf x_k,\mathsf x_k+b]$ and summing over $k$,
\[
\sum_{k\in\IT}\mathsf F(\mathsf x_k)\le\frac1b\int_\R\tilde n(u)\Bigl(\frac1b\int_{u-b}^uF+\int_{u-b}^u|\mathsf F'|\Bigr)du\le\frac1b\|\tilde n\|_2\bigl(\|\mathsf F\|_2+b\|\mathsf F'\|_2\bigr),
\]
by Cauchy--Schwarz and because the averaging operator $h\mapsto\frac1b\int_{u-b}^uh$ has norm $\le1$ on $L^2(\R)$. Here $\|\mathsf F\|_2=\|g\|_{8r}^{4r}\le(\mathcal A_{8r}\|\varphi S'_X\|_{8r})^{4r}\le(C\sqrt r(1+\log X))^{4r}T^{1/2}$ and $\|\mathsf F'\|_2\le4r\|g\|_{8r}^{4r-1}\|g'\|_{8r}\le4r(1+\log X)\,(C\sqrt r(1+\log X))^{4r}T^{1/2}$, by Theorem~\CGref{thm:known_maximal} ($p=8r$, $\mathcal A_p$ bounded) and Lemma~\CGref{lem:higher_moments} (third bound with $2r$ in place of $r$, which needs $X^{4r}\le T$). Since $b\cdot4r(1+\log X)\le8\pi r\,c\,(1+o(1))\le7$ for $c=\log X/\log T\le1/(4r)$, we get $\sum_kF(\mathsf x_k)\le8\,b^{-1}(\mathfrak k b\mathsf M)^{1/2}T^{1/2}(C\sqrt r(1+\log X))^{4r}\le C'\sqrt{\mathfrak k}\,\mathsf M\,(C\sqrt r(1+\log X))^{4r}$, since $b^{-1}T\asymp\mathsf M=\#\IT$; take the $4r$-th root. For the $x^{(1)}_k$, $\mathfrak k\ll1$ by Lemma~\CGref{lem:multiplicity}.
\end{proof}

\subsection{The second step}
\begin{theorem}[Second Newton-type step]\label{thm:second_step}
Let $H_0=\log\log T$, $H_1=\log H_0$, $H_2=\log H_1$, $X_1=T^{H_0^{-1/2}}$, $X_2=T^{H_1^{-1/2}}$, and let $x^{(2)}_k$ be given by \CGeqref{eq:newton_iteration}. Then
\begin{equation}
\begin{split}
&\bigl\langle\lL(x^{(1)}_k)^2e_{2,k}^2\bigr\rangle_{\IT}=\frac1{4\pi^2}H_2+O\bigl(1+\sqrt{H_2}\bigr),\\
&\frac1{\#\IT}\sum_{k\in\IT}\bigl(x^{(2)}_k-\gamma_k\bigr)^2=\frac{\log\log\log\log T+O\bigl(\sqrt{\log\log\log\log T}\bigr)}{\log^2T}.
\end{split}
\label{eq:second_step}
\end{equation}
\end{theorem}

\begin{proof}
\emph{Fourth moment of $e_1$.} By \CGeqref{eq:recursion_identity} with $j=1$ (that is \CGeqref{eq:r_decomposition}), $\lL(\tau_k)e_{1,k}=u_k(X_1)+\mathfrak t_k(X_1)+\chi_k$. Lemma~\CGref{lem:higher_moments} with $r=2$ gives $\langle u_k(X_1)^4\rangle\ll(1+\log(1/c_1))^2\ll H_1^2$. For the transport, $|\mathfrak t_k|\le2|\delta_k|\mathsf g_1(\tau_k)$ with $\mathsf g_1=\mathcal M(\varphi S'_{X_1})$, so $\langle\mathfrak t_k^4\rangle\le16\langle\delta_k^8\rangle^{1/2}\langle \mathsf g_1(\tau_k)^8\rangle^{1/2}$; here $\langle\delta_k^8\rangle\ll H_0^4/\log^8T$ by \CGeqref{eq:s_moments} with $j=4$, and $\langle \mathsf g_1(\tau_k)^8\rangle\ll(c_1\log T)^8$ by the sampling argument of Lemma~\CGref{lem:transport} in $L^8$ (Sobolev on the intervals $\mathcal U_k$ with $g^8$, Theorem~\CGref{thm:known_maximal} with $p=8$, Lemma~\CGref{lem:higher_moments} with $r=2$). Hence $\langle\mathfrak t_k^4\rangle\ll(c_1^2H_0)^2=1$, and $\chi_k\ll1/T$. Dividing by $\lL(\tau_k)^4\asymp\log^4T$,
\begin{equation}
\langle e_{1,k}^4\rangle_{\IT}\ll\frac{H_1^2}{\log^4T}.
\label{eq:e1_fourth}
\end{equation}
\emph{Transport of the second step.} The interval centred at $x^{(1)}_k$ of radius $|e_{1,k}|$ contains $\gamma_k$ and lies in $[T-2T^{1/2},2T+2T^{1/2}]$, so $|S_{X_2}(\gamma_k)-S_{X_2}(x^{(1)}_k)|\le2|e_{1,k}|\mathsf g_2(x^{(1)}_k)$ with $\mathsf g_2=\mathcal M(\varphi S'_{X_2})$, and by Cauchy--Schwarz, \CGeqref{eq:e1_fourth} and Lemma~\CGref{lem:sampling_rho} (applicable since $X_2\le T^{1/8}$),
\[
\bigl\langle(S_{X_2}(\gamma_k)-S_{X_2}(x^{(1)}_k))^2\bigr\rangle\le4\langle e_{1,k}^4\rangle^{1/2}\langle \mathsf g_2(x^{(1)}_k)^4\rangle^{1/2}\ll\frac{H_1}{\log^2T}(c_2\log T)^2=c_2^2H_1=1 .
\]
\emph{Conclusion.} By \CGeqref{eq:recursion_identity} with $j=2$, $0\le\chi_{2,k}\le e_{1,k}^2/(2\pi T)$, and $\langle\chi_{2,k}^2\rangle\ll T^{-2}$. Theorem~\CGref{thm:main} with $X=X_2$ (allowed, $2\le X_2\le T^{1/2}$) gives $\langle u_k(X_2)^2\rangle=\frac1{2\pi^2}\log\frac1{c_2}+O(1+\sqrt{\log(1/c_2)})$ with $\log\frac1{c_2}=\frac12H_2$. Combining as in the proof of Theorem~\CGref{thm:computable_clock} gives the first statement of \CGeqref{eq:second_step}, and $\lL(x^{(1)}_k)=\frac{\log T}{2\pi}(1+O(\frac1{\log T}))$ on $\IT$, since $x^{(1)}_k\in[T-T^{1/2},2T+T^{1/2}]$, gives the second.
\end{proof}

\subsection{All further steps: the exact identity, the local deformation, and the multiplicity}
\label{subsec:further_steps}
This subsection proves the Newton-type steps $m\ge3$ with constants uniform in $m$, up to $m=\log^*T+O(1)$ steps, where the error becomes $O(1)$ in units of zeros.

\begin{lemma}[The second-level identity]\label{lem:second_identity}
For $j\ge1$, with $y^{(j)}_k=\tau_k-x^{(j)}_k$ and $\chi''_{j,k}=\int_{x^{(j-1)}_k}^{\tau_k}(\lL(\tau_k)-\lL(y))\,dy$,
\begin{equation}
\begin{split}
&y^{(j)}_k=\Bigl(1-\frac{\lL(\tau_k)}{\lL(x^{(j-1)}_k)}\Bigr)y^{(j-1)}_k+\frac{S_{X_j}(x^{(j-1)}_k)+\chi''_{j,k}}{\lL(x^{(j-1)}_k)},\\
&\text{hence}\qquad\lL(\tau_k)y^{(j)}_k=S_{X_j}(x^{(j-1)}_k)+\mathfrak r_{j,k},
\end{split}
\label{eq:Y_identity}
\end{equation}
with $|\mathfrak r_{j,k}|\ll T^{c_j-1}$ if $|y^{(j-1)}_k|\le T^{c_{j-1}/2}$ and $c_{j-1}\le c_j$ (true by \CGeqref{eq:sigma_M_def} and induction): the three contributions are $\lL(\tau_k)(1-\lL(\tau_k)/\lL(x))y^{(j-1)}_k\ll(y^{(j-1)}_k)^2/T$, $(\lL(\tau_k)/\lL(x)-1)S_{X_j}(x)\ll|y^{(j-1)}_k|\sigma_{X_j}/(T\log T)$ and $\chi''_{j,k}\le(y^{(j-1)}_k)^2/(4\pi(T-T^{1/2}))$. Newton's step does not accumulate the corrections of the previous levels: it replaces the previous one by the phase of the new cutoff evaluated at the previous point.
\end{lemma}

\begin{proof}
$\NO(x^{(j-1)}_k)-(k-\frac12)=\NO(x^{(j-1)}_k)-\NO(\tau_k)=-\lL(\tau_k)y^{(j-1)}_k+\chi''_{j,k}$, with $0\le\chi''_{j,k}\le(y^{(j-1)}_k)^2/(4\pi(T-T^{1/2}))$; insert in \CGeqref{eq:newton_iteration} and use $1-\lL(\tau_k)/\lL(x)=O(|y^{(j-1)}_k|/(T\log T))$ and $|S_{X_j}|\le\sigma_{X_j}$.
\end{proof}

\begin{proposition}[Local deformation of the second clock]\label{prop:deformation}
Let $X_1=T^{c_1}$, $X_2=T^{c_2}$ with $c_1=H_0^{-1/2}$, $c_2=H_1^{-1/2}$, and $d_{j,k}=x^{(j)}_{k+1}-x^{(j)}_k$ ($d_{0,k}=\tau_{k+1}-\tau_k$). Then, for $T\ge T_0$,
\begin{equation}
\Bigl\langle\Bigl|\frac{d_{2,k}-d_{0,k}}{d_{0,k}}\Bigr|^4\Bigr\rangle_{\IT}\ll c_2^4=\frac1{H_1^2},\qquad
\frac1{\#\IT}\#\Bigl\{k\in\IT:\Bigl|\frac{d_{2,k}}{d_{0,k}}-1\Bigr|>\alpha\Bigr\}\ll\frac1{\alpha^4H_1^2}\quad(0<\alpha\le1),
\label{eq:deformation}
\end{equation}
in particular $\#\{k\in\IT:d_{2,k}\le0\}\ll\#\IT/(\log\log\log T)^2$: the second clock preserves the order and a spacing comparable to that of the RvM clock outside a proportion $O(H_1^{-2})$ of the indices.
\end{proposition}

\begin{proof}
By \CGeqref{eq:Y_identity} with $j=2$, $d_{2,k}-d_{0,k}=-\nabla_k\bigl(S_{X_2}(x^{(1)})/\lL(\tau)\bigr)-\nabla_k(\mathfrak r_2/\lL(\tau))$, where $\mathfrak r_2$ is the sequence $\mathfrak r_{2,k}$ of Lemma~\CGref{lem:second_identity}, $\nabla_ka=a_{k+1}-a_k$ is the forward difference ($\Delta_k$ is reserved for the Gram displacement); the last term is $\ll T^{c_1-1}$, negligible against $d_{0,k}\asymp1/\log T$, and $\nabla_k(1/\lL(\tau))\ll1/(T\log^3T)$ is negligible too, so it suffices to bound $\nabla_kS_{X_2}(x^{(1)})/(d_{0,k}\lL(\tau_k))$ in $\ell^4(\IT)$. Let $\mathsf g_j=\mathcal M(\varphi S'_{X_j})$ with $\varphi=1$ on $[T-2T^{1/2},2T+2T^{1/2}]$. Two applications of $|S(a)-S(b)|\le2|a-b|\,\mathcal M(S')(a)$:
\[
|x^{(1)}_{k+1}-x^{(1)}_k|\le d_{0,k}+|y_{k+1}-y_k|\le d_{0,k}\Bigl(1+\frac{2\mathsf g_1(\tau_k)}{\lL(\tau_k)}\Bigr)+O\bigl(T^{c_1/2-1}\bigr),
\]
\[
|\nabla_kS_{X_2}(x^{(1)})|\le2|x^{(1)}_{k+1}-x^{(1)}_k|\,\mathsf g_2(x^{(1)}_k).
\]
Hence, with $\lL(\tau_k)\asymp\log T$,
\[
\Bigl\|\frac{\nabla S_{X_2}(x^{(1)})}{d_0\,\lL(\tau)}\Bigr\|_{4,\IT}\ll\Bigl\|\frac{\mathsf g_2(x^{(1)})}{\log T}\Bigr\|_{4,\IT}+\Bigl\|\frac{\mathsf g_2(x^{(1)})}{\log T}\Bigr\|_{8,\IT}\Bigl\|\frac{\mathsf g_1(\tau)}{\log T}\Bigr\|_{8,\IT}+o(1).
\]
Sampling at the $\tau_k$ in $L^8$ (Sobolev on the disjoint intervals $\mathcal U_k$ as in Lemma~\CGref{lem:transport}, with Lemma~\CGref{lem:higher_moments} for $\int|S'_{X_1}|^8$, $\int|S''_{X_1}|^8$ and Theorem~\CGref{thm:known_maximal} with $p=8$) gives $\|\mathsf g_1(\tau)/\log T\|_8\ll c_1$. Sampling at the points $x^{(1)}_k$ in $L^8$ is Lemma~\CGref{lem:sampling_rho} with $\mathsf x_k=x^{(1)}_k$ and $\mathsf F=\mathsf g_2^8$: $\|\mathsf F\|_2=\|\mathsf g_2\|_{16}^8$ and $\|\mathsf F'\|_2\le8\|\mathsf g_2\|_{16}^7\|\mathsf g_2'\|_{16}$, which need Theorem~\CGref{thm:known_MV} with products of $8$ primes ($X_2^{16}\le T$, true for $T\ge T_0$ since $c_2\to0$) and Theorem~\CGref{thm:known_maximal} with $p=16$; the result is $\|\mathsf g_2(x^{(1)})/\log T\|_8\ll c_2$. The first inequality of \CGeqref{eq:deformation} follows, and the second is Markov's inequality; the last statement is $\alpha=1$.
\end{proof}

Proposition~\CGref{prop:deformation} controls consecutive differences; the multiplicity of the points $x^{(j)}_k$ at scale $\mathfrak b_T$, which Lemma~\CGref{lem:sampling_rho} needs, is obtained for every $j$ by comparing $x^{(j)}_k$ with $\gamma_k$ through two unconditional counts of near pairs and near triples of ordinates. The ordinates enter the proof of the bound, as in Theorem~\CGref{thm:main}, and not the construction of the points.

\begin{lemma}[Near pairs and triples of ordinates]\label{lem:zero_pairs}
Let $b=\mathfrak b_T$, $\mathsf M=\#\IT$ and, for $\mathsf R\ge1$ and real $u$, $n_{\mathsf R}(u)=\#\{l\in\IT:|\gamma_l-u|\le\mathsf R b\}$. There are absolute constants $C_1,C_2$ such that for $T\ge T_0$ and every $\mathsf R\ge1$,
\begin{equation}
\#\{(k,l)\in\IT^2:|\gamma_k-\gamma_l|\le\mathsf R b\}\le C_1\mathsf R\mathsf M,\qquad\sum_{k\in\IT}n_{\mathsf R}(\gamma_k)^2\le C_2\mathsf R^2\mathsf M .
\label{eq:zero_pairs}
\end{equation}
\end{lemma}

\begin{proof}
If $4\mathsf R b>1$, then $N(u+h)-N(u)\le Ch\log T$ for $1\le h\le T$ and $u\in[T-1,2T]$, by \CGeqref{eq:N_decomposition} and Theorem~\CGref{thm:known_basic}. Since only $l\in\IT$ are counted, $n_{\mathsf R}(\gamma_k)$ is at most the number of ordinates in an interval of length $h=\min(2\mathsf R b,T)$ contained in $[T-1,2T]$, which is $\le Ch\log T\le C'\mathsf R$ (for $h=T$ use $T\log T\le2\mathsf R b\log T\le4\pi\mathsf R$); both bounds follow. Let $4\mathsf R b\le1$ and $h\in\{2\mathsf R b,4\mathsf R b\}$. For $u\in[T-1,2T]$, $N(u+h)-N(u)=h\lL(u)+S(u+h)-S(u)+O(1/T)$ with $h\lL(u)\le5\mathsf R$, so \CGeqref{eq:fujii_difference} and $\log(2+h\log T)\le\log(30\mathsf R)\le30\mathsf R$ give
\[
\int_{T-1}^{2T}\bigl(N(u+h)-N(u)\bigr)^{2k}du\le C\mathsf R^{2k}T\qquad(k=1,2).
\]
A pair $(k,l)$ with $|\gamma_k-\gamma_l|\le\mathsf R b$ lies in $[u,u+2\mathsf R b]$ for every $u\in[\min(\gamma_k,\gamma_l)-\mathsf R b,\min(\gamma_k,\gamma_l)]$, a set of measure $\mathsf R b$ contained in $[T-1,2T]$; hence the number of pairs is at most $\frac1{\mathsf R b}\int_{T-1}^{2T}(N(u+2\mathsf R b)-N(u))^2du\le C\mathsf R T/b\le C_1\mathsf R\mathsf M$, as $\mathsf M\asymp T/b$. A triple $(k,l,l')$ with $|\gamma_l-\gamma_k|,|\gamma_{l'}-\gamma_k|\le\mathsf R b$ lies in $[u,u+4\mathsf R b]$ for $u$ in a set of measure $2\mathsf R b$; by H\"older's inequality the number of triples is at most $\frac1{2\mathsf R b}\int_{T-1}^{2T}(N(u+4\mathsf R b)-N(u))^3du\le\frac{T^{1/4}(C\mathsf R^4T)^{3/4}}{2\mathsf R b}\le C_2\mathsf R^2\mathsf M$; the triples with $l,l'\in\IT$ are among them.
\end{proof}

\begin{theorem}[Multiplicity of the corrected clock points]\label{thm:multiplicity_all}
Let $j\ge0$ and let $x^{(j)}_k$ be given by \CGeqref{eq:newton_iteration} with cutoffs $X_i=T^{c_i}$, $c_i\le\frac12$, so that $x^{(j)}_k\in[T-2T^{1/2},2T+2T^{1/2}]$ for $k\in\IT$ and $T\ge T_0$ (Lemma~\CGref{lem:second_identity}). Let $\tilde n_j(u)=\#\{k\in\IT:u-\mathfrak b_T\le x^{(j)}_k\le u\}$, and let $\mathsf A\ge1$, $\mathsf H\ge1$ be such that $\|\lL(\tau_k)e_{j,k}\|_{4,\IT}\le \mathsf A\sqrt \mathsf H$. Then, with an absolute constant $C_3$,
\begin{equation}
\int_\R\tilde n_j(u)^2\,du\le C_3\,\mathsf A\sqrt \mathsf H\;\mathfrak b_T\,\#\IT .
\label{eq:multiplicity_all}
\end{equation}
\end{theorem}

\begin{proof}
Write $b=\mathfrak b_T$, $\mathsf M=\#\IT$, $x_k=x^{(j)}_k$, $e_k=e_{j,k}$. As in Lemma~\CGref{lem:multiplicity}, $\int\tilde n_j^2=\sum_{k,l\in\IT}(b-|x_k-x_l|)_+\le b\,\mathsf P$ with $\mathsf P=\#\{(k,l)\in\IT^2:|x_k-x_l|<b\}$. Let $\mathsf R_0\ge2$. The pairs with $|\gamma_k-\gamma_l|\le\mathsf R_0b$ number at most $C_1\mathsf R_0\mathsf M$ by \CGeqref{eq:zero_pairs}. For the others, let $\mathsf R=2^i\mathsf R_0$, $i\ge0$, and consider the pairs with $\mathsf R b<|\gamma_k-\gamma_l|\le2\mathsf R b$ and $|x_k-x_l|<b$. For them $|e_k-e_l|\ge(\mathsf R-1)b\ge\mathsf R b/2$, and $b\lL(\tau_k)\ge\frac12$ for $T\ge T_0$, so $\lL(\tau_k)|e_k-e_l|\ge\mathsf R/4$ and
\[
1\le\Bigl(\frac{4\lL(\tau_k)|e_k-e_l|}{\mathsf R}\Bigr)^2\le\frac{32}{\mathsf R^2}\Bigl(\lL(\tau_k)^2e_k^2+\lL(\tau_l)^2e_l^2\Bigr)\bigl(1+O(1/\log T)\bigr).
\]
Summing over the ordered pairs of the block and symmetrizing, their number $\mathsf P_{\mathsf R}$ satisfies
\[
\mathsf P_{\mathsf R}\le\frac{65}{\mathsf R^2}\sum_{k\in\IT}\lL(\tau_k)^2e_k^2\,n_{2\mathsf R}(\gamma_k)\le\frac{65}{\mathsf R^2}\Bigl(\sum_k\lL(\tau_k)^4e_k^4\Bigr)^{1/2}\Bigl(\sum_kn_{2\mathsf R}(\gamma_k)^2\Bigr)^{1/2}\le\frac{130\sqrt{C_2}\,\mathsf A^2\mathsf H}{\mathsf R}\,\mathsf M,
\]
by Cauchy--Schwarz, the hypothesis and \CGeqref{eq:zero_pairs} (both $k$ and $l$ lie in $\IT$, so $n_{2\mathsf R}$ counts every $l$ of the block). Summing over $i\ge0$ gives at most $260\sqrt{C_2}\,\mathsf A^2\mathsf H\mathsf M/\mathsf R_0$. With $\mathsf R_0=\mathsf A\sqrt \mathsf H\,(260\sqrt{C_2}/C_1)^{1/2}$ (which is $\ge2$ for $C_2$ large enough) the two parts are equal and $\mathsf P\le2\sqrt{260C_1\sqrt{C_2}}\,\mathsf A\sqrt \mathsf H\,\mathsf M$.
\end{proof}

Only the fourth moment of the error of the $j$-th clock enters. For $j=2$, with $\mathsf A$ bounded (shown below), the multiplicity of the second clock is $\ll\sqrt{H_2}\,\mathfrak b_T\#\IT$. The factor $\sqrt \mathsf H$ is not a loss in the iteration: it is absorbed by the cutoff of the next step.

\subsubsection*{All further steps}
Let $H_0=\log\log T$, $c_1=H_0^{-1/2}$, $c_2=H_1^{-1/2}$ with $H_1=\log H_0$, and define recursively, with an absolute constant $0<C_\star\le1$ fixed in the proof of Theorem~\CGref{thm:uniform_iteration},
\begin{equation}
H_j=2\log\frac1{c_j}\quad(j\ge1),\qquad c_j=C_\star\,H_{j-1}^{-9/16}\quad(j\ge3),
\label{eq:cutoff_law}
\end{equation}
so that $H_2=\log H_1$ as in Theorem~\CGref{thm:second_step}, $H_j=\frac98\log H_{j-1}+2\log(1/C_\star)$ for $j\ge3$, and $\log H_{j-1}\le H_j$ for all $j$. All cutoffs depend only on $T$. Put
\[
A_j(r)=\frac{\|\lL(\tau_k)e_{j,k}\|_{2r,\IT}}{\sqrt{H_j}}\quad(j\ge1),\qquad A_0(r)=\frac{\|\lL(\tau_k)\delta_k\|_{2r,\IT}}{\sqrt{H_0}},
\]
and let $m^*=m^*(T)$ be the first index $m$ with $H_m<H^*$, where $H^*$ is the absolute constant of the next theorem. Since $\log H_{j-1}\le H_j\le\frac98\log H_{j-1}+2\log(1/C_\star)$, the sequence $H_j$ decreases from $\log\log T$ to below $H^*$ in $m^*=\log^*T+O(1)$ steps, where $\log^*$ is the iterated logarithm.

\begin{theorem}[All Newton-type steps, uniformly]\label{thm:uniform_iteration}
There are absolute constants $C_\star$, $H^*$, $C$ and $T_0$ such that, with the cutoffs \CGeqref{eq:cutoff_law}, for every $T\ge T_0$ and every $1\le m\le m^*$,
\begin{equation}
\bigl\langle\lL(\tau_k)^2e_{m,k}^2\bigr\rangle_{\IT}=\frac{H_m}{4\pi^2}+\vartheta_m\,C\bigl(1+\sqrt{H_m}\bigr),\qquad|\vartheta_m|\le1,
\label{eq:uniform_iteration}
\end{equation}
and $\|\lL(\tau_k)e_{j,k}\|_{4,\IT}\le C\sqrt{H_j}$ for $1\le j<m^*$. In particular, at $m=m^*=\log^*T+O(1)$,
\begin{equation}
\frac1{\#\IT}\sum_{k\in\IT}\bigl(x^{(m^*)}_k-\gamma_k\bigr)^2\le\frac{C'}{\log^2T}:
\label{eq:final_clock}
\end{equation}
the clock corrected $\log^*T+O(1)$ times with the Euler phases of the primes up to $T^{c_j}$ locates the ordinates with a mean-square error of $O(1)$ mean spacings, with an absolute constant, using no ordinate in its construction. The cutoffs \CGeqref{eq:cutoff_law} and the number $m^*$ of steps involve the absolute constants $C_\star$ and $H^*$ fixed in the proof, which are not computed here; once they are fixed, every $x^{(m)}_k$ is an explicit expression in $k$, $\tau_k$ and the primes (Question~\CGref{q:explicit}).
\end{theorem}

\begin{proof}
\emph{Multiplicities.} Let $\mathfrak k_j$ be the least $\mathfrak k\ge1$ with $\int\tilde n_j^2\le\mathfrak k\,\mathfrak b_T\#\IT$. Then $\mathfrak k_0=1$ (the $\tau_k$ are separated by at least $\mathfrak b_T$, so only the diagonal contributes), $\mathfrak k_1\le C$ (Lemma~\CGref{lem:multiplicity}) and $\mathfrak k_j\le C_3\max(A_j(2),1)\sqrt{H_j}$ for $j\ge2$ (Theorem~\CGref{thm:multiplicity_all} with $\mathsf A=\max(A_j(2),1)$ and $\mathsf H=H_j$).

\emph{Base.} By Lemma~\CGref{lem:higher_moments} with $X=2$, $|S_2|\le\sigma_2<1$ and $0\le\chi_k\le1/T$, $\|\lL(\tau_k)\delta_k\|_{2r}\le\|u_k(2)\|_{2r}+2\le C_0(r^2+\sqrt{rH_0})$, so $A_0(r)\le C_0r^2$ for $r\ge1$.

\emph{One step.} Let $j\ge1$ and $r\ge1$ with $X_j\le T^{1/(160r)}$ and $2r\le\kappa_{\rm F}\log T$. By \CGeqref{eq:recursion_norms}, $\lL(x^{(j-1)}_k)=\lL(\tau_k)(1+O(T^{-1/2}))$, Lemma~\CGref{lem:higher_moments} (with $V\le\log\frac1{c_j}+O(1)\le H_j$), Lemma~\CGref{lem:sampling_rho} at the points $x^{(j-1)}_k$ with multiplicity constant $\mathfrak k_{j-1}$, and $\|\chi_j\|_{2r}\ll T^{c_{j-1}-1}$,
\[
\sqrt{H_j}\,A_j(r)\le C\bigl(r^2+\sqrt{rH_j}\bigr)+2\sqrt{H_{j-1}}\,A_{j-1}(2r)\cdot\frac{C\mathfrak k_{j-1}^{1/(8r)}\sqrt r\,(1+\log X_j)}{\lL(\tau_k)}\bigl(1+o(1)\bigr)+o(1),
\]
and $(1+\log X_j)/\lL(\tau_k)\le2\pi c_j(1+o(1))$, $\mathfrak k_{j-1}^{1/(8r)}\le\mathfrak k_{j-1}^{1/8}$. For $j\le2$, $\sqrt{H_{j-1}}c_j=1$ and $\mathfrak k_{j-1}\le C$; for $j\ge3$, $\mathfrak k_{j-1}^{1/8}\sqrt{H_{j-1}}c_j\le(C_3\max(A_{j-1}(2),1))^{1/8}H_{j-1}^{1/16+1/2-9/16}C_\star=C_\star(C_3\max(A_{j-1}(2),1))^{1/8}$, and $\max(A_{j-1}(2),1)\le\max(A_{j-1}(2r),1)$ (monotonicity of the discrete norms), so that $A_{j-1}(2r)\max(A_{j-1}(2r),1)^{1/8}=\max(A_{j-1}(2r),A_{j-1}(2r)^{9/8})$. Fix $C_\star=\min(1,(4\pi C\,C_3^{1/8})^{-1})$. Then, for all $j\ge1$ and $r\ge1$,
\begin{equation}
A_j(r)\le C\Bigl(\sqrt r+\frac{r^2}{\sqrt{H_j}}\Bigr)+\frac{C_5\sqrt r}{\sqrt{H_j}}\max\bigl(A_{j-1}(2r),A_{j-1}(2r)^{9/8}\bigr),
\label{eq:uniform_recursion}
\end{equation}
with absolute $C,C_5$, provided $X_j\le T^{1/(160r)}$ and $2r\le\kappa_{\rm F}\log T$.

\emph{Induction.} Let $B_\star=\max(2C,C_0,1)$ and $C_6=\bigl(4C/B_\star+4C_5(4B_\star)^{9/8}/B_\star\bigr)^2$. We claim that for every $m\le m^*$,
\begin{equation}
A_j(r)\le B_\star r^2\qquad(0\le j\le m-1,\ 1\le r\le2^{m-j}).
\label{eq:induction_claim}
\end{equation}
For $j=0$ this is the base. Suppose it holds at $j-1$ for $r\le2^{m-j+1}$, let $1\le r\le2^{m-j}$, and assume $H_j\ge C_64^{m-j}$, $X_j\le T^{1/(160r)}$ and $2r\le\kappa_{\rm F}\log T$. Since $A_{j-1}(2r)\le4B_\star r^2$ and $4B_\star r^2\ge1$, \CGeqref{eq:uniform_recursion} gives
\[
A_j(r)\le C\sqrt r+\frac{Cr^2+C_5(4B_\star)^{9/8}r^{11/4}}{\sqrt{H_j}}\le\frac{B_\star}2r^2+\frac{\bigl(C+C_5(4B_\star)^{9/8}\bigr)r^{3}}{\sqrt{H_j}}\le B_\star r^2,
\]
because $r^3/\sqrt{H_j}\le r^2\cdot2^{m-j}/\sqrt{H_j}\le r^2/\sqrt{C_6}$ and $\sqrt{C_6}\ge2(C+C_5(4B_\star)^{9/8})/B_\star$.

\emph{The three conditions.} Let $H^*$ be an absolute constant with $H^*\ge\max(4C_6,160^2\cdot4)$ and so large that $e^{\frac89(H-2\log(1/C_\star))}\ge4H$ and $e^H\ge4H$ for $H\ge H^*$. If $m\le m^*$ then $H_{m-1}\ge H^*$ (by definition of $m^*$), and from $H_{j-1}\ge e^{H_j}$ ($j\le2$) or $H_{j-1}=e^{\frac89(H_j-2\log(1/C_\star))}$ ($j\ge3$) we get $H_{j-1}\ge4H_j$ whenever $H_j\ge H^*$; by downward induction, $H_{m-1-i}\ge4^iH^*$ for $0\le i\le m-1$. Hence $H_j\ge4^{m-1-j}H^*\ge C_64^{m-j}$ for $1\le j\le m-1$, and $H_{j-1}\ge4^{m-j}H^*\ge160^2\,4^{m-j}$ for $1\le j\le m$, which gives $c_j\le H_{j-1}^{-1/2}\le2^{-(m-j)}/160$, i.e.\ $X_j\le T^{1/(160r)}$ for $r\le2^{m-j}$, and $X_j\ge2$ because $c_j\ge C_\star(\log\log T)^{-9/16}$. Finally $2r\le2^{m-j+1}\le2^{m^*+1}$, and $m^*=\log^*T+O(1)$ with an absolute $O(1)$, so $2r\le\kappa_{\rm F}\log T$ for $T\ge T_0$. So \CGeqref{eq:induction_claim} holds for every $m\le m^*$, which is the statement $\|\lL e_j\|_{4,\IT}\le4B_\star\sqrt{H_j}$ for $j<m^*$.

\emph{The last step.} For $m\le m^*$ and $r=1$, \CGeqref{eq:induction_claim} at $j=m-1$ gives $A_{m-1}(2)\le4B_\star$, so the transport term of the $m$-th step, which is already in units of zeros, satisfies $\|\mathfrak t^{(m)}_k\|_{2,\IT}\le C_5(4B_\star)^{9/8}$ (the transport bound in the proof of \CGeqref{eq:uniform_recursion}, with $r=1$), with $\mathfrak t^{(m)}_k=S_{X_m}(\gamma_k)-S_{X_m}(x^{(m-1)}_k)$. Theorem~\CGref{thm:main} at $X=X_m$ ($2\le X_m\le T^{1/160}$) gives $\langle u_k(X_m)^2\rangle=\frac1{2\pi^2}\log\frac1{c_m}+O(1+\sqrt{H_m})=\frac{H_m}{4\pi^2}+O(1+\sqrt{H_m})$, the cross term is at most $2\|u(X_m)\|_2\,C_5(4B_\star)^{9/8}\ll1+\sqrt{H_m}$, and $\chi_{m,k}\ll T^{-1/2}$; \CGeqref{eq:recursion_identity} gives \CGeqref{eq:uniform_iteration}. At $m=m^*$, $H_{m^*}<H^*$ and $\lL(\tau_k)=\frac{\log T}{2\pi}(1+O(1/\log T))$ give \CGeqref{eq:final_clock}.
\end{proof}

\begin{corollary}[Third step]\label{cor:third_step}
With $c_3=C_\star H_2^{-9/16}$ and $H_3=\frac98\log H_2+2\log(1/C_\star)$,
\[
\bigl\langle\lL(\tau_k)^2e_{3,k}^2\bigr\rangle_{\IT}=\frac{H_3}{4\pi^2}+O\bigl(1+\sqrt{H_3}\bigr),\qquad\frac1{\#\IT}\sum_{k\in\IT}\bigl(x^{(3)}_k-\gamma_k\bigr)^2=\frac{H_3+O(\sqrt{H_3})}{\log^2T},
\]
and the multiplicity of the second clock satisfies $\int\tilde n_2^2\ll\sqrt{H_2}\,\mathfrak b_T\#\IT$.
\end{corollary}

\begin{proof}
Theorem~\CGref{thm:uniform_iteration} with $m=3$ (for $T\ge T_0$, $H_2\ge H^*$), and Theorem~\CGref{thm:multiplicity_all} with $\mathsf A=A_2(2)\le4B_\star$ and $\mathsf H=H_2$.
\end{proof}

The induction \CGeqref{eq:induction_claim} bounds the moments of every order $2r\le2^{m^*-j+1}$ of the $j$-th error, not only its mean square, and the number of available moments grows with $T$. This gives tails.

\begin{corollary}[Tails of the corrected clocks]\label{cor:clock_tails}
With the cutoffs \CGeqref{eq:cutoff_law}, for $T\ge T_0$, every $1\le j<m^*$, every integer $1\le q\le2^{m^*-j}$ and every $\lambda>0$,
\begin{equation}
\#\bigl\{k\in\IT:\lL(\tau_k)|e_{j,k}|>\lambda\sqrt{H_j}\bigr\}\le\Bigl(\frac{B_\star q^2}{\lambda}\Bigr)^{2q}\#\IT .
\label{eq:clock_tails}
\end{equation}
For fixed $j$ and $q$ the right side is $\ll_q\lambda^{-2q}\#\IT$ for $T\ge T_0(j,q)$, since $2^{m^*-j}\to\infty$: for every fixed $q$, and $T$ large in terms of $q$, the $j$-th corrected clock misses its ordinate by more than $\lambda\sqrt{H_j}$ positions on a proportion $\ll_q\lambda^{-2q}$ of the indices. The orders available at a given $T$ are $q\le2^{m^*-j}$; nothing is claimed uniformly in $q$.
\end{corollary}

\begin{proof}
\CGeqref{eq:induction_claim} with $m=m^*$ gives $\|\lL(\tau_k)e_{j,k}\|_{2q,\IT}\le B_\star q^2\sqrt{H_j}$ for $j\le m^*-1$ and $q\le2^{m^*-j}$; apply Chebyshev's inequality to the $2q$-th moment.
\end{proof}

The Newton points of a level need not be in increasing order (Proposition~\CGref{prop:deformation} bounds the proportion of inversions of the second clock). Sorting the points of the final level does not increase the unweighted error.

\begin{lemma}[Sorting the final points]\label{lem:sorted}
Let $\gamma_1\le\dots\le\gamma_{\mathsf M}$, let $z_1,\dots,z_{\mathsf M}$ be real, and let $z^\uparrow$ be the same values in nondecreasing order. Then $\sum_k|z^\uparrow_k-\gamma_k|^q\le\sum_k|z_k-\gamma_k|^q$ for every $q\ge1$, and $\{z^\uparrow_k\}=\{z_k\}$ as multisets. In particular, sorting the points $x^{(j)}_k$, $k\in I$, of one level does not increase the unweighted means $\langle|x^{(j)}_k-\gamma_k|^q\rangle_I$ and does not change the multiplicity function $\tilde n_j$; it removes the strict inversions $d_{j,k}<0$ of Proposition~\CGref{prop:deformation}, but not the coincidences $d_{j,k}=0$: points that coincide remain coincident.
\end{lemma}

\begin{proof}
For $a\le b$, $x\le y$ and $\varphi$ convex, $\varphi(a-x)+\varphi(b-y)\le\varphi(a-y)+\varphi(b-x)$ (the points $a-y\le a-x,b-y\le b-x$ have the same sum in both pairings, and the outer pair majorizes the inner one). With $\varphi(t)=|t|^q$, each transposition of an inverted pair of $z$'s does not increase the sum, and finitely many transpositions sort $z$. Sorting permutes the values, so the multiset, and with it $\tilde n_j$, is unchanged, and equal values stay equal.
\end{proof}

The lemma is used only for the final level (Corollary~\CGref{cor:sorted_clock}). It does not show that the Newton-type iteration continued from sorted points satisfies Theorems~\CGref{thm:second_step}, \CGref{thm:multiplicity_all} and \CGref{thm:uniform_iteration}: their proofs use the weighted norms $\|\lL(\tau_k)e_{j,k}\|_{r,I}$, in which the weight is attached to the index, and the identity \CGeqref{eq:recursion_identity}, in which the index enters through $k-\frac12$; the effect of sorting on these is not controlled here.

\begin{corollary}[A nondecreasing clock built without ordinates]\label{cor:sorted_clock}
Let $\mathcal I_T^\tau=\{k:T\le\tau_k<2T\}$, a set of consecutive indices defined by the clock alone. The points $x^{(m^*)}_k$ of Theorem~\CGref{thm:uniform_iteration} are defined by \CGeqref{eq:newton_iteration} for every index $k$ with $\tau_k\ge T_0$, in particular for $k\in\mathcal I_T^\tau$; let $x^\uparrow_k$, $k\in\mathcal I_T^\tau$, be these values in nondecreasing order. Then $x^\uparrow_k\le x^\uparrow_{k+1}$, no ordinate enters the construction of the $x^\uparrow_k$ (they depend on the $\tau_k$, the primes and the constants $C_\star$, $H^*$ of Theorem~\CGref{thm:uniform_iteration}), and for $T\ge T_0$
\[
\frac1{\#\mathcal I_T^\tau}\sum_{k\in\mathcal I_T^\tau}\bigl(x^\uparrow_k-\gamma_k\bigr)^2\le\frac{C'+1}{\log^2T}.
\]
\end{corollary}

\begin{proof}
By Lemma~\CGref{lem:sorted} with $q=2$ on the index set $\mathcal I_T^\tau$, the sum is at most $\sum_{k\in\mathcal I_T^\tau}(x^{(m^*)}_k-\gamma_k)^2$. Theorem~\CGref{thm:uniform_iteration} is a statement over $\IT$; for the indices of $\mathcal I_T^\tau\setminus\IT$ only the localization of the points proved below is used. Since $|\delta_k|\le C_{10}$ for all $k$ (Theorem~\CGref{thm:known_basic}), the symmetric difference $\mathcal I_T^\tau\triangle\IT$ consists of indices with $|\gamma_k-T|\le C_{10}$ or $|\gamma_k-2T|\le C_{10}$, and has $O(\log T)$ elements by \CGeqref{eq:N_decomposition}. For every $k\in\mathcal I_T^\tau$ the points $x^{(j)}_k$ lie in $[T-2T^{1/2},2T+2T^{1/2}]$ and Lemma~\CGref{lem:second_identity} gives $|x^{(j)}_k-\tau_k|=|y^{(j)}_k|\le(\sigma_{X_j}+1)/\lL(\tau_k)\ll T^{1/320}/\log T$, because $X_j\le T^{1/160}$ for every $j\le m^*$ (proof of Theorem~\CGref{thm:uniform_iteration}, \emph{the three conditions}, with $r=1$) and $\sigma_X\ll\sqrt X/\log X$. Hence $|x^{(m^*)}_k-\gamma_k|\le C_{10}+O(T^{1/320}/\log T)$ for $k\in\mathcal I_T^\tau\triangle\IT$, these indices contribute $O(\log T\cdot T^{1/160}/\log^2T)=O(T^{1/160})$ to the sum, and \CGeqref{eq:final_clock} gives
\[
\sum_{k\in\mathcal I_T^\tau}\bigl(x^\uparrow_k-\gamma_k\bigr)^2\le\frac{C'\,\#\IT}{\log^2T}+O\bigl(T^{1/160}\bigr)\le\frac{(C'+1)\,\#\mathcal I_T^\tau}{\log^2T}
\]
for $T\ge T_0$, since $\#\mathcal I_T^\tau=\#\IT+O(\log T)$ and $\#\IT\asymp T\log T$.
\end{proof}

The cutoff of the third step, $T^{c_3}$ with $c_3\asymp H_2^{-9/16}$ instead of $H_2^{-1/2}$, is what the multiplicity costs: the exponent $\frac9{16}=\frac12+\frac1{16}$ comes from the factor $\sqrt{H_2}$ of \CGeqref{eq:multiplicity_all} through the exponent $\frac1{8r}\le\frac18$ of Lemma~\CGref{lem:sampling_rho}. The constants of Fujii's and Montgomery--Vaughan's bounds, which grow with the moment order $2^{m-j}$ used at level $j$, are absorbed by the normalization $H_j^{-1/2}$, because $H_j$ grows as a tower in $m-j$ while the orders grow geometrically.

A bounded error in units of zeros with a \emph{fixed} final cutoff $T^{c_0}$ is not obtained by this method: for a fixed cutoff exponent the error in units of zeros is bounded in mean square if and only if the transport is (Remark~\CGref{rem:newton}), and the only available bound for the transport, Lemma~\CGref{lem:transport}, grows like $c^2\log\log T$.

%======================================================================
\section{Theorem E: the pointwise envelope}
\label{sec:envelope}

Let $\mathfrak a=\limsup_{t\to\infty}|S(t)|/\log t$. By Theorem~\CGref{thm:known_BW}, $\mathfrak a\le0.10076$, and $\mathfrak a\le0.1013$ also follows from \cite{Trudgian2010}; whether $\mathfrak a=0$ is not known.

\begin{proposition}[The envelope of the clock is the envelope of $S$]\label{prop:envelope_equivalence}
(a) $\limsup_{k\to\infty}|\tau_k-\gamma_k|=2\pi\mathfrak a$. (b) If $X_k\ge2$ and $\sigma_{X_k}=o(\log\tau_k)$, in particular if $X_k\le(\log\tau_k)^{2-\varepsilon}$, then $\limsup_{k\to\infty}|\tau_k^{X_k}-\gamma_k|=2\pi\mathfrak a$.
\end{proposition}

\begin{proof}
(a) The values $S(\gamma_k^\pm)$ are limits of values of $S$ at non-ordinates. For $\varepsilon>0$, \CGeqref{eq:s_S_bound} gives $|s_k|\le(\mathfrak a+\varepsilon)\log\gamma_k+O_\varepsilon(1)$, and $|\delta_k|=|s_k|/\lL(\xi_k)$ with $\lL(\xi_k)=\frac{\log\gamma_k}{2\pi}(1+o(1))$ since $\delta_k=O(1)$; so $\limsup|\delta_k|\le2\pi\mathfrak a$. Conversely, let $t>\gamma_1$ be a non-ordinate, let $\gamma<t<\gamma'$ be the nearest ordinates and $k=N(t)$, so that $\gamma_k=\gamma$ and $\gamma_{k+1}=\gamma'$. On $(\gamma,\gamma')$ the function $G$ is continuous and decreasing ($G'=-\lL$), so $G(\gamma'^-)\le G(t)\le G(\gamma^+)$, while $s_k=k-\frac12-\NO(\gamma)=G(\gamma^+)-\frac12$ and $s_{k+1}=k+\frac12-\NO(\gamma')=G(\gamma'^-)+\frac12$. Hence $s_k\ge G(t)-\frac12$ and $s_{k+1}\le G(t)+\frac12$, so $\max(|s_k|,|s_{k+1}|)\ge|G(t)|-\frac12$, and $|G(t)-S(t)|=|\varpi(t)|/\pi\to0$. Both ordinates are within bounded distance of $t$ (Theorem~\CGref{thm:known_basic}). Choosing non-ordinates $t\to\infty$ with $|S(t)|\ge(\mathfrak a-\varepsilon)\log t$ gives indices $k\to\infty$ with $|\delta_k|=|s_k|/\lL(\xi_k)\ge2\pi(\mathfrak a-\varepsilon)+o(1)$.

(b) $|\tau_k^{X_k}-\tau_k|=|S_{X_k}(\tau_k)|/\lL(\tau_k)\le\sigma_{X_k}/\lL(\tau_k)\to0$, and $\sigma_X\ll\sqrt X/\log X$.
\end{proof}

So the corrected clocks with slowly growing cutoffs have exactly the pointwise envelope of the bare clock, and for every cutoff the envelope of $\tau_k$ is that of $S$. A smaller coefficient for the clock is a smaller coefficient for $S(t)$.

\begin{theorem}[Theorem E]\label{thm:envelope}
Let $\mathfrak e$ be continuous on $[\frac12,1]$ with $\mathfrak e(1)=0$, such that for every $\varepsilon>0$, $\zeta(\sigma+it)\ll_\varepsilon|t|^{\mathfrak e(\sigma)+\varepsilon}$ uniformly for $\frac12\le\sigma\le1$ and $|t|\ge2$, and put $\mathfrak e(\sigma)=\frac12-\sigma+\mathfrak e(1-\sigma)$ for $0\le\sigma<\frac12$. Then
\begin{equation}
\mathfrak a\le\frac1{2\pi\log2}\int_0^1\frac{\mathfrak e(\sigma)}{\sqrt{\sigma(2-\sigma)}}\,d\sigma .
\label{eq:envelope_formula}
\end{equation}
For $J\ge0$ let $(\mathsf k_0,\mathsf l_0)=(\frac{13}{84},\frac{55}{84})$, $(\mathsf k_j,\mathsf l_j)=\mathscr A^j(\mathsf k_0,\mathsf l_0)$ with $\mathscr A(\mathsf k,\mathsf l)=\bigl(\frac{\mathsf k}{2\mathsf k+2},\frac12+\frac{\mathsf l}{2\mathsf k+2}\bigr)$, $\sigma_j=\mathsf l_j-\mathsf k_j$, and let $\mathfrak e_J$ be the piecewise linear function through the points $(\sigma_j,\mathsf k_j)$, $0\le j\le J$, and $(1,0)$; the first vertices are $(\frac12,\frac{13}{84})$, $(\frac{139}{194},\frac{13}{194})$, $(\frac{173}{207},\frac{13}{414})$, $(\frac{773}{854},\frac{13}{854})$. Let $C_J$ be the right side of \CGeqref{eq:envelope_formula} for $\mathfrak e=\mathfrak e_J$. Then $\mathfrak a\le C_J$ for every $J$; the sequence $C_J$ decreases to a limit $C_\infty$ with $0\le C_J-C_\infty\le3\mathsf k_J\sqrt{1-\sigma_J}/(2\pi\log2)$, and
\begin{equation}
C_\infty=0.0968151142790\ldots,\qquad\limsup_{k\to\infty}|\tau_k-\gamma_k|=2\pi\mathfrak a\le2\pi C_\infty=0.608307303\ldots ;
\label{eq:envelope_value}
\end{equation}
for comparison $C_0=0.1010228\ldots$ (the exponent $\frac{13}{84}$ alone) and $C_7=0.0968160\ldots$.
\end{theorem}

\begin{proof}
\emph{Uniform bounds.} Extend $\mathfrak e$ by $0$ on $[1,\infty)$ and by $\frac12-\sigma$ on $(-\infty,0)$; the extension $\mathfrak e^*$ is continuous ($\mathfrak e(0)=\frac12+\mathfrak e(1)=\frac12$). For $\sigma\ge1$, $\zeta(\sigma+it)\ll\log|t|$ uniformly \cite[Thm.~3.5]{Titchmarsh1986}. For $-\frac1{50}\le\sigma<\frac12$, $\zeta(\sigma+it)=\chi(\sigma+it)\zeta(1-\sigma-it)$ with $|\chi(\sigma+it)|=(|t|/2\pi)^{1/2-\sigma}(1+O(1/|t|))$ uniformly \cite[(4.12.3)]{Titchmarsh1986}, and $|\zeta(1-\sigma-it)|=|\zeta(1-\sigma+it)|$. Hence $\zeta(\sigma+it)\ll_\varepsilon|t|^{\mathfrak e^*(\sigma)+\varepsilon}$ uniformly for $-\frac1{50}\le\sigma\le4$, $|t|\ge2$.

\emph{Backlund's method.} Let $0<\varepsilon_\circ<\frac1{100}$, $\sigma_\circ=1+\varepsilon_\circ$, $\varrho=1+2\varepsilon_\circ$. Then $1<\sigma_\circ<\varrho$, $\sigma^\sharp=\frac54+\frac32\varepsilon_\circ$, $\delta^\sharp=\frac14+\frac{\varepsilon_\circ}2\in(0,\frac12)$, and $\varrho/(\sigma_\circ-\frac12)=2$. For $s=\sigma_\circ+\varrho e^{i\theta}$, $\Re s\in[-\varepsilon_\circ,2+3\varepsilon_\circ]$ and $|\Im s|\le2$, so for $T$ large
\[
\begin{aligned}
\frac1\varkappa\log|\mathfrak f_\varkappa(s)|
&\le\log\max_\pm|(s\pm iT-1)\zeta(s\pm iT)|\\
&\le\log(T+4)+(\mathfrak e^*(\Re s)+\varepsilon)\log(T+2)+C_\varepsilon
=:F_\circ(\theta).
\end{aligned}
\]
an even continuous function. By Theorem~\CGref{thm:known_backlund}(a), along $\varkappa_m\to\infty$,
\[
\limsup_m\mathcal Z_{\varkappa_m}(\sigma_\circ,\varrho)\le\log T\cdot\frac1\pi\int_0^\pi\bigl(\mathfrak e^*(\sigma_\circ+\varrho\cos\theta)+\varepsilon\bigr)d\theta+O_{\varepsilon,\varepsilon_\circ}(1),
\]
the terms $\log(T+4)$ and $-\log\sqrt{\varepsilon_\circ^2+T^2}$ cancelling up to $O(1/T)$. The argument defining $S(T)$ runs along $2\to2+iT\to\sigma^\sharp+iT\to\frac12+iT$, and on the first two segments $\Re s>1$, where $\log\zeta$ is holomorphic; so $\pi S(T)=\arg\zeta(\sigma^\sharp+iT)+\Delta_{\sigma=\sigma^\sharp}^{1/2}\arg((\sigma+iT-1)\zeta(\sigma+iT))-\Delta_{\sigma=\sigma^\sharp}^{1/2}\arg(\sigma+iT-1)$, with $|\arg\zeta(\sigma^\sharp+iT)|\le\log\zeta(\sigma^\sharp)$ and the last term bounded by $\pi$. Theorem~\CGref{thm:known_backlund}(b) with $\varkappa=\varkappa_m\to\infty$ gives, for $T$ large,
\[
|S(T)|\le\frac{\log T}{2\log2}\cdot\frac1\pi\int_0^\pi\bigl(\mathfrak e^*(1+\varepsilon_\circ+(1+2\varepsilon_\circ)\cos\theta)+\varepsilon\bigr)d\theta+O_{\varepsilon,\varepsilon_\circ}(1).
\]
\emph{Limit.} Divide by $\log T$, let $T\to\infty$, then $\varepsilon\to0$ and $\varepsilon_\circ\to0$; $\mathfrak e^*$ is bounded and continuous, so by dominated convergence $\mathfrak a\le\frac1{2\pi\log2}\int_0^\pi\mathfrak e^*(1+\cos\theta)\,d\theta$. The integrand vanishes for $0\le\theta\le\frac\pi2$; on $[\frac\pi2,\pi]$ substitute $\sigma=1+\cos\theta$, $d\sigma=-\sin\theta\,d\theta$, $\sin\theta=\sqrt{\sigma(2-\sigma)}$, which gives \CGeqref{eq:envelope_formula}.

\emph{The profile.} Let $\varepsilon'>0$, $(\mathsf k_0,\mathsf l_0)=(\frac{13}{84}+\varepsilon',\frac{55}{84}+\varepsilon')$ and $(\mathsf k_j,\mathsf l_j)=\mathscr A^j(\mathsf k_0,\mathsf l_0)$ for $1\le j\le J$, exponent pairs by Theorem~\CGref{thm:known_exponents}(a),(c), with $\mathsf l_j-\mathsf k_j\ge\frac12$; by (b), $|\zeta(\mathsf l_j-\mathsf k_j+it)|\ll|t|^{\mathsf k_j}\log|t|$, and $|\zeta(1+it)|\ll\log|t|$. The points $(\mathsf l_j-\mathsf k_j,\mathsf k_j)$ and $(1,0)$ have decreasing ordinates, and Theorem~\CGref{thm:known_PL} between consecutive points gives the hypothesis of the theorem for the piecewise linear function through them. For $\varepsilon'=0$ these are the vertices of $\mathfrak e_J$; the right side of \CGeqref{eq:envelope_formula} depends continuously on $\varepsilon'$, so $\mathfrak a\le C_J$.

\emph{The limit.} From $\mathsf k_{j+1}=\mathsf k_j/(2\mathsf k_j+2)$ and $\mathsf l_{j+1}-\mathsf k_{j+1}=\frac12+(\mathsf l_j-\mathsf k_j)/(2\mathsf k_j+2)$ one gets $\mathsf k_j\le2^{-j}\mathsf k_0$, $\sigma_j<\sigma_{j+1}<1$ and $1-\sigma_{j+1}=(1-\sigma_j+\mathsf k_j)/(2\mathsf k_j+2)\le\frac12(1-\sigma_j)+\mathsf k_j$, so $\sigma_j\to1$. Each new vertex $(\sigma_{J+1},\mathsf k_{J+1})$ lies below the chord from $(\sigma_J,\mathsf k_J)$ to $(1,0)$: this is $\mathsf k_{J+1}(1-\sigma_J)\le\mathsf k_J(1-\sigma_{J+1})$, i.e.\ $\frac{1-\sigma_J}{2\mathsf k_J+2}\le\frac{1-\sigma_J+\mathsf k_J}{2\mathsf k_J+2}$, which holds. Hence $\mathfrak e_{J+1}\le\mathfrak e_J$ on $[\frac12,1]$ (they agree on $[\frac12,\sigma_J]$), the reflected parts are ordered in the same way, and $C_{J+1}\le C_J$; the limit $C_\infty$ exists. Since $0\le\mathfrak e_J-\mathfrak e_{J'}\le\mathsf k_J$ for $J'\ge J$, supported on $[\sigma_J,1]$ and, after reflection, on $[0,1-\sigma_J]$, and $\sigma(2-\sigma)\ge\sigma$ on $[0,1]$,
\[
\begin{split}0\le C_J-C_{J'}&\le\frac{\mathsf k_J}{2\pi\log2}\Bigl(\int_{\sigma_J}^1\frac{d\sigma}{\sqrt{\sigma(2-\sigma)}}+\int_0^{1-\sigma_J}\frac{d\sigma}{\sqrt\sigma}\Bigr)\\&\le\frac{\mathsf k_J\bigl((1-\sigma_J)\sqrt2+2\sqrt{1-\sigma_J}\bigr)}{2\pi\log2}\le\frac{3\mathsf k_J\sqrt{1-\sigma_J}}{2\pi\log2}.\end{split}
\]
\emph{The values.} On each linear piece $\mathfrak e(\sigma)=\alpha\sigma+\beta$ of $[0,1]$ a primitive of $\mathfrak e(\sigma)/\sqrt{\sigma(2-\sigma)}$ is $-\alpha\sqrt{\sigma(2-\sigma)}+(\alpha+\beta)\arcsin(\sigma-1)$; the vertices are rational, and the program \texttt{opc\_envelope\_iterations\_p1.py} (Data availability) evaluates $C_J$ from this primitive at $50$ digits for $0\le J\le30$. It gives $C_{30}=0.09681511427905377\ldots$ and $3\mathsf k_{30}\sqrt{1-\sigma_{30}}/(2\pi\log2)<5\cdot10^{-15}$, hence $C_\infty\in[C_{30}-5\cdot10^{-15},C_{30}]$ and \CGeqref{eq:envelope_value}; $C_7$ agrees with the $40$-digit quadrature of \texttt{opc\_envelope\_coefficient\_v14.py}. Multiply by $2\pi$ and use Proposition~\CGref{prop:envelope_equivalence}(a).
\end{proof}

\begin{remark}\label{rem:envelope_inputs}
(a) Trudgian \cite[\S5 of the preprint]{Trudgian2010} obtained $\mathfrak a\le0.1013$ from Huxley\textquoteright s exponent $\frac{32}{205}$ at $\sigma=\frac12$ by the same method; Theorem~\CGref{thm:envelope} replaces the exponents and lets the parameters of the disc tend to their limits. With the explicit exponents used by Bellotti and Wong ($\frac{27}{164}$ at $\sigma=\frac12$, Yang\textquoteright s lines $\sigma_j=1-j/(2^j-2)$ with exponents $1/(2^j-2)$, $4\le j\le9$ \cite{BellottiWong2025}) and linear interpolation between $\sigma=0$ and $\sigma=\frac12$, as in \cite[\S3.1]{BellottiWong2025}, the right side of \CGeqref{eq:envelope_formula} is $0.10074$, the limit of their $C_1=0.10076$ as their parameters tend to $(1,1)$; the same right-hand profile reflected by the functional equation gives $0.09783$. (b) Profiles of the growth exponent of $\zeta(\sigma+it)$ below $\mathfrak e_J$ are known: Trudgian and Yang \cite{TrudgianYang2023} close a convex hull of known exponent pairs, Bourgain\textquoteright s among them, under van der Corput\textquoteright s $\mathscr A$- and $\mathscr B$-processes and the $\mathscr C$-process and deduce piecewise bounds for that exponent on $[\frac12,1]$. Any such profile can be inserted in \CGeqref{eq:envelope_formula}; we do not evaluate it here. (c) For the Lindel\"of profile ($\mathfrak e=0$ on $[\frac12,1]$) the right side is $0.07863$: Backlund\textquoteright s method with a disc cannot give less, whatever the exponents. (d) Theorem~\CGref{thm:envelope} is asymptotic: the explicit constants of the bounds for $\zeta$ do not enter, which is why non-explicit exponents may be used, and it does not replace the explicit bounds of \cite{BellottiWong2025}, valid for all $T\ge e$.
\end{remark}

%======================================================================
\section{Conclusions}
\label{sec:conclusions}

\subsection{What was known}
Pointwise, $\delta_k=\tau_k-\gamma_k$ is bounded, with $\limsup|\delta_k|\le0.6331$ from the explicit bound of Bellotti and Wong \cite{BellottiWong2025}, and $\limsup|\delta_k|\le0.6365$ from Trudgian's asymptotic bound \cite{Trudgian2010}. On average, $s_k$ is asymptotically Gaussian with variance $\frac1{2\pi^2}\log\log T$, by Selberg's formulae for the Gram displacement proved by Korolev \cite{KorolevIzv2012,KorolevPIM2012}, so $\langle\delta_k^2\rangle\sim2\log\log T/\log^2T$: the clock is close to the zeros in height but misplaces almost every zero by an unbounded number of positions. Selberg's approximation \cite{Selberg1946a,Fujii1974} replaces $S(t)$ by short prime sums in $t$-mean.

\subsection{What this paper adds}
\begin{sloppypar}
\begin{enumerate}[leftmargin=1.6em]
\item \emph{An exact dictionary between clock-type models and the zeros} (Theorems~\CGref{thm:identity_general} and \CGref{thm:identity_g}). For any $C^1$ model $F$ and any continuous $\mathfrak g$, sums of $\mathfrak g$-increments of $k-\frac12-F(\gamma_k)$ over the ordinates are exactly $dF$-integrals of $\mathfrak g(N-F)$ up to two boundary terms; for powers this gives the even moments with the constant $\frac1{12}$ per zero, for exponentials it transfers the signed law of $S(t)$ to $s_k$ (Corollary~\CGref{cor:signed_law}).
\item \emph{The prime information needed by the clock} (Theorems~\CGref{thm:main} and \CGref{thm:variable_cutoff}). The corrected clock equation $\NO(t)+S_X(t)=k-\frac12$ holds at the true ordinates with mean-square ordinal defect $\frac1{2\pi^2}\log\frac{\log T}{\log X}+O(1+\sqrt{\log(\log T/\log X)})$, also with a cutoff varying at every index; the defect is bounded in mean square if and only if the primes reach a fixed power of the height.
\item \emph{An exact finite-height account} (Theorem~\CGref{thm:accounting}). The gain of the one-step clock is a diagonal prime response minus a cross term, and its residual differs from the defect by a transport term and an explicit nonnegative second-order term.
\item \emph{A computable clock that approaches the zeros} (Theorems~\CGref{thm:computable_clock}, \CGref{thm:second_step} and \CGref{thm:uniform_iteration}). With $X=T^{1/\sqrt{\log\log T}}$ the one-step clock has mean-square error $(\log\log\log T+O(\sqrt{\log\log\log T}))/\log^2T$ in height, against $(2+o(1))\log\log T/\log^2T$ for the bare clock; the $m$-th Newton-type step has error $\frac{H_m}{4\pi^2}+O(1+\sqrt{H_m})$ in units of zeros with absolute constants, and after $\log^*T+O(1)$ steps the error is $O(1)$ mean spacings, using no ordinate; the cutoffs involve absolute constants that are not computed. The proof needs an exact moment recursion, a multiplicity bound for the corrected points of every level (Theorem~\CGref{thm:multiplicity_all}) and a cutoff law that absorbs the multiplicity.
\item \emph{Tails and order} (Theorem~\CGref{thm:tails}, Corollaries~\CGref{cor:clock_tails} and \CGref{cor:sorted_clock}). The defect at the ordinates exceeds $\lambda$ on at most $2^{2-\sqrt{\lambda/C}}$ of the zeros when the primes reach $T^{c}$ with $c$ in a range fixed by $\lambda$; the errors of the corrected clocks have all moments bounded in units of $\sqrt{H_j}$; and sorting the final clock makes it nondecreasing without increasing its mean-square error.
\item \emph{The pointwise envelope} (Proposition~\CGref{prop:envelope_equivalence}, Theorem~\CGref{thm:envelope}). $\limsup|\tau_k-\gamma_k|=2\pi\limsup|S(t)|/\log t$ exactly, no slowly growing prime correction changes it, and Backlund's method with the profiles generated from Bourgain's exponent pair gives $\limsup|\tau_k-\gamma_k|\le2\pi C_\infty=0.6083073\ldots$.
\end{enumerate}
\end{sloppypar}

\subsection{Open problems}
\begin{question}\label{q:constant}
For fixed $c\in(0,\frac12]$, do the bounded means of $u_k(T^c)^2$ converge, and to $(2\pi^2)^{-1}\log(1/c)$ plus which constant?
\end{question}

\begin{question}\label{q:fixed_c}
Does $\langle r_k(T^c)^2\rangle_{\IT}$ stay bounded for fixed $c$, that is (Remark~\CGref{rem:newton}), is $\langle\mathfrak t_k(T^c)^2\rangle_{\IT}$ bounded? Is the error of the corrected clock after $\log^*T+O(1)$ steps bounded in units of zeros for \emph{almost every} index, and not only in mean square? Can Theorems~\CGref{thm:main} and \CGref{thm:computable_clock} be extended to $c>\frac12$?
\end{question}

\begin{question}\label{q:explicit}
Can the absolute constants $C_\star$, $H^*$ and $B_\star$ of Theorem~\CGref{thm:uniform_iteration}, and those of Fujii's and Montgomery--Vaughan's bounds on which they rest, be made explicit, so that the cutoff law \CGeqref{eq:cutoff_law} becomes a formula that can be evaluated at a given height?
\end{question}

Theorem~\CGref{thm:identity_general} applies to any counting staircase and any $C^1$ model, in particular to the zeros of a primitive Dirichlet $L$-function; an analogue of Theorem~\CGref{thm:main} would additionally need Selberg's approximation for $L(s,\chi)$ \cite{Selberg1946b,Fujii1974} and uniformity in the conductor.

\section*{Declaration of competing interest}
The author declares no competing interests.

\section*{Data availability}
No zero dataset was used to choose a construction or parameter. Two programs, \texttt{opc\_\allowbreak envelope\_\allowbreak coefficient\_\allowbreak v14.py} ($40$-digit quadrature) and \texttt{opc\_\allowbreak envelope\_\allowbreak iterations\_\allowbreak p1.py} (closed form at $50$ digits, $0\le J\le30$), evaluate the integral \CGeqref{eq:envelope_formula} for piecewise linear exponent profiles with rational vertices; they and their outputs are supplied in the accompanying compilation archive, in \texttt{code/} and \texttt{outputs/}. A third program, \texttt{verify\_\allowbreak envelope\_\allowbreak intervals\_\allowbreak v5.py}, encloses $C_{30}$ and the analytic tail by interval arithmetic at $80$ digits; its output \texttt{envelope\_\allowbreak intervals\_\allowbreak v5.json} confirms the upper bound $2\pi C_\infty<0.6083074$. The numerical manifest identifies each program, its precision and parameters, and the output supporting each displayed decimal. Python 3 and mpmath 1.3.0 were used. The programs evaluate consequences of the analytic formulae; no observed zero ordinate chooses the construction or a parameter.

\section*{Declaration of generative AI and AI-assisted technologies in the writing process}
During the preparation of this work the author used AI assistants (Claude, Anthropic, and ChatGPT, OpenAI) to edit the text, check proofs and write programs. The author reviewed and edited the content and takes full responsibility for it.

%======================================================================

\begin{thebibliography}{99}

\bibitem{BellottiWong2025}
C. Bellotti, P.-J. Wong, Improved estimates for the argument and zero-counting function of the Riemann zeta-function, Math. Comp. (2025) (statements numbered as in arXiv:2412.15470v2), \href{https://doi.org/10.1090/mcom/4133}{https://doi.org/10.1090/mcom/4133}.

\bibitem{Billingsley1995}
P. Billingsley, Probability and Measure, third ed., Wiley, New York, 1995.

\bibitem{Bourgain2017}
J. Bourgain, Decoupling, exponential sums and the Riemann zeta function, J. Amer. Math. Soc. 30 (1) (2017) 205--224 (theorems numbered as in arXiv:1408.5794v2), \href{https://doi.org/10.1090/jams/860}{https://doi.org/10.1090/jams/860}.

\bibitem{Carleson1966}
L. Carleson, On convergence and growth of partial sums of Fourier series, Acta Math. 116 (1966) 135--157, \href{https://doi.org/10.1007/BF02392815}{https://doi.org/10.1007/BF02392815}.

\bibitem{Corless1996}
R.M. Corless, G.H. Gonnet, D.E.G. Hare, D.J. Jeffrey, D.E. Knuth, On the Lambert $W$ function, Adv. Comput. Math. 5 (1996) 329--359, \href{https://doi.org/10.1007/BF02124750}{https://doi.org/10.1007/BF02124750}.

\bibitem{FazzariGerspachMinelli2025}
A. Fazzari, M. Gerspach, P. Minelli, Selberg's central limit theorem weighted by linear statistics of zeta zeros, preprint, arXiv:2507.04150, 2025, \href{https://doi.org/10.48550/arXiv.2507.04150}{https://doi.org/10.48550/arXiv.2507.04150}.

\bibitem{Folland1999}
G.B. Folland, Real Analysis: Modern Techniques and Their Applications, second ed., Wiley, New York, 1999.




\bibitem{Fujii1974}
A. Fujii, On the zeros of Dirichlet $L$-functions. I, Trans. Amer. Math. Soc. 196 (1974) 225--235, \href{https://doi.org/10.1090/S0002-9947-1974-0349603-2}{https://doi.org/10.1090/S0002-9947-1974-0349603-2}.

\bibitem{Fujii1975}
A. Fujii, On the distribution of the zeros of the Riemann zeta function in short intervals, Bull. Amer. Math. Soc. 81 (1) (1975) 139--142, \href{https://doi.org/10.1090/S0002-9904-1975-13674-3}{https://doi.org/10.1090/S0002-9904-1975-13674-3}.

\bibitem{Fujii1981}
A. Fujii, On the zeros of Dirichlet $L$-functions. II (with corrections to ``On the zeros of Dirichlet $L$-functions. I\textquotedblright\ and the subsequent papers), Trans. Amer. Math. Soc. 267 (1) (1981) 33--40.

\bibitem{HasanalizadeShenWong2022}
E. Hasanalizade, Q. Shen, P.-J. Wong, Counting zeros of the Riemann zeta function, J. Number Theory 235 (2022) 219--241, \href{https://doi.org/10.1016/j.jnt.2021.06.032}{https://doi.org/10.1016/j.jnt.2021.06.032}.

\bibitem{Hunt1968}
R.A. Hunt, On the convergence of Fourier series, in: Orthogonal Expansions and their Continuous Analogues (Proc. Conf., Edwardsville, Ill., 1967), Southern Illinois Univ. Press, Carbondale, 1968, pp. 235--255.

\bibitem{Ivic1985}
A. Ivi\'c, The Riemann Zeta-Function, Wiley, New York, 1985 (reprinted by Dover, Mineola, NY, 2003).

\bibitem{Kinnunen1997}
J. Kinnunen, The Hardy--Littlewood maximal function of a Sobolev function, Israel J. Math. 100 (1997) 117--124, \href{https://doi.org/10.1007/BF02773636}{https://doi.org/10.1007/BF02773636}.

\bibitem{KorolevPIM2012}
M. Korolev, Gram's law and the argument of the Riemann zeta function, Publ. Inst. Math. (Beograd) (N.S.) 92(106) (2012) 53--78, \href{https://doi.org/10.2298/PIM1206053K}{https://doi.org/10.2298/PIM1206053K}.

\bibitem{Korolev2016}
M.A. Korolev, Gram\textquoteright s law in the theory of the Riemann zeta-function. Part 1, Sovrem. Probl. Mat. 20 (2015) 3--161, \href{https://doi.org/10.4213/spm53}{https://doi.org/10.4213/spm53}; English transl.\ Proc. Steklov Inst. Math. 292, suppl.~2 (2016) S1--S146, \href{https://doi.org/10.1134/S0081543816030019}{https://doi.org/10.1134/S0081543816030019}.

\bibitem{IvicKorolev2018}
A. Ivi\'c, M.A. Korolev, On the distribution of values of the argument of the Riemann zeta-function, preprint, arXiv:1808.10768 (2018), \href{https://doi.org/10.48550/arXiv.1808.10768}{https://doi.org/10.48550/arXiv.1808.10768}.

\bibitem{KorolevIzv2012}
M.A. Korolev, On Gram's law in the theory of the Riemann zeta function, Izv. Math. 76 (2) (2012) 275--309, \href{https://doi.org/10.1070/IM2012v076n02ABEH002583}{https://doi.org/10.1070/IM2012v076n02ABEH002583}.


\bibitem{DiPlinioLerner2014}
F. Di Plinio, A.K. Lerner, On weighted norm inequalities for the Carleson and Walsh--Carleson operators, J. Lond. Math. Soc. (2) 90 (3) (2014) 654--674 (pages cited from arXiv:1312.0833v1), \href{https://doi.org/10.1112/jlms/jdu049}{https://doi.org/10.1112/jlms/jdu049}.

\bibitem{Littlewood1924}
J.E. Littlewood, On the zeros of the Riemann zeta-function, Proc. Cambridge Philos. Soc. 22 (1924) 295--318, \href{https://doi.org/10.1017/S0305004100014225}{https://doi.org/10.1017/S0305004100014225}.

\bibitem{Mertens1874}
F. Mertens, Ein Beitrag zur analytischen Zahlentheorie, J. Reine Angew. Math. 78 (1874) 46--62, \href{https://doi.org/10.1515/crll.1874.78.46}{https://doi.org/10.1515/crll.1874.78.46}.

\bibitem{MontgomeryVaughan1974}
H.L. Montgomery, R.C. Vaughan, Hilbert's inequality, J. Lond. Math. Soc. (2) 8 (1974) 73--82, \href{https://doi.org/10.1112/jlms/s2-8.1.73}{https://doi.org/10.1112/jlms/s2-8.1.73}.

\bibitem{Selberg1946a}
A. Selberg, Contributions to the theory of the Riemann zeta-function, Arch. Math. Naturvid. 48 (5) (1946) 89--155.

\bibitem{Selberg1946b}
A. Selberg, Contributions to the theory of Dirichlet's $L$-functions, Skr. Norske Vid.-Akad. Oslo I 1946 (3) (1946) 2--62.

\bibitem{Stein1970}
E.M. Stein, Singular Integrals and Differentiability Properties of Functions, Princeton University Press, Princeton, NJ, 1970.

\bibitem{Tenenbaum2015}
G. Tenenbaum, Introduction to Analytic and Probabilistic Number Theory, third ed., Grad. Stud. Math., vol. 163, Amer. Math. Soc., Providence, RI, 2015, \href{https://doi.org/10.1090/gsm/163}{https://doi.org/10.1090/gsm/163}.

\bibitem{Titchmarsh1986}
E.C. Titchmarsh, The Theory of the Riemann Zeta-Function, second ed., revised by D.R. Heath-Brown, Clarendon Press, Oxford, 1986.

\bibitem{Trudgian2010}
T.S. Trudgian, An improved upper bound for the argument of the Riemann zeta-function on the critical line, Math. Comp. 81 (278) (2012) 1053--1061, \href{https://doi.org/10.1090/S0025-5718-2011-02537-8}{https://doi.org/10.1090/S0025-5718-2011-02537-8}; preprint: A modest improvement on the function $S(T)$, arXiv:1010.4596 (2010).

\bibitem{Trudgian2014}
T.S. Trudgian, An improved upper bound for the argument of the Riemann zeta-function on the critical line II, J. Number Theory 134 (2014) 280--292, \href{https://doi.org/10.1016/j.jnt.2013.07.017}{https://doi.org/10.1016/j.jnt.2013.07.017}.

\bibitem{TrudgianYang2023}
T.S. Trudgian, A. Yang, Toward optimal exponent pairs, preprint, arXiv:2306.05599 (2023), \href{https://doi.org/10.48550/arXiv.2306.05599}{https://doi.org/10.48550/arXiv.2306.05599}.

\bibitem{Tsang1984}
K.-M. Tsang, The distribution of the values of the Riemann zeta-function, Ph.D. thesis, Princeton University, 1984.

\end{thebibliography}

\end{document}
